% Lesson 19 — Reading: Texts on helping with fund-raising activities/support groups in the acquisition of community resources/services — Anglais Tle A — Cours signé OnBuch+
% Programme officiel MINESEC. Compiler avec : tectonic EN19-reading-texts-on-helping-with-fund-raising-activities-s.tex
\newcommand{\DOCMATIERE}{Anglais}
\newcommand{\DOCNIVEAU}{Tle A}
\newcommand{\DOCMODULE}{Chapter 2 : Economic Life and Occupations}
\newcommand{\DOCLECON}{EN19}
\newcommand{\DOCTITRE}{Lesson 19 — Reading: Texts on helping with fund-raising activities/support groups in the acquisition of community resources/services}
% OnBuch+ — préambule « cours signés » ENRICHI (compilé avec tectonic / XeLaTeX)
% Système visuel « Pop » d'OnBuch+ : fond crème (jamais blanc), bordures encre
% épaisses, ombres « dures » décalées, titres Archivo Black en capitales.
% Version MATHÉMATIQUES : graphes (pgfplots/tikz), boîtes pédagogiques
% (définition, propriété, méthode, attention, le savais-tu, exercice résolu,
% exercice, TP, bilan), page de garde et signature OnBuch+ ancrée en bas.
%
% Macros attendues dans main.tex AVANT \input{preamble} :
%   \DOCMATIERE  \DOCNIVEAU  \DOCMODULE  \DOCLECON (numéro)  \DOCTITRE
\documentclass[11pt]{article}

\usepackage{fontspec}
\usepackage[english,french]{babel} % français = langue principale ; l'anglais via \eng{..} / textebox
\usepackage[a4paper,margin=2cm,top=3.4cm,bottom=2.8cm,headheight=1.4cm,footskip=1.3cm]{geometry}
\usepackage{amsmath,amssymb,amsfonts}
\usepackage{mathtools} % \xmapsto, \coloneqq... souvent utilisés par les modèles
\usepackage{cancel} % simplifications barrées (bilans de forces)
% \abs{z} (module, valeur absolue) et \norm{u} (norme), utilisés par les modèles
\providecommand{\abs}{}\let\abs\relax\DeclarePairedDelimiter{\abs}{\lvert}{\rvert}
\providecommand{\norm}{}\let\norm\relax\DeclarePairedDelimiter{\norm}{\lVert}{\rVert}
\usepackage[table]{xcolor} % [table] : fournit aussi \rowcolors (alternance de lignes)
\usepackage{pagecolor}
\usepackage{tikz}
\usetikzlibrary{arrows.meta,positioning,calc,shapes.geometric,shapes.misc,shapes.arrows,decorations.pathmorphing,decorations.pathreplacing,decorations.markings,patterns,shadows,mindmap,trees,fit,backgrounds,matrix,intersections,angles,quotes}
\usepackage{pgfplots}
\usepackage[version=4]{mhchem} % équations nucléaires \ce{^{235}_{92}U}
\usepackage[european,siunitx]{circuitikz} % schémas électriques
\pgfplotsset{compat=1.18}
\usepgfplotslibrary{fillbetween,groupplots} % aires entre courbes (calcul intégral)
\usepackage{enumitem}
\usepackage{fancyhdr}
% [most] charge toutes les bibliothèques tcolorbox (listings, minted, raster,
% documentation...) et gonfle inutilement la mémoire TeX sur un document long
% (~300 boîtes) : on ne charge que ce qui est réellement utilisé.
\usepackage[skins,breakable]{tcolorbox}
\usepackage{graphicx}
\usepackage{colortbl}
\usepackage{booktabs}
\usepackage{tabularx}
\usepackage{diagbox} % case d'en-tête diagonale des tableaux à double entrée
% Colonnes extensibles centrée / gauche / droite (C, L, R), souvent utilisées par les modèles
\newcolumntype{C}{>{\centering\arraybackslash}X}
\newcolumntype{L}{>{\raggedright\arraybackslash}X}
\newcolumntype{R}{>{\raggedleft\arraybackslash}X}
\usepackage{array}
\usepackage{multicol}
\usepackage{multirow}
\usepackage{siunitx}
% parse-numbers=false : accepte aussi \SI{2,5 \times 10^{-3}}{...}
\sisetup{locale=FR,output-decimal-marker={,},group-minimum-digits=4,parse-numbers=false}
\DeclareSIUnit\ohm{\ensuremath{\symup{\Omega}}} % U+2126 absent de la police
\DeclareSIUnit\division{div} % réglages d'oscilloscope (ms/div, V/div)
\DeclareSIUnit\tr{tr} % tours (vitesse de rotation, ex. \SI{1500}{\tr\per\minute})
\DeclareSIUnit\molar{M} % concentration molaire (ex. \SI{3}{\molar}, \SI{1}{\micro\molar})
\DeclareSIUnit\an{an} % année (le modèle confond parfois avec \year, un compteur TeX) — ex. \SI{5}{\kilo\meter\per\an}
\DeclareSIUnit\FCFA{FCFA} % franc CFA (ex. \SI{500}{\FCFA\per\kilo\gram})
\DeclareSIUnit\degreeBrix{\textdegree Brix} % teneur en sucre d'un sirop/jus (réfractométrie)
\DeclareSIUnit\month{mois} % le modèle utilise parfois \month, un compteur TeX, comme unité de durée
\usepackage{qrcode}
% \degree : commande fréquemment utilisée par les modèles pour un angle ou une
% température en dehors de \SI{}{} (non fournie par les paquets chargés).
\providecommand{\degree}{^{\circ}}
% \qed (fin de démonstration), utilisé par les modèles sans amsthm
\providecommand{\qed}{\hfill$\square$}
% \barre : double barre des valeurs interdites dans un tableau de variations (tkz-tab)
\providecommand{\barre}{\ensuremath{\|}}
% environnement proof (amsthm non chargé) utilisé par les modèles
\ifdefined\proof\else\newenvironment{proof}[1][Démonstration]{\par\noindent\textit{#1.}\ }{\qed\par}\fi
\usepackage{lastpage}
\usepackage{hyperref}
\hypersetup{hidelinks}

% ============================================================
% Palette « Pop » OnBuch+ — mêmes hex que la classe P (plus_ui.dart)
% ============================================================
\definecolor{poporange}{HTML}{F59321}
\definecolor{popdark}{HTML}{E07A0C}
\definecolor{popcream}{HTML}{FFF6E8}
\definecolor{popink}{HTML}{1C1714}
\definecolor{poppurple}{HTML}{7A5AE0}
\definecolor{popgreen}{HTML}{2E9E6A}
\definecolor{poppink}{HTML}{E0457B}
\definecolor{popblue}{HTML}{2F7DE1}
\definecolor{popgold}{HTML}{C98A12}
\definecolor{popmuted}{HTML}{8A7F76}
\definecolor{popline}{HTML}{EADFD2}
\definecolor{popgray}{HTML}{8A7F76}
\colorlet{popgrey}{popgray}
% Alias tolérants (le modèle nomme parfois les couleurs différemment)
\colorlet{popwhite}{white}
\colorlet{popblack}{popink}
\colorlet{popred}{poppink}
\colorlet{popyellow}{popgold}
% Teintes claires (fonds de tableaux/encadrés) — le modèle nomme parfois
% les couleurs avec un suffixe L (ex. popblueL) sans les avoir définies.
\colorlet{poporangeL}{poporange!20}
\colorlet{popdarkL}{popdark!20}
\colorlet{poppurpleL}{poppurple!20}
\colorlet{popgreenL}{popgreen!20}
\colorlet{poppinkL}{poppink!20}
\colorlet{popblueL}{popblue!20}
\colorlet{popgoldL}{popgold!20}
\colorlet{popredL}{popred!20}
\colorlet{popgrayL}{popgray!20}
% Teintes claires pour fonds de tableaux / zones
\colorlet{popblueL}{popblue!10!white}
\colorlet{poporangeL}{poporange!14!white}
\colorlet{popgreenL}{popgreen!12!white}
\colorlet{poppurpleL}{poppurple!12!white}
\colorlet{poppinkL}{poppink!12!white}
\colorlet{popgoldL}{popgold!14!white}

\pagecolor{popcream}

% ============================================================
% Typographie
% ============================================================
\defaultfontfeatures{Ligatures=TeX}
\setmainfont{PlusJakartaSans-Medium.ttf}[Path=../../../fonts/,BoldFont=PlusJakartaSans-Bold.ttf]
% Maths en sans-serif (Fira Math) avec chiffres et lettres droites de Plus
% Jakarta, pour que les formules se fondent dans le texte.
\usepackage[math-style=ISO,bold-style=ISO]{unicode-math}
\setmathfont{FiraMath-Regular.otf}[Path=../../../fonts/]
\setmathfont{PlusJakartaSans-Medium.ttf}[Path=../../../fonts/,range={up/num,up/{Latin,latin},\mathup}]
\usepackage{icomma} % virgule décimale sans espace en mode math
% Caractères mathématiques Unicode absents des polices de texte (Plus Jakarta,
% Archivo Black) : rendus par leur équivalent mathématique, en texte comme en
% formule — sinon ils s'affichent en carré vide (ex. « Arithmétique dans ℤ »).
\usepackage{newunicodechar}
\newunicodechar{ℝ}{\ensuremath{\mathbb{R}}}
\newunicodechar{ℤ}{\ensuremath{\mathbb{Z}}}
\newunicodechar{ℂ}{\ensuremath{\mathbb{C}}}
\newunicodechar{ℕ}{\ensuremath{\mathbb{N}}}
\newunicodechar{ℚ}{\ensuremath{\mathbb{Q}}}
\newunicodechar{𝔻}{\ensuremath{\mathbb{D}}}
\newunicodechar{α}{\ensuremath{\alpha}}
\newunicodechar{β}{\ensuremath{\beta}}
\newunicodechar{γ}{\ensuremath{\gamma}}
\newunicodechar{δ}{\ensuremath{\delta}}
\newunicodechar{ε}{\ensuremath{\varepsilon}}
\newunicodechar{θ}{\ensuremath{\theta}}
\newunicodechar{λ}{\ensuremath{\lambda}}
\newunicodechar{μ}{\ensuremath{\mu}}
\newunicodechar{σ}{\ensuremath{\sigma}}
\newunicodechar{τ}{\ensuremath{\tau}}
\newunicodechar{φ}{\ensuremath{\varphi}}
\newunicodechar{ψ}{\ensuremath{\psi}}
\newunicodechar{ω}{\ensuremath{\omega}}
\newunicodechar{Σ}{\ensuremath{\Sigma}}
% majuscules produites par \MakeUppercase dans les titres (α-aminés -> Α)
\newunicodechar{Α}{\ensuremath{\alpha}}
\newunicodechar{Β}{\ensuremath{\beta}}
\newunicodechar{Γ}{\ensuremath{\Gamma}}
\newunicodechar{Δ}{\ensuremath{\Delta}}
\newunicodechar{Ε}{\ensuremath{\varepsilon}}
\newunicodechar{Θ}{\ensuremath{\Theta}}
\newunicodechar{Λ}{\ensuremath{\Lambda}}
\newunicodechar{Μ}{\ensuremath{\mu}}
\newunicodechar{Τ}{\ensuremath{\tau}}
\newunicodechar{Φ}{\ensuremath{\Phi}}
\newunicodechar{Ψ}{\ensuremath{\Psi}}
\newunicodechar{Ω}{\ensuremath{\Omega}}
\newunicodechar{↦}{\ensuremath{\mapsto}}
\newunicodechar{∈}{\ensuremath{\in}}
\newunicodechar{∉}{\ensuremath{\notin}}
\newunicodechar{∧}{\ensuremath{\wedge}}
\newunicodechar{⇔}{\ensuremath{\Leftrightarrow}}
\newunicodechar{⟺}{\ensuremath{\Longleftrightarrow}}
\newunicodechar{⇒}{\ensuremath{\Rightarrow}}
\newunicodechar{⟹}{\ensuremath{\Longrightarrow}}
\newunicodechar{∘}{\ensuremath{\circ}}
\newunicodechar{∪}{\ensuremath{\cup}}
\newunicodechar{∩}{\ensuremath{\cap}}
\newunicodechar{≡}{\ensuremath{\equiv}}
\newunicodechar{⊂}{\ensuremath{\subset}}
\newunicodechar{⊥}{\ensuremath{\perp}}
\newunicodechar{∃}{\ensuremath{\exists}}
\newunicodechar{∀}{\ensuremath{\forall}}
\newunicodechar{∛}{\ensuremath{\sqrt[3]{\,}}}
\newunicodechar{∜}{\ensuremath{\sqrt[4]{\,}}}
\newunicodechar{⃗}{\ensuremath{{}^{\to}}}
\newunicodechar{ⁿ}{\ifmmode^{n}\else\textsuperscript{\ensuremath{n}}\fi}
\newunicodechar{ˣ}{\ifmmode^{x}\else\textsuperscript{\ensuremath{x}}\fi}
\newunicodechar{ᵇ}{\ifmmode^{b}\else\textsuperscript{\ensuremath{b}}\fi}
\newunicodechar{ʳ}{\ifmmode^{r}\else\textsuperscript{\ensuremath{r}}\fi}
\newunicodechar{ᵘ}{\ifmmode^{u}\else\textsuperscript{\ensuremath{u}}\fi}
\newunicodechar{ʸ}{\ifmmode^{y}\else\textsuperscript{\ensuremath{y}}\fi}
\newunicodechar{ᵃ}{\ifmmode^{a}\else\textsuperscript{\ensuremath{a}}\fi}
\newunicodechar{ᵉ}{\ifmmode^{e}\else\textsuperscript{\ensuremath{e}}\fi}
\newunicodechar{ᵏ}{\ifmmode^{k}\else\textsuperscript{\ensuremath{k}}\fi}
\newunicodechar{ᵖ}{\ifmmode^{p}\else\textsuperscript{\ensuremath{p}}\fi}
\newunicodechar{ᴺ}{\ifmmode^{N}\else\textsuperscript{\ensuremath{N}}\fi}
\newunicodechar{ᶜ}{\ifmmode^{c}\else\textsuperscript{\ensuremath{c}}\fi}
\newunicodechar{ʰ}{\ifmmode^{h}\else\textsuperscript{\ensuremath{h}}\fi}
\newunicodechar{ᵠ}{\ifmmode^{q}\else\textsuperscript{\ensuremath{q}}\fi}
\newunicodechar{ₙ}{\ifmmode_{n}\else\textsubscript{\ensuremath{n}}\fi}
\newunicodechar{ₐ}{\ifmmode_{a}\else\textsubscript{\ensuremath{a}}\fi}
\newunicodechar{ᵢ}{\ifmmode_{i}\else\textsubscript{\ensuremath{i}}\fi}
\newunicodechar{ᵣ}{\ifmmode_{r}\else\textsubscript{\ensuremath{r}}\fi}
\newunicodechar{ₖ}{\ifmmode_{k}\else\textsubscript{\ensuremath{k}}\fi}
\newunicodechar{ᵦ}{\ifmmode_{\beta}\else\textsubscript{\ensuremath{\beta}}\fi}
\newunicodechar{ₓ}{\ifmmode_{x}\else\textsubscript{\ensuremath{x}}\fi}
\newfontfamily\popdisplay{ArchivoBlack-Regular.ttf}[Path=../../../fonts/]
\newfontfamily\poplogo{SpaceGrotesk-Bold.ttf}[Path=../../../fonts/]
\newfontfamily\popsemi{PlusJakartaSans-SemiBold.ttf}[Path=../../../fonts/]
\newfontfamily\popxbold{PlusJakartaSans-ExtraBold.ttf}[Path=../../../fonts/]
\newfontfamily\popxboldls{PlusJakartaSans-ExtraBold.ttf}[Path=../../../fonts/,LetterSpace=3.0]

\setlist[enumerate]{leftmargin=1.7em,itemsep=5pt,topsep=4pt}
\setlist[itemize]{leftmargin=1.7em,itemsep=4pt,topsep=4pt,label={\color{poporange}\textbullet}}
\setlength{\parindent}{0pt}
\setlength{\parskip}{5pt}
\renewcommand{\arraystretch}{1.25}

% pgfplots : style Pop par défaut
\pgfplotsset{
  every axis/.append style={axis line style={popink,line width=1pt},tick style={popink},
    grid style={popline,line width=0.5pt},label style={font=\small},tick label style={font=\footnotesize}},
  popcurve/.style={poporange,line width=1.8pt,no markers,smooth},
}
% Même style disponible dans un \draw TikZ simple (hors pgfplots)
\tikzset{popcurve/.style={poporange,line width=1.8pt}}
\tikzset{popgrid/.style={popline,very thin}}
\colorlet{popcurve}{poporange} % le nom du style est parfois utilisé comme couleur

% ============================================================
% Pill / squiggle
% ============================================================
\newcommand{\poppill}[3]{%
  \tikz[baseline=-0.55ex]\node[draw=popink,line width=1.2pt,rounded corners=2.6pt,%
    fill=#2,inner xsep=6.5pt,inner ysep=3pt]{%
    \popxboldls\fontsize{7}{7}\selectfont\color{#3}\MakeUppercase{#1}};%
}
\newcommand{\popsquiggle}[1]{%
  \tikz[baseline=-0.4ex]{
    \draw[#1,line width=2.6pt,line cap=round]
      plot [smooth] coordinates {(0,0.09) (0.34,0.01) (0.68,0.21) (1.02,0.01) (1.36,0.10)};
  }%
}
% Mot-clé surligné (marqueur orange sous le texte)
\newcommand{\cle}[1]{{\popxbold\color{popdark}#1}}

% ============================================================
% En-tête / pied
% ============================================================
\pagestyle{fancy}
\fancyhf{}
\renewcommand{\headrulewidth}{0pt}
\renewcommand{\footrulewidth}{0pt}
\lhead{\raisebox{-0.15ex}{\poplogo\fontsize{13}{13}\selectfont\color{popink}OnBuch\color{poporange}+}}
\rhead{\poppill{\DOCMATIERE\ \textbullet\ \DOCNIVEAU\ \textbullet\ Leçon \DOCLECON}{white}{popink}}
\lfoot{\popxbold\fontsize{7.6}{7.6}\selectfont\color{popmuted}\MakeUppercase{Cours signé OnBuch+}\ \textbullet\ {\popsemi\DOCTITRE}}
\rfoot{\tikz[baseline=-0.6ex]\node[rounded corners=8.5pt,fill=popink,minimum height=17pt,minimum width=17pt,inner xsep=5pt,inner sep=0]{\popxbold\fontsize{8}{8}\selectfont\color{popcream}\thepage\,/\,\pageref*{LastPage}};}

% ============================================================
% Titres
% ============================================================
\newcommand{\chaptitre}[2]{%
  \begin{tcolorbox}[enhanced,colback=poporange,colframe=popink,boxrule=2.6pt,arc=11pt,
      drop shadow=popink,left=15pt,right=15pt,top=12pt,bottom=11pt,before skip=3pt,after skip=18pt]
    \poppill{#2}{popink}{popcream}\par\vspace{9pt}
    {\raggedright\hyphenpenalty=10000\exhyphenpenalty=10000\popdisplay\fontsize{22}{24}\selectfont\color{popcream}\MakeUppercase{#1}\par}
  \end{tcolorbox}
}
% Section numérotée : pastille orange avec le numéro + titre Archivo
\newcounter{popsec}
\newcommand{\coursec}[1]{%
  \stepcounter{popsec}\setcounter{popsub}{0}%
  \par\vspace{16pt}\noindent
  \begin{minipage}{\linewidth}
  \tikz[baseline=-0.7ex]\node[draw=popink,line width=1.6pt,fill=poporange,rounded corners=4pt,minimum size=22pt,inner sep=0]{\popdisplay\fontsize{12}{12}\selectfont\color{popcream}\thepopsec};%
  \hspace{8pt}{\popdisplay\fontsize{15}{17}\selectfont\color{popink}\MakeUppercase{#1}}\par
  \vspace{4pt}\noindent{\color{poporange}\rule{36pt}{3.6pt}}
  \end{minipage}\par\nopagebreak\vspace{8pt}
}
\newcounter{popsub}
\newcommand{\courssub}[1]{%
  \stepcounter{popsub}%
  \par\vspace{9pt}\noindent{\popxbold\fontsize{11.5}{13}\selectfont\color{popink}{\color{poporange}\thepopsec.\thepopsub}\hspace{6pt}#1}\par\nopagebreak\vspace{4pt}
}

% ============================================================
% Boîtes pédagogiques — même recette : fond blanc, bordure épaisse colorée,
% ombre dure encre, badge plein. Titre optionnel [#1].
% ============================================================
\tcbset{popbase/.style={breakable,enhanced,colback=white,boxrule=2.3pt,arc=9pt,
  shadow={3pt}{-3pt}{0pt}{popink},left=14pt,right=14pt,top=11pt,bottom=11pt,before skip=11pt,after skip=15pt,
  fontupper=\popsemi\fontsize{10.6}{14.5}\selectfont\color{popink},
  fontlower=\popsemi\fontsize{10.6}{14.5}\selectfont\color{popink}}}
% Coche dessinée (le glyphe ✓ n'existe pas dans les polices du document)
\newcommand{\popcheck}[1]{\tikz[baseline=-0.5ex]\draw[#1,line width=1.6pt,line cap=round,line join=round](-0.13,0.0)--(-0.04,-0.09)--(0.13,0.1);}
\providecommand{\coche}{\popcheck{popgreen}}
\providecommand{\resultat}[1]{\par\smallskip\noindent\fcolorbox{popgreen}{popgreenL}{\parbox{\dimexpr\linewidth-2\fboxsep-2\fboxrule\relax}{\textbf{Résultat.} #1}}\par\smallskip}
\newcommand{\popboxhead}[3]{\poppill{#1}{#2}{white}\ifx\relax#3\relax\else\hspace{6pt}{\popxbold\fontsize{10.6}{12}\selectfont\color{popink}#3}\fi\par\vspace{7pt}}

\newtcolorbox{aretenir}[1][]{popbase,colframe=popgreen,before upper={\popboxhead{À retenir}{popgreen}{#1}}}
% Texte en anglais (document de lecture, dialogue, extrait) : cadre
% d'étude. Un titre optionnel (source, auteur) s'affiche dans l'en-tête.
\newtcolorbox{textebox}[1][]{popbase,colframe=popblue,before upper={\popboxhead{Text}{popblue}{#1}\begin{otherlanguage}{english}},after upper={\end{otherlanguage}}}
% Marques ✓ / ✗ (forme correcte / fautive) souvent utilisées par les modèles
\providecommand{\cross}{\textcolor{poppink}{\ensuremath{\times}}}
\providecommand{\xmark}{\cross}\providecommand{\crossmark}{\cross}\providecommand{\cmark}{\ensuremath{\checkmark}}
% Ligne de réponse / justification à compléter (inventée par les modèles)
\providecommand{\justification}[1]{\par\noindent\textit{Justification~:} \dotfill\par}
% \textipa (package tipa non chargé) : la police ne couvre pas l'API -> simple passage
\providecommand{\textipa}[1]{#1}
% Soulignement (ulem non chargé) : \uline, \ul
\providecommand{\uline}[1]{\underline{#1}}\providecommand{\ul}[1]{\underline{#1}}
% \glossaire{\item ...} inventée par les modèles : liste de glossaire
\providecommand{\glossaire}[1]{\begin{itemize}#1\end{itemize}}
% \alert (beamer) inventée par les modèles : texte mis en évidence
\providecommand{\alert}[1]{\textbf{\textcolor{poppink}{#1}}}
% variantes ASCII d'ordinaux écrites en macro
\providecommand{\iere}{\textsuperscript{re}}\providecommand{\ieme}{\textsuperscript{e}}\providecommand{\ere}{\textsuperscript{re}}\providecommand{\eme}{\textsuperscript{e}}\providecommand{\er}{\textsuperscript{er}}\providecommand{\ie}{\textsuperscript{e}}
% Ordinaux français écrits en macro par les modèles : 1\ière, 3\ième
\newcommand{\ière}{\textsuperscript{re}}\newcommand{\ième}{\textsuperscript{e}}
% \exemple (inventée par les modèles) : étiquette « Exemple : »
\providecommand{\exemple}{\textbf{Exemple~:}\ }
% \square utilisable aussi hors mode math (cases à cocher)
\AtBeginDocument{\let\squareMath\square\renewcommand{\square}{\ensuremath{\squareMath}}}
% Mot ou phrase en anglais à l'intérieur d'un paragraphe français
\newcommand{\eng}[1]{\ifmmode\text{\textit{#1}}\else\begingroup\selectlanguage{english}\itshape #1\endgroup\fi}
\let\esp\eng % alias toléré
\newtcolorbox{exemplebox}[1][]{popbase,colframe=poporange,before upper={\popboxhead{Exemple}{poporange}{#1}}}
\newtcolorbox{definition}[1][]{popbase,colframe=popblue,before upper={\popboxhead{Définition}{popblue}{#1}}}
\newtcolorbox{propriete}[1][]{popbase,colframe=poppurple,before upper={\popboxhead{Propriété}{poppurple}{#1}}}
\newtcolorbox{methode}[1][]{popbase,colframe=popgold,before upper={\popboxhead{Méthode}{popgold}{#1}}}
\newtcolorbox{attention}[1][]{popbase,colframe=poppink,before upper={\popboxhead{Attention}{poppink}{#1}}}
\newtcolorbox{experience}[1][]{popbase,colframe=popblue,colback=popblueL,before upper={\popboxhead{Expérience}{popblue}{#1}}}
% « Le savais-tu ? » : carte violette pleine (contexte, culture, Cameroun)
\newtcolorbox{savaistu}[1][]{popbase,colback=poppurple,colframe=popink,
  fontupper=\popsemi\fontsize{10.4}{14.2}\selectfont\color{white},
  before upper={\poppill{Le savais-tu\,?}{white}{poppurple}\ifx\relax#1\relax\else\hspace{6pt}{\popxbold\color{white}#1}\fi\par\vspace{7pt}}}
% Exercice résolu : énoncé en haut, \tcblower puis la solution sur fond crème
\newcounter{popexo}\newcounter{popexr}
\newtcolorbox{exoresolu}[1][]{popbase,colframe=poporange,colbacklower=popcream,
  segmentation style={popink,line width=1.2pt,dashed},
  before upper={\stepcounter{popexr}\popboxhead{Exercice résolu \thepopexr}{poporange}{#1}},
  before lower={\poppill{Solution}{popink}{popcream}\par\vspace{6pt}}}
% Exercice (non résolu en place) — la correction est dans la partie Corrigés
\newtcolorbox{exercice}[1][]{popbase,colframe=popink,
  before upper={\stepcounter{popexo}\popboxhead{Exercice \thepopexo}{popink}{#1}}}
% Corrigé
\newtcolorbox{corrige}[1][]{popbase,colframe=popgreen,colback=popgreenL,
  before upper={\popboxhead{Corrigé}{popgreen}{#1}}}
% Objectifs / prérequis / bilan
\newtcolorbox{objectifs}{popbase,colframe=popink,colback=poporangeL,
  before upper={\popboxhead{Objectifs de la leçon}{popink}{}}}
\newtcolorbox{prerequis}{popbase,colframe=popink,colback=popblueL,
  before upper={\popboxhead{Prérequis}{popblue}{}}}
\newtcolorbox{bilan}{popbase,colframe=popink,colback=white,boxrule=2.8pt,
  before upper={\popboxhead{Fiche bilan}{poporange}{L'essentiel à connaître pour l'examen}}}
% Figure encadrée avec légende
\newtcolorbox{popfigure}[1][]{enhanced,breakable=false,colback=white,colframe=popline,boxrule=1.4pt,arc=8pt,
  left=10pt,right=10pt,top=10pt,bottom=8pt,before skip=10pt,after skip=12pt,halign=center,
  lower separated=false,halign lower=center,
  fontlower=\popsemi\fontsize{9}{11.5}\selectfont\color{popmuted},
  #1}
\newcommand{\legende}[1]{\par\vspace{6pt}{\popsemi\fontsize{9}{11.5}\selectfont\color{popmuted}#1}}

% ============================================================
% Signature OnBuch+
% ============================================================
\newcommand{\poplogoinline}[1]{{\poplogo\fontsize{#1}{#1}\selectfont\color{popink}OnBuch\color{poporange}+}}
\newcommand{\signatureonbuch}{%
  \begin{tcolorbox}[enhanced,colback=popink,colframe=popink,boxrule=2.2pt,arc=10pt,
      drop shadow=poporange,left=14pt,right=14pt,top=10pt,bottom=10pt,before skip=0pt,after skip=0pt]
    \tikz[baseline=-0.7ex]\node[circle,fill=popgreen,draw=popcream,line width=1.3pt,minimum size=22pt,inner sep=0]{\popcheck{white}};%
    \hspace{9pt}\begin{minipage}[c]{0.62\linewidth}
      {\popxbold\fontsize{12}{13}\selectfont\color{popcream}Signé\ \textbullet\ {\poplogo OnBuch\color{poporange}+}}\par\vspace{2pt}
      {\raggedright\popsemi\fontsize{8}{10.5}\selectfont\color{popline}\MakeUppercase{Équipe pédagogique OnBuch+}\\\MakeUppercase{Conforme au programme officiel MINESEC}\par}
    \end{minipage}\hfill
    {\poplogo\fontsize{22}{22}\selectfont\color{popcream}OnBuch\color{poporange}+}
  \end{tcolorbox}%
}

% ============================================================
% Page de garde
% #1 titre · #2 matière · #3 niveau · #4 module · #5 n° leçon · #6 durée ·
% #7 description / objectifs (texte ou itemize)
% ============================================================
\newcommand{\pagegarde}[7]{%
  \thispagestyle{empty}
  \newgeometry{a4paper,margin=1.7cm,top=1.6cm,bottom=1.4cm}%
  \begin{tikzpicture}[remember picture,overlay]
    \fill[poporange!18] ([xshift=-3.2cm,yshift=-3.6cm]current page.north east) circle (4.6cm);
    \fill[poppurple!14] ([xshift=2.4cm,yshift=7.6cm]current page.south west) circle (3.8cm);
  \end{tikzpicture}%
  {\poplogo\fontsize{30}{30}\selectfont\color{popink}OnBuch\color{poporange}+}\hfill
  \raisebox{0.6ex}{\poppill{Cours signé \textbullet\ Premium}{popink}{popcream}}\par
  \vspace{1pt}\popsquiggle{poppurple}\par
  \vspace{8pt}
  \begin{tcolorbox}[enhanced,colback=poporange,colframe=popink,boxrule=2.8pt,arc=13pt,
      drop shadow=popink,left=18pt,right=18pt,top=12pt,bottom=13pt,after skip=10pt]
    \poppill{#2 \textbullet\ #3}{popink}{popcream}\hspace{5pt}\poppill{#4}{white}{popink}\par\vspace{8pt}
    {\popdisplay\fontsize{14}{14}\selectfont\color{popink}LEÇON #5}\par\vspace{4pt}
    {\raggedright\hyphenpenalty=10000\exhyphenpenalty=10000\popdisplay\fontsize{25}{27}\selectfont\color{popcream}\MakeUppercase{#1}\par}
    \vspace{8pt}
    {\popxbold\fontsize{9.5}{9.5}\selectfont\color{popink}\MakeUppercase{Durée conseillée : #6}}
  \end{tcolorbox}
  \begin{tcolorbox}[enhanced,colback=white,colframe=popink,boxrule=2.2pt,arc=10pt,
      drop shadow=popink,left=16pt,right=16pt,top=10pt,bottom=10pt,after skip=8pt]
    \poppill{Dans cette leçon}{poppurple}{white}\par\vspace{5pt}
    \popsemi\fontsize{9.3}{12.6}\selectfont\color{popink}#7
  \end{tcolorbox}
  \noindent\begin{minipage}[t]{0.487\linewidth}
    \begin{tcolorbox}[enhanced,colback=popgreen,colframe=popink,boxrule=2.2pt,arc=10pt,
        drop shadow=popink,left=11pt,right=11pt,top=10pt,bottom=10pt,height=3.1cm,valign=center]
      \tcbox[colback=white,colframe=popink,boxrule=1.3pt,arc=5pt,boxsep=0pt,nobeforeafter,
        left=3pt,right=3pt,top=3pt,bottom=3pt,valign=center]{\qrcode[height=1.9cm]{https://play.google.com/store/apps/details?id=com.luvvix.onbuch&hl=fr}}%
      \hspace{8pt}\begin{minipage}[c]{0.5\linewidth}
        {\popdisplay\fontsize{11}{12.5}\selectfont\color{white}\MakeUppercase{Télécharge OnBuch}}\par\vspace{4pt}
        {\raggedright\popsemi\fontsize{8.4}{11}\selectfont\color{white}Gratuit \textbullet\ Google Play\\Tuteur IA, annales, concours, résultats\par}
      \end{minipage}
    \end{tcolorbox}
  \end{minipage}\hfill
  \begin{minipage}[t]{0.487\linewidth}
    \begin{tcolorbox}[enhanced,colback=poppurple,colframe=popink,boxrule=2.2pt,arc=10pt,
        drop shadow=popink,left=11pt,right=11pt,top=10pt,bottom=10pt,height=3.1cm,valign=center]
      \tikz[baseline=-0.7ex]\node[circle,fill=white,draw=popink,line width=1.3pt,minimum size=1.6cm,inner sep=0]{\popdisplay\fontsize{24}{24}\selectfont\color{poppurple}+};%
      \hspace{8pt}\begin{minipage}[c]{0.6\linewidth}
        {\popdisplay\fontsize{11}{12.5}\selectfont\color{white}\MakeUppercase{Passe à OnBuch+}}\par\vspace{4pt}
        {\raggedright\popsemi\fontsize{8.4}{11}\selectfont\color{white}Tous les cours signés, Tuteur IA illimité, fiches PDF — onglet OnBuch+ de l'app\par}
      \end{minipage}
    \end{tcolorbox}
  \end{minipage}
  \vspace{8pt}
  \popcontenu
  \vfill
  \signatureonbuch
  \restoregeometry
  \newpage
}

% Rangée « Ce cours contient » (page de garde)
\newcommand{\poptile}[3]{% #1 couleur #2 gros texte #3 légende
  \begin{tcolorbox}[enhanced,colback=#1,colframe=popink,boxrule=2pt,arc=9pt,shadow={2.5pt}{-2.5pt}{0pt}{popink},
      left=6pt,right=6pt,top=9pt,bottom=9pt,halign=center,valign=center,height=1.75cm,width=\linewidth,nobeforeafter]
    {\popdisplay\fontsize{12.5}{13}\selectfont\color{white}\MakeUppercase{#2}}\par\vspace{3pt}
    {\centering\popsemi\fontsize{7.8}{9.5}\selectfont\color{white}#3\par}
  \end{tcolorbox}}
\newcommand{\popcontenu}{%
  \par\noindent{\popxboldls\fontsize{7.5}{7.5}\selectfont\color{popmuted}CE COURS SIGNÉ CONTIENT}\par\vspace{7pt}
  \noindent\begin{minipage}[t]{0.235\linewidth}\poptile{popblue}{Cours}{complet, illustré, pas à pas}\end{minipage}\hfill
  \begin{minipage}[t]{0.235\linewidth}\poptile{poporange}{Méthodes}{et exercices résolus}\end{minipage}\hfill
  \begin{minipage}[t]{0.235\linewidth}\poptile{poppink}{Exercices}{type examen + corrigés}\end{minipage}\hfill
  \begin{minipage}[t]{0.235\linewidth}\poptile{popgreen}{Fiche bilan}{l'essentiel à retenir}\end{minipage}\par}

% Sommaire
\newtcolorbox{sommaire}{enhanced,colback=white,colframe=popink,boxrule=2.2pt,arc=10pt,
  drop shadow=popink,left=18pt,right=18pt,top=13pt,bottom=13pt,before skip=6pt,after skip=20pt,
  before upper={\poppill{Sommaire}{poppurple}{white}\par\vspace{9pt}\popsemi\fontsize{10.6}{14.5}\selectfont\color{popink}}}

% Signature de fin de cours — poussée tout en bas de la dernière page
\newcommand{\findecours}{%
  \par\vfill\nopagebreak
  \begin{center}\popsquiggle{poporange}\end{center}\vspace{4pt}
  \signatureonbuch
}
\AtBeginDocument{\def\llbracket{\mathopen{[\mkern-3.5mu[}}\def\rrbracket{\mathclose{]\mkern-3.5mu]}}} % intervalles d'entiers (glyphe ⟦ absent de la police)
% Glyphes absents de Fira Math : redéfinis à partir de glyphes disponibles
\AtBeginDocument{%
  \def\setminus{\mathbin{\backslash}}\let\smallsetminus\setminus
  \def\vdots{\mathord{\vbox{\baselineskip3.2pt\lineskiplimit0pt\kern4pt\hbox{.}\hbox{.}\hbox{.}}}}%
  \def\ddots{\mathinner{\mkern1mu\raise6.5pt\hbox{.}\mkern2mu\raise3.6pt\hbox{.}\mkern2mu\raise0.7pt\hbox{.}\mkern1mu}}%
  \def\triangle{\mathord{\scalebox{0.9}{$\symup{\Delta}$}}}%
  \def\nabla{\mathord{\raisebox{\depth}{\scalebox{1}[-1]{$\symup{\Delta}$}}}}%
  \def\checkmark{\popcheck{popdark}}%
  \def\star{\mathbin{\ast}}\def\bigstar{\mathord{\ast}}%
  \def\diamond{\mathbin{\scalebox{0.6}{$\Diamond$}}}%
  \def\bigcup{\mathop{\mathchoice{\raisebox{-0.3em}{\scalebox{1.7}{$\cup$}}}{\scalebox{1.25}{$\cup$}}{\cup}{\cup}}\displaylimits}%
  \def\bigcap{\mathop{\mathchoice{\raisebox{-0.3em}{\scalebox{1.7}{$\cap$}}}{\scalebox{1.25}{$\cap$}}{\cap}{\cap}}\displaylimits}%
  \def\lesseqqgtr{\mathrel{\lessgtr}}\def\bigoplus{\mathop{\oplus}\displaylimits}%
}
\providecommand{\ssi}{si et seulement si} % abréviation fréquente
\AtBeginDocument{\providecommand{\wideparen}{\overparen}} % arc de cercle

%% FIGFIX-BEGIN v3 — correctifs globaux des figures (docs/onbuch-generation/tools/figfix)
\usepackage{adjustbox}\usepackage{etoolbox}
% toute figure TikZ/pgfplots plus large que la ligne est réduite à la largeur disponible
\BeforeBeginEnvironment{tikzpicture}{\begin{adjustbox}{max width=\linewidth}}
\AfterEndEnvironment{tikzpicture}{\end{adjustbox}}
% légendes pgfplots lisibles par-dessus les courbes
\pgfplotsset{every axis/.append style={legend style={fill=white,fill opacity=0.92,text opacity=1,draw=black!15,inner sep=3pt}}}
% étiquettes d'arêtes (syntaxe edge["..."]) avec fond blanc
\tikzset{every edge quotes/.append style={fill=white,inner sep=1.2pt}}
% bulles de cartes mentales : texte à la ligne dans la bulle, bulle un peu plus grande
\tikzset{every concept/.append style={text width=2.0cm,align=center,minimum size=2.3cm,inner sep=2pt}}
%% FIGFIX-END
\begin{document}

\pagegarde{Lesson 19 — Reading: Texts on helping with fund-raising activities/support groups in the acquisition of community resources/services}{Anglais}{Tle A}{Chapter 2 : Economic Life and Occupations}{EN19}{2 h}{Cette leçon de compréhension écrite (Reading) porte sur les textes décrivant les activités de collecte de fonds (fund-raising) et le rôle des groupes de soutien dans l'acquisition de ressources et services communautaires. Elle développe les stratégies de lecture (skimming, scanning, inférence) et le lexique thématique économique et social.
\vspace{4pt}
\begin{multicols}{2}\small
\begin{itemize}
\item À la fin de cette leçon, je sais identifier le thème principal et la structure d'un texte sur la collecte de fonds communautaire (skimming).
\item Je sais localiser des informations précises : montants, organisateurs, bénéficiaires, méthodes (scanning).
\item Je sais déduire le sens de mots inconnus grâce au contexte (indices lexicaux, morphologie).
\item Je sais maîtriser le vocabulaire clé : fund-raising, grant, sponsorship, donation, community resources, outreach.
\item Je sais répondre à des questions de compréhension globale, détaillée et d'interprétation (Vrai/Faux, QCM, Réponses courtes).
\item Je sais résumer le texte en 3-4 phrases en reprenant les idées forces.
\end{itemize}\end{multicols}}

\begin{sommaire}
\begin{multicols}{2}
\begin{enumerate}
\item 1. Découverte du texte support \& Stratégies de lecture
\item 2. Compréhension du texte : Questions variées
\item 3. Lexique Thématique : Fund-raising \& Community Resources
\item 4. Grammaire en contexte : Discours rapporté \& Passif (Révision Bac)
\item 5. Production Guidée : Écriture d'un Appel aux Dons (Fund-raising Appeal Letter)
\item 6. Production Orale : Simulation de Réunion Communautaire (Role-play)
\item Activité d'intégration
\item Méthodes et exercices résolus
\item Exercices
\item Corrigés des exercices
\item Fiche bilan
\end{enumerate}
\end{multicols}
\end{sommaire}

\begin{objectifs}
\begin{itemize}
\item À la fin de cette leçon, je sais identifier le thème principal et la structure d'un texte sur la collecte de fonds communautaire (skimming).
\item Je sais localiser des informations précises : montants, organisateurs, bénéficiaires, méthodes (scanning).
\item Je sais déduire le sens de mots inconnus grâce au contexte (indices lexicaux, morphologie).
\item Je sais maîtriser le vocabulaire clé : fund-raising, grant, sponsorship, donation, community resources, outreach.
\item Je sais répondre à des questions de compréhension globale, détaillée et d'interprétation (Vrai/Faux, QCM, Réponses courtes).
\item Je sais résumer le texte en 3-4 phrases en reprenant les idées forces.
\item Je sais exprimer mon opinion personnelle sur l'efficacité de ces actions dans mon contexte camerounais.
\item Je sais rédiger un court appel aux dons ou un rapport d'activité en utilisant les connecteurs logiques appropriés.
\end{itemize}
\end{objectifs}

\begin{prerequis}
\begin{itemize}
\item Maîtrise des temps du récit (Past Simple, Present Perfect) pour lire des comptes-rendus d'activités.
\item Connaissance du vocabulaire de base sur l'économie (money, bank, project, development) vu en Première.
\item Utilisation des connecteurs logiques (addition, cause, consequence, contrast).
\item Techniques de base de la lecture rapide (repérage mots-clés, titres, images).
\item Culture générale sur la vie associative au Cameroun (GIC, associations de développement, tontines).
\end{itemize}
\end{prerequis}

\newpage

\coursec{Découverte du texte support \& Stratégies de lecture}

Imaginez la scène : dans votre quartier à Yaoundé ou à Bamenda, l'école primaire du coin a besoin de nouveaux bancs, et le centre de santé intégré manque cruellement de médicaments. La mairie, sollicitée de toutes parts, ne peut pas tout financer seule. Alors que faire ? Attendre ? Non. Les habitants se réunissent un dimanche après-midi sous le manguier du chef de quartier. L'association des parents d'élèves propose une tombola, les femmes de la \eng{tontine} promettent une contribution, un jeune diplômé propose de lancer une campagne sur les réseaux sociaux, et un commerçant offre gratuitement des sacs de ciment. En quelques mois, grâce à ces \eng{fund-raising activities} (activités de collecte de fonds) et à ces \eng{support groups} (groupes de soutien), la communauté acquiert des \eng{community resources} (ressources communautaires) : des bancs, des médicaments, parfois même une salle polyvalente ou un forage d'eau potable.

C'est exactement le monde dans lequel cette leçon vous plonge : celui des initiatives citoyennes de financement participatif, racontées dans des textes authentiques — articles de blog d'ONG, rapports communautaires, appels aux dons. La problématique qui guidera notre lecture est la suivante : \emph{comment lire efficacement un texte anglais sur la collecte de fonds pour en extraire rapidement l'idée générale, les informations précises et l'intention de l'auteur ?}

\courssub{Le texte support : un article de blog d'ONG locale}

Avant toute explication, lisons. Voici un article original, tel qu'il pourrait paraître sur le blog d'une organisation non gouvernementale camerounaise. Lisez-le une première fois \emph{sans dictionnaire}, simplement pour le plaisir de comprendre.

\begin{textebox}[Building a Brighter Future: The Ngoma Community Hall Project]
For years, the villagers of Ngoma, a rural community in the Northwest Region, dreamed of a multipurpose hall. The old mission building had collapsed during the heavy rains of 2021, leaving the women's cooperative without a meeting place and the youth without a venue for skills training.

In January 2023, the Ngoma Development Association (NDA) launched a fund-raising campaign. Their strategy was threefold: first, a “Buy-a-Brick” campaign in which each family contributed 1,000 FCFA for a symbolic brick; second, a cultural festival featuring local dance troupes and a football tournament, which raised 2.5 million FCFA in gate fees and sponsorships from local businesses such as Ngoma Agro Ltd; third, a grant application to the Council and to an international NGO, Rural Aid International.

By December 2023, the roof was on. The hall now hosts adult literacy classes and a weekly market for farm produce, and it also serves as a vaccination centre. “This hall is our bank,” says Mama Agnes, president of the women's group. “We deposit our efforts and withdraw development.” The NDA is now raising funds for solar panels to ensure a sustainable energy supply.

\emph{Adapted from a community report, 2024.}
\end{textebox}

\begin{definition}[Glossaire du texte]
\begin{itemize}
\item \cle{villager} (n.) : villageois, habitant d'un village.
\item \cle{multipurpose hall} (n.) : salle polyvalente.
\item \cle{to collapse} (v.) : s'effondrer, s'écrouler.
\item \cle{venue} (n.) : lieu (d'un événement), endroit prévu pour une activité.
\item \cle{skills training} (n.) : formation professionnelle, apprentissage de compétences.
\item \cle{to launch} (v.) : lancer (une campagne, un projet).
\item \cle{threefold} (adj.) : triple, en trois volets.
\item \cle{brick} (n.) : brique.
\item \cle{gate fees} (n.) : recettes des entrées (billetterie).
\item \cle{sponsorship} (n.) : parrainage, mécénat.
\item \cle{grant} (n.) : subvention (argent donné, non remboursable).
\item \cle{literacy classes} (n.) : cours d'alphabétisation.
\item \cle{farm produce} (n.) : produits agricoles.
\item \cle{solar panels} (n.) : panneaux solaires.
\item \cle{sustainable} (adj.) : durable.
\end{itemize}
\end{definition}

\courssub{Stratégie 1 : le skimming, ou l'art de survoler un texte}

Observez ce que vous venez de faire spontanément : avant de lire chaque mot, vos yeux ont d'abord croisé le titre, puis les premières phrases des paragraphes. C'est exactement la stratégie du \cle{skimming}.

\begin{propriete}[Le skimming : règle et emploi]
Le \cle{skimming} (lecture rapide, « écrémage ») consiste à lire un texte \emph{en diagonale} pour en saisir l'idée générale (\eng{the gist}), sans s'arrêter sur les détails ni sur les mots inconnus. On lit uniquement :
\begin{enumerate}
\item le \textbf{titre} et les sous-titres ;
\item la \textbf{première phrase} de chaque paragraphe (souvent la phrase qui introduit le sujet du paragraphe, \eng{topic sentence}) ;
\item la \textbf{dernière phrase} de chaque paragraphe (souvent la conclusion ou la transition) ;
\item les mots en \textbf{gras} ou en \emph{italique}, les chiffres et les noms propres qui sautent aux yeux.
\end{enumerate}
Question-clé du skimming : \eng{What is the text mainly about?} (« De quoi parle principalement le texte ? »)
\end{propriete}

Appliquons cette stratégie au texte de Ngoma :

\begin{exemplebox}[Skimming appliqué au texte]
\begin{itemize}
\item \textbf{Titre :} \eng{Building a Brighter Future: The Ngoma Community Hall Project} → le texte parle d'un \emph{projet de salle communautaire} à Ngoma.
\item \textbf{Première phrase du §1 :} \eng{For years, the villagers of Ngoma dreamed of a multipurpose hall.} → un rêve, un besoin ancien.
\item \textbf{Première phrase du §2 :} \eng{In January 2023, the NDA launched a fund-raising campaign.} → une campagne de collecte de fonds est lancée.
\item \textbf{Première phrase du §3 :} \eng{By December 2023, the roof was on.} → le projet a abouti.
\item \textbf{Conclusion du skimming :} le texte raconte comment une communauté rurale a collecté des fonds pour construire une salle polyvalente. Idée générale acquise en moins de deux minutes !
\end{itemize}
\end{exemplebox}

\courssub{Stratégie 2 : le scanning, ou la recherche ciblée}

Le \cle{scanning} est l'inverse du skimming : on ne cherche plus l'idée générale, mais une \emph{information précise} — un chiffre, une date, un nom — comme un aigle qui repère sa proie.

\begin{propriete}[Le scanning : règle et emploi]
Le \cle{scanning} (lecture sélective) consiste à parcourir le texte \emph{les questions en main}, en cherchant uniquement les mots-clés correspondants : chiffres, dates, noms propres (faciles à repérer grâce à la majuscule), montants en FCFA, verbes d'action. On ignore tout le reste. Questions-clés : \eng{How much? When? Who? Where?} (« Combien ? Quand ? Qui ? Où ? »)
\end{propriete}

\begin{exemplebox}[Scanning appliqué au texte]
\begin{itemize}
\item \eng{When did the old building collapse?} → on scanne les dates : \textbf{2021} (\eng{the heavy rains of 2021}).
\item \eng{When was the campaign launched?} → \textbf{January 2023}.
\item \eng{How much did each family contribute?} → on scanne les montants : \textbf{1,000 FCFA}.
\item \eng{How much did the festival raise?} → \textbf{2.5 million FCFA}.
\item \eng{Who sponsored the festival?} → on scanne les noms propres : \textbf{Ngoma Agro Ltd}.
\item \eng{Which organisations received the grant application?} → \textbf{the Council} and \textbf{Rural Aid International}.
\end{itemize}
\end{exemplebox}

\begin{popfigure}
\begin{tikzpicture}
\node[draw, rounded corners, fill=popblueL, align=center] (title) at (0,0) {\textbf{Skimming vs Scanning}};
\node[draw, below left=1.2cm and -0.5cm of title, fill=popgreenL, align=center] (skim) {\emph{Skimming :} vue d'ensemble\\ « Quelle est l'idée principale ? »};
\node[draw, below right=1.2cm and -0.5cm of title, fill=poporangeL, align=center] (scan) {\emph{Scanning :} recherche ciblée\\ « Combien ? Quand ? Qui ? »};
\draw[->, thick] (title) -- (skim);
\draw[->, thick] (title) -- (scan);
\end{tikzpicture}
\legende{Les deux stratégies complémentaires de lecture rapide : le skimming pour l'idée générale, le scanning pour les détails précis.}
\end{popfigure}

\begin{center}
\begin{tabularx}{\linewidth}{|X|X|X|}
\hline
\rowcolor{popblueL} Critère & Skimming & Scanning \\
\hline
Objectif & Idée générale (\eng{gist}) & Information précise (\eng{detail}) \\
\hline
Ce qu'on lit & Titre, 1\textsuperscript{re}/dernière phrase de chaque § & Mots-clés : chiffres, dates, noms propres \\
\hline
Vitesse & Très rapide (2 min) & Rapide et sélective (3 min) \\
\hline
Question typique & \eng{What is the text about?} & \eng{How much was raised?} \\
\hline
Équivalent français & Survoler, lire en diagonale & Repérer, balayer du regard \\
\hline
\end{tabularx}
\end{center}

\courssub{Stratégie 3 : déduire le sens des mots inconnus}

Au Baccalauréat, le dictionnaire est interdit. Il faut donc deviner le sens des mots difficiles grâce au \cle{contexte}. Trois indices vous aident :

\begin{enumerate}
\item \textbf{La famille du mot :} \eng{fund-raising} contient \eng{fund} (fonds) + \eng{raising} (lever, collecter) → « collecte de fonds ».
\item \textbf{Les explications implicites :} \eng{a “Buy-a-Brick” campaign in which each family contributed 1,000 FCFA} → le mot \eng{contributed} est éclairé par le contexte : les familles \emph{donnent} de l'argent → « contribuer ».
\item \textbf{La logique de la phrase :} \eng{The old mission building had collapsed during the heavy rains} → sous de fortes pluies, un vieux bâtiment ne peut que « s'effondrer ».
\end{enumerate}

\begin{methode}[La méthode des « 3 Lectures » pour le Bac]
\begin{enumerate}
\item \textbf{Lecture 1 (Skimming — 2 min) :} titre + première et dernière phrase de chaque paragraphe = idée générale. Notez-la en une phrase au brouillon.
\item \textbf{Lecture 2 (Scanning — 3 min) :} questions en main, cherchez les mots-clés (chiffres, noms propres, verbes d'action) et surlignez les passages qui contiennent les réponses.
\item \textbf{Lecture 3 (Détaillée — 5 min) :} relisez pour vérifier vos réponses, devinez le vocabulaire difficile grâce au contexte, et repérez les nuances d'opinion de l'auteur.
\end{enumerate}
\end{methode}

\courssub{Identifier la structure du texte}

Un texte de collecte de fonds suit presque toujours la même architecture, qu'il faut savoir reconnaître pour répondre aux questions du type \eng{What is the function of paragraph 2?}

\begin{center}
\begin{tabularx}{\linewidth}{|X|X|X|}
\hline
\rowcolor{popblueL} Partie & Fonction & Dans notre texte \\
\hline
Introduction & Contexte et problème & §1 : le vieux bâtiment s'est effondré, la communauté n'a plus de salle \\
\hline
Actions & Événements et démarches & §2 : campagne « Buy-a-Brick », festival, demande de subvention \\
\hline
Résultats / Impact & Ce qui a été obtenu & §3 : toit posé, cours d'alphabétisation, marché, vaccinations \\
\hline
Appel final & Incitation à agir & §3 fin : nouvelle collecte pour les panneaux solaires \\
\hline
\end{tabularx}
\end{center}

\courssub{Analyser le ton et le public visé}

Deux questions essentielles au Bac : \eng{What is the writer's purpose?} (« Quel est le but de l'auteur ? ») et \eng{Who is the text addressed to?} (« À qui le texte s'adresse-t-il ? »)

\begin{aretenir}[Ton et public du texte]
\begin{itemize}
\item \textbf{Ton informatif :} l'auteur donne des faits, des dates et des chiffres précis (\eng{January 2023}, \eng{2.5 million FCFA}).
\item \textbf{Ton persuasif :} la citation émotive de Mama Agnes (\eng{“This hall is our bank”}) et l'annonce d'une nouvelle collecte visent à toucher le lecteur et à l'inciter à donner.
\item \textbf{Public visé :} les donateurs potentiels (\eng{potential donors}), les membres de la communauté et les partenaires (conseil municipal, ONG internationales).
\end{itemize}
Indices linguistiques du ton persuasif : adjectifs positifs (\eng{brighter}, \eng{sustainable}), pronoms inclusifs (\eng{our}, \eng{we}), citations directes, verbes d'action (\eng{contribute}, \eng{support}, \eng{join}).
\end{aretenir}

\begin{attention}[Faux amis et pièges de traduction]
\begin{itemize}
\item Ne confondez pas \eng{fund-raising} (nom composé : « collecte de fonds ») avec une forme incorrecte comme \eng{rising funds}. On dit \eng{a fund-raising campaign} = « une campagne de collecte de fonds ». Ne traduisez jamais par « lever des fonds » (anglicisme syntaxique) ; préférez « collecte de fonds » ou « recherche de financement ».
\item \eng{Grant} = subvention, don \emph{non remboursable} ; \eng{loan} = prêt, \emph{remboursable}. \eng{The NDA applied for a grant, not a loan.} « La NDA a demandé une subvention, pas un prêt. »
\item \eng{To raise money} = collecter de l'argent ; \eng{to rise} (rise, rose, risen) = s'élever, se lever — verbe \emph{intransitif} qui ne prend jamais de complément d'objet. On dit \eng{They raised 2.5 million FCFA}, jamais \eng{They rised money} ni \eng{They rose money}.
\item \eng{Actually} ne veut pas dire « actuellement » mais « en fait » ; « actuellement » se dit \eng{currently} ou \eng{at present}.
\item \eng{To assist} signifie « aider » ; pour « assister à » (une réunion), on dit \eng{to attend}. \eng{Many villagers attended the fund-raising meeting.} « De nombreux villageois ont assisté à la réunion de collecte de fonds. »
\end{itemize}
\end{attention}

\begin{exoresolu}[Questions de compréhension sur le texte]
Répondez en anglais, en phrases complètes : (a) What problem did the villagers of Ngoma face? (b) Mention the three fund-raising strategies used by the NDA. (c) How much money did the cultural festival raise? (d) What is the hall used for today? (e) What is the NDA's next project?
\tcblower
\textbf{Corrigé détaillé :}
\begin{itemize}
\item (a) \eng{The old mission building had collapsed during the heavy rains of 2021, so the community had no meeting place and no venue for skills training.} (Réponse trouvée au §1 par scanning de la date et du problème.)
\item (b) \eng{The NDA used a “Buy-a-Brick” campaign, a cultural festival with a football tournament, and a grant application to the Council and to Rural Aid International.} (§2 : repérer les marqueurs \eng{first}, \eng{second}, \eng{third}.)
\item (c) \eng{The festival raised 2.5 million FCFA in gate fees and sponsorships.} (Scanning du montant.)
\item (d) \eng{Today, the hall hosts adult literacy classes and a weekly market for farm produce, and it also serves as a vaccination centre.} (§3.)
\item (e) \eng{The NDA is now raising funds for solar panels to ensure a sustainable energy supply.} (Dernière phrase : l'appel final.)
\end{itemize}
\end{exoresolu}

\begin{exercice}[À vous de jouer]
Relisez le texte en appliquant la méthode des 3 lectures, puis : (1) résumez le texte en deux phrases en anglais (\eng{The text is about... It explains how...}) ; (2) trouvez dans le texte un mot ou une expression signifiant : « s'effondrer », « subvention », « panneaux solaires », « parrainage » ; (3) dites en une phrase si le ton du texte est plutôt informatif, persuasif ou les deux, et justifiez avec une citation du texte.
\end{exercice}

\begin{corrige}[Exercice : éléments de réponse]
(1) \eng{The text is about a fund-raising campaign in Ngoma, a village in the Northwest Region. It explains how the community built a multipurpose hall thanks to a “Buy-a-Brick” campaign, a cultural festival and grants.} (2) « s'effondrer » = \eng{to collapse} ; « subvention » = \eng{grant} ; « panneaux solaires » = \eng{solar panels} ; « parrainage » = \eng{sponsorship}. (3) Le ton est les deux : informatif (\eng{“which raised 2.5 million FCFA in gate fees”}) et persuasif (\eng{“This hall is our bank... We deposit our efforts and withdraw development.”}).
\end{corrige}

\begin{savaistu}[La collecte de fonds, une tradition anglophone]
Dans les pays anglophones, la collecte de fonds communautaire est une véritable institution : les écoles britanniques organisent des \eng{school fairs} (fêtes scolaires), les Américains des \eng{bake sales} (ventes de gâteaux) et des \eng{charity runs} (courses caritatives). Au Cameroun anglophone, on retrouve la même énergie dans les \eng{contributions} de village, les \eng{njangi} (tontines) et les journées de \eng{community work}. Le mot \eng{fund-raising} lui-même s'est répandu au XX\textsuperscript{e} siècle avec le développement des grandes organisations caritatives comme la Croix-Rouge.
\end{savaistu}

\coursec{Compréhension du texte : Questions variées}

Dans la section précédente, nous avons découvert le texte support — l'histoire de la \eng{Nkwen Development Association (NDA)} et de sa campagne de collecte de fonds pour rénover le \eng{community hall} de Bamenda — et nous avons appliqué les stratégies de lecture (\emph{skimming} pour l'idée générale, \emph{scanning} pour les détails, déduction du sens des mots par le contexte). Il est temps maintenant de vérifier notre compréhension en profondeur, exactement comme à l'épreuve d'anglais du Baccalauréat : à travers cinq types d'exercices classiques.

Pour travailler sereinement, relisons d'abord le texte support dans sa version complète.

\begin{textebox}[Reading text — “A Roof Over Our Heads: How Nkwen Rebuilt Its Community Hall”]
When the roof of the Nkwen community hall collapsed during the heavy rains of September 2022, the quarter lost more than a building: it lost its meeting place, its polling station and its celebration ground. The Nkwen Development Association (NDA) decided to act immediately rather than wait for help that might never come.

In January 2023, the NDA launched a threefold fund-raising strategy. First, the “Buy-a-Brick” campaign invited every villager, including those living in Douala and abroad, to sponsor bricks at 1,000 CFA francs each. Second, a cultural festival with a football tournament was organised during the Easter holidays; gate fees and food stalls brought in nearly two million francs. Third, the association applied for grants from two NGOs and the regional council.

Mama Agnes, president of the women's cooperative, played a decisive role. Her group contributed 500,000 francs from their savings and cooked for the volunteers who worked on the site every weekend. “Nobody paid us,” she says proudly. “We did it for our children.”

Not everything was easy. Some villagers complained that the money was being mismanaged, so the NDA published its accounts on the church notice board every month. This transparency restored confidence and donations increased again.

By December 2023, the new roof was finished. The hall now hosts a nursery school in the morning, adult literacy classes in the evening, and weddings on Saturdays. “The hall belongs to everybody,” the NDA chairman declared at the inauguration, “because everybody built it.”
\end{textebox}

\textbf{Glossaire :}

\begin{center}
\begin{tabularx}{\linewidth}{|X|X|X|}\hline
\rowcolor{popblueL} Mot anglais & Nature & Traduction française \\ \hline
roof & nom & toit, toiture \\ \hline
to collapse & verbe & s'effondrer \\ \hline
polling station & nom & bureau de vote \\ \hline
threefold & adjectif & triple, en trois volets \\ \hline
fund-raising & nom / adjectif & collecte de fonds \\ \hline
to sponsor & verbe & parrainer, financer \\ \hline
gate fees & nom & prix d'entrée (billetterie) \\ \hline
grant & nom & subvention \\ \hline
to mismanage & verbe & mal gérer, détourner (des fonds) \\ \hline
notice board & nom & panneau d'affichage \\ \hline
literacy classes & nom & cours d'alphabétisation \\ \hline
\end{tabularx}
\end{center}

\courssub{Les types de questions au Bac : savoir où chercher}

Avant de répondre, il faut comprendre \emph{ce que la question attend}. Toutes les questions de compréhension se rangent en quatre familles.

\begin{propriete}[Types de questions Bac \& astuces]
\begin{itemize}
\item \textbf{Explicit (« Right there ») :} la réponse est littéralement dans le texte, avec les mêmes mots ou des synonymes. Stratégie : repérer le mot-clé de la question, puis faire du \emph{scanning}.
\item \textbf{Implicit (« Think \& Search ») :} il faut combiner des informations de plusieurs phrases ou paragraphes. Stratégie : souligner tous les passages liés avant de formuler la réponse.
\item \textbf{Vocabulary in context :} pour trouver le sens d'un mot, remplacez-le par chacune des propositions ; celle qui conserve le sens de la phrase est la bonne.
\item \textbf{Reference words :} pour \eng{“their strategy”}, \eng{“this hall”}, \eng{“it”}, identifiez l'antécédent (ici : la NDA, le hall). Cherchez toujours le nom \emph{avant} le pronom.
\end{itemize}
\end{propriete}

\begin{attention}[Piège fréquent des francophones]
Ne répondez jamais d'après vos connaissances générales ou votre expérience personnelle : répondez \textbf{uniquement d'après le texte}. Si le texte dit que le toit s'est effondré à cause des pluies, une réponse évoquant un tremblement de terre est fausse, même si elle semble « plausible ». Autre piège : recopier une longue phrase du texte qui ne répond pas à la question posée. Citez court, citez juste.
\end{attention}

\courssub{Exercice Type 1 : Vrai / Faux / Non mentionné (True / False / Not stated)}

\textbf{Méthode.} Pour chaque affirmation : (1) repérez le passage correspondant ; (2) comparez mot à mot ; (3) écrivez \textbf{TRUE}, \textbf{FALSE} ou \textbf{NOT STATED}, puis \textbf{justifiez toujours par une courte citation} entre guillemets anglais. « Non mentionné » signifie que le texte ne dit \emph{rien} sur le sujet — ni vrai ni faux.

\begin{exercice}[True, False or Not stated? Justify with a quotation]
\begin{enumerate}
\item The old building was destroyed by an earthquake.
\item The “Buy-a-Brick” campaign also concerned villagers living outside Nkwen.
\item The cultural festival raised exactly three million francs.
\item Mama Agnes was paid for cooking for the volunteers.
\item The NDA published its accounts every month.
\item The new hall is used by a secondary school.
\end{enumerate}
\end{exercice}

\begin{corrige}[Exercice Type 1]
\begin{enumerate}
\item \textbf{FALSE.} Text: “the roof ... collapsed during the heavy rains of September 2022” (pluies, pas tremblement de terre).
\item \textbf{TRUE.} Text: “including those living in Douala and abroad”.
\item \textbf{FALSE.} Text: “brought in nearly two million francs” (\eng{nearly} = presque ; ni exactement, ni trois millions).
\item \textbf{FALSE.} Text: “Nobody paid us.”
\item \textbf{TRUE.} Text: “published its accounts on the church notice board every month”.
\item \textbf{NOT STATED.} Le texte mentionne une \eng{nursery school} (maternelle) et des cours pour adultes, mais aucun collège ou lycée.
\end{enumerate}
\end{corrige}

\courssub{Exercice Type 2 : QCM — vocabulaire en contexte et inférence}

Le QCM teste deux choses : le sens d'un mot dans son contexte, et votre capacité à \emph{déduire} (inférer) ce que le texte suggère sans le dire.

\begin{exemplebox}[QCM commenté]
\eng{What does “threefold” mean in paragraph 2?}\par
A) Expensive \quad B) In three parts \quad C) Difficult \quad D) Quick\par
\textbf{Réponse : B} — « en trois volets ». Indice contextuel : le paragraphe énumère \eng{“First ... Second ... Third ...”}. Le préfixe \emph{three-} et le suffixe \emph{-fold} (comme \emph{twofold}) confirment le sens.
\end{exemplebox}

\begin{exercice}[Multiple Choice Questions]
\begin{enumerate}
\item \eng{“Gate fees” (paragraph 2) refers to:}
\begin{enumerate}
\item the price of building a gate
\item money paid by visitors to enter the festival
\item taxes paid to the regional council
\item the cost of food at the stalls
\end{enumerate}
\item \eng{Why did donations increase again? (inference)}
\begin{enumerate}
\item because the NGOs gave more money
\item because people trusted the association after seeing the accounts
\item because the roof was already finished
\item because Mama Agnes asked the church for help
\end{enumerate}
\item \eng{“This transparency” (paragraph 4) refers to:}
\begin{enumerate}
\item the new roof
\item the publication of the accounts
\item the football tournament
\item the women's savings
\end{enumerate}
\end{enumerate}
\end{exercice}

\begin{corrige}[Exercice Type 2]
\begin{enumerate}
\item \textbf{b} — les \eng{gate fees} sont la billetterie, l'argent payé à l'entrée du festival.
\item \textbf{b} — inférence : le texte enchaîne “This transparency restored confidence and donations increased again” ; la confiance restaurée explique la hausse des dons.
\item \textbf{b} — mot de référence : \eng{this transparency} renvoie à la phrase précédente, la publication mensuelle des comptes.
\end{enumerate}
\end{corrige}

\courssub{Exercice Type 3 : Questions à réponses courtes (Short answers)}

\textbf{Méthode.} Répondez par une phrase courte et complète, en anglais, avec vos propres mots autant que possible. Une réponse d'un seul mot est acceptée si la question le permet (\eng{Who...?}, \eng{When...?}), mais elle doit être exacte.

\begin{exercice}[Answer in short complete sentences]
\begin{enumerate}
\item Who is Mama Agnes?
\item When did the NDA launch its fund-raising strategy?
\item How much did one brick cost in the “Buy-a-Brick” campaign?
\item Why did some villagers complain?
\item What activities take place in the new hall?
\end{enumerate}
\end{exercice}

\begin{corrige}[Exercice Type 3]
\begin{enumerate}
\item She is the president of the women's cooperative (women's group).
\item It launched the strategy in January 2023.
\item One brick cost 1,000 CFA francs.
\item They complained because they thought the money was being mismanaged.
\item The hall hosts a nursery school in the morning, adult literacy classes in the evening and weddings on Saturdays.
\end{enumerate}
\end{corrige}

\courssub{Exercice Type 4 : Réordonnancement chronologique (Timeline)}

Ce type d'exercice vérifie que vous savez suivre la trame temporelle d'un texte. Repérez les marqueurs de temps : \eng{in January 2023}, \eng{during the Easter holidays}, \eng{by December 2023}, \eng{first / second / third}.

\begin{exoresolu}[Put the events in chronological order]
Énoncé — Voici cinq événements dans le désordre. Remettez-les dans l'ordre chronologique (1 à 5) :\par
a) The NDA applied for grants from NGOs and the regional council.\par
b) The “Buy-a-Brick” campaign was launched.\par
c) The new roof was finished.\par
d) The roof of the old hall collapsed.\par
e) A cultural festival and football tournament were organised.
\tcblower
Solution détaillée — On suit les marqueurs temporels du texte :\par
1 → d) “collapsed during the heavy rains of September 2022” (tout commence par la catastrophe).\par
2 → b) “In January 2023, the NDA launched ... First, the ‘Buy-a-Brick’ campaign”.\par
3 → e) “Second, a cultural festival with a football tournament ... during the Easter holidays”.\par
4 → a) “Third, the association applied for grants”.\par
5 → c) “By December 2023, the new roof was finished.”\par
Ordre final : \textbf{d – b – e – a – c}.
\end{exoresolu}

\begin{popfigure}
\begin{center}
\begin{tikzpicture}[node distance=0.4cm]
\node[draw, rectangle, rounded corners, fill=popgoldL, align=center] (start) {Sept 2022\\Effondrement};
\node[draw, rectangle, rounded corners, right=of start, fill=popgoldL, align=center] (brick) {Jan 2023\\“Buy-a-Brick”};
\node[draw, rectangle, rounded corners, right=of brick, fill=popgoldL, align=center] (fest) {Pâques 2023\\Festival + Tournoi};
\node[draw, rectangle, rounded corners, right=of fest, fill=popgoldL, align=center] (grant) {2023\\Demandes de\\subventions};
\node[draw, rectangle, rounded corners, right=of grant, fill=popgreenL, align=center] (end) {Déc 2023\\Toiture finie};
\draw[->, thick, popdark] (start) -- (brick);
\draw[->, thick, popdark] (brick) -- (fest);
\draw[->, thick, popdark] (fest) -- (grant);
\draw[->, thick, popdark] (grant) -- (end);
\end{tikzpicture}
\end{center}
\legende{Frise chronologique de la campagne de collecte de fonds de la NDA (septembre 2022 – décembre 2023).}
\end{popfigure}

\courssub{Exercice Type 5 : Association idées forces / paragraphes (Matching)}

\textbf{Méthode.} Lisez d'abord toutes les idées forces proposées, puis relisez chaque paragraphe en vous demandant : « \emph{What is this paragraph mainly about?} » Méfiez-vous des détails qui attirent l'œil mais ne résument pas le paragraphe.

\begin{exercice}[Match each main idea with the correct paragraph (1–5)]
Idées forces :\par
A) The results and the new uses of the hall.\par
B) The disaster and the decision to act.\par
C) The contribution of the women's cooperative.\par
D) A problem of confidence and its solution.\par
E) The three fund-raising actions.
\end{exercice}

\begin{corrige}[Exercice Type 5]
\begin{itemize}
\item Paragraphe 1 → \textbf{B} (effondrement du toit, décision immédiate de la NDA).
\item Paragraphe 2 → \textbf{E} (la stratégie \eng{threefold} : briques, festival, subventions).
\item Paragraphe 3 → \textbf{C} (Mama Agnes et la coopérative des femmes).
\item Paragraphe 4 → \textbf{D} (plaintes de mauvaise gestion → transparence des comptes).
\item Paragraphe 5 → \textbf{A} (toiture terminée, nouvelles activités du hall).
\end{itemize}
\end{corrige}

\courssub{Bilan : la boîte à outils du lecteur}

\begin{center}
\begin{tabularx}{\linewidth}{|X|X|X|}\hline
\rowcolor{popblueL} Type de question & Ce qu'elle teste & Stratégie gagnante \\ \hline
True / False / Not stated & Compréhension littérale & Citer le texte entre guillemets \\ \hline
QCM vocabulaire & Sens en contexte & Remplacer le mot par chaque option \\ \hline
QCM inférence & Déduction & Combiner plusieurs phrases \\ \hline
Short answers & Compréhension détaillée & Phrase courte, mots propres \\ \hline
Timeline & Trame chronologique & Repérer les marqueurs de temps \\ \hline
Matching & Idées principales & Demander “What is it mainly about?” \\ \hline
\end{tabularx}
\end{center}

\begin{aretenir}[À retenir]
Toute réponse de compréhension doit pouvoir être \textbf{prouvée par le texte}. Citez court, citez exact, et distinguez toujours ce qui est dit (\emph{explicit}), ce qui est suggéré (\emph{implicit}) et ce qui n'est pas dit du tout (\emph{not stated}). Ces cinq formats d'exercices sont exactement ceux de l'épreuve de \emph{Reading Comprehension} du Baccalauréat : les maîtriser, c'est sécuriser des points.
\end{aretenir}

\coursec{Lexique Thématique : Fund-raising \& Community Resources}

Après avoir lu et questionné le texte support dans la section précédente, vous avez remarqué qu'il regorge de mots spécifiques : \eng{donors}, \eng{grant}, \eng{pledge}, \eng{outreach}\ldots{} Cette section les organise méthodiquement. Un bon lecteur ne mémorise pas des mots isolés : il les classe par \cle{champs lexicaux} (acteurs, actions, argent, ressources) et par \cle{collocations} (mots qui vont naturellement ensemble). C'est exactement ce que nous allons faire, en partant d'un court texte d'observation.

\courssub{Texte d'observation : un article de presse local}

Lisez ce texte une première fois sans dictionnaire. Soulignez mentalement tous les mots liés à l'argent et à l'organisation.

\begin{textebox}[The Bamenda Community Hall Project — a press report]
The Bamenda Hill Development Association has just launched a fund-raising campaign to build a new community hall. The organizers hope to raise 25 million FCFA before the end of the year. So far, local donors have pledged 8 million FCFA, and a microfinance institution has promised a grant of 5 million FCFA for vocational training equipment.

“We want full transparency,” said Mrs. Nfor, the project coordinator. “Every contribution will be recorded, and an independent audit will be published at the end of the campaign. Our beneficiaries are the 3,000 inhabitants of the quarter, especially young people who need literacy programs and skills training.”

The association is also counting on its members in the diaspora. A crowdfunding page has been opened online, and a charity auction and a raffle will take place during the Annual Convention in Buea. Volunteers will run an outreach program in remote villages to raise awareness about the project.

“This is not only about money,” added a sponsor from a local brewery. “It is about sustainable development: better amenities, better healthcare access, and a stronger community. If we mobilize our resources well, the campaign will yield results.”
\end{textebox}

\textbf{Glossaire (mots difficiles)}

\begin{center}
\begin{tabularx}{\linewidth}{|l|l|X|}
\hline
\rowcolor{popblueL} Mot anglais & Nature & Traduction française \\
\hline
to pledge & verbe & promettre (un don), s'engager à donner \\
\hline
grant & nom & subvention \\
\hline
audit & nom & audit, vérification des comptes \\
\hline
raffle & nom & tombola \\
\hline
auction & nom & vente aux enchères \\
\hline
outreach & nom & action de terrain, sensibilisation de proximité \\
\hline
amenities & nom (pluriel) & équipements collectifs, commodités \\
\hline
to yield results & expression & produire des résultats, porter ses fruits \\
\hline
\end{tabularx}
\end{center}

\textbf{Questions rapides (entraînement au scanning)}
\begin{enumerate}
\item How much money do the organizers hope to raise?
\item Who has pledged 8 million FCFA?
\item What will the grant of 5 million FCFA pay for?
\item Name two fund-raising events mentioned in the text.
\item What does the sponsor say the project is really about?
\end{enumerate}

\textit{Réponses : 1. 25 million FCFA. 2. Local donors. 3. Vocational training equipment. 4. A charity auction and a raffle (also a crowdfunding page). 5. Sustainable development: amenities, healthcare access, a stronger community.}

\courssub{Champ lexical 1 : Les acteurs (the people involved)}

\begin{definition}[Les acteurs d'une collecte de fonds]
\begin{itemize}
\item \cle{Organizers} (n. pl.) : les organisateurs. \eng{The organizers planned the campaign carefully.}
\item \cle{Donors} (n. pl.) : les donateurs, ceux qui donnent. \eng{Donors pledged 10 million FCFA.}
\item \cle{Beneficiaries} (n. pl.) : les bénéficiaires, ceux qui profitent du projet. \eng{The beneficiaries are the villagers.}
\item \cle{Volunteers} (n. pl.) : les bénévoles, qui travaillent gratuitement. \eng{Volunteers distributed the posters.}
\item \cle{Stakeholders} (n. pl.) : les parties prenantes (tous ceux qui ont un intérêt dans le projet). \eng{All stakeholders attended the meeting.}
\item \cle{Sponsors} (n. pl.) : les sponsors, parrains financiers (souvent des entreprises). \eng{A local bank became our main sponsor.}
\item \cle{Grant-makers} (n. pl.) : les organismes qui accordent des subventions (fondations, États, ONG).
\end{itemize}
\end{definition}

\begin{attention}[Volunteer ou « volontaire » ?]
Un \eng{volunteer} est un \textbf{bénévole} (qui travaille sans salaire). L'adjectif français « volontaire » (qui a de la bonne volonté) se dit \eng{willing} : \eng{She is willing to help.} « Elle est volontaire pour aider. » Ne traduisez jamais « bénévole » par \eng{voluntary worker} dans une copie de lycée : le mot attendu est \eng{volunteer}. Prononciation : « vo-lon-TIRE » (accent sur la dernière syllabe).
\end{attention}

\courssub{Champ lexical 2 : Les actions et les événements}

\begin{definition}[Fund-raising events and actions]
\begin{itemize}
\item \cle{Fund-raising campaign} : campagne de collecte de fonds. \eng{They launched a fund-raising campaign in March.}
\item \cle{Drive} : collecte (courte et ciblée). \eng{A blood drive} = une collecte de sang ; \eng{a food drive} = une collecte de denrées.
\item \cle{Rally} : rassemblement, meeting de mobilisation.
\item \cle{Auction} : vente aux enchères. \eng{The painting was sold at the charity auction.}
\item \cle{Raffle} : tombola (on vend des tickets, on tire au sort un gagnant).
\item \cle{Crowdfunding} : financement participatif (en ligne). \eng{They raised the money through crowdfunding.}
\item \cle{Sponsorship deal} : accord de parrainage (avec une entreprise).
\item \cle{Outreach} : action de terrain, « aller vers » les populations. \eng{The outreach program targets remote villages.}
\end{itemize}
\end{definition}

\begin{definition}[Grant, pledge, outreach : trois mots-clés]
\cle{Grant} (n.) : \textbf{subvention}. Somme d'argent accordée par une institution (gouvernement, ONG, fondation) pour un projet précis, sans obligation de remboursement. \eng{The NGO secured a \$50,000 grant.} « L'ONG a obtenu une subvention de 50 000 dollars. »\par
\cle{Pledge} (n. et v.) : \textbf{promesse de don / s'engager à donner}. \eng{Donors pledged 10 million FCFA.} « Les donateurs ont promis 10 millions de FCFA. »\par
\cle{Outreach} (n.) : \textbf{action de terrain}, sensibilisation de proximité. \eng{The outreach program targets remote villages.} « Le programme de terrain cible les villages reculés. »
\end{definition}

\begin{attention}[« Subvention » n'est pas \eng{subvention} !]
Le mot \eng{subvention} existe en anglais mais est rare et soutenu ; le mot courant est \cle{grant}. De même, une « collecte » se dit \eng{fund-raising} ou \eng{drive}, jamais \eng{a collection} dans ce contexte (\eng{collection} = collection d'objets). Enfin, « parrainer » un projet = \eng{to sponsor}, et le « parrainage » = \eng{sponsorship}.
\end{attention}

\courssub{Champ lexical 3 : L'argent et la gestion}

\begin{definition}[Money and management]
\begin{itemize}
\item \cle{Budget} : budget. \eng{The project budget is 30 million FCFA.}
\item \cle{Donation} : don (l'argent donné). \eng{Every donation, big or small, counts.}
\item \cle{Contribution} : contribution, apport. \eng{Thank you for your generous contribution.}
\item \cle{Revenue} : revenus, recettes. \eng{The revenue from the raffle reached 2 million FCFA.}
\item \cle{Expenses} : dépenses. \eng{Expenses must not exceed the budget.}
\item \cle{Accountability} : obligation de rendre compte, redevabilité. \eng{Accountability builds trust.}
\item \cle{Transparency} : transparence. \eng{We publish our accounts for the sake of transparency.}
\item \cle{Audit} : audit, vérification comptable. \eng{An independent audit confirmed the figures.}
\end{itemize}
\end{definition}

\begin{aretenir}[Accountability : un mot intraduisible en un seul mot]
\eng{Accountability} vient de \eng{account} (compte) : c'est littéralement « l'obligation de rendre des comptes ». En français on dit « redevabilité » ou « reddition des comptes ». Au Bac, c'est un mot très valorisé dans les sujets sur la gouvernance et les ONG. Prononciation : « a-kaon-te-BI-li-ti ».
\end{aretenir}

\courssub{Champ lexical 4 : Ressources et services communautaires}

\begin{definition}[Community resources and services]
\begin{itemize}
\item \cle{Community hall} : salle communautaire, salle des fêtes du quartier.
\item \cle{Infrastructure} : infrastructures (routes, ponts, écoles). Toujours au singulier en anglais : \eng{The infrastructure is poor.}
\item \cle{Amenities} : équipements collectifs, commodités (eau, latrines, terrain de sport).
\item \cle{Utilities} : services publics (eau, électricité). \eng{The village lacks basic utilities.}
\item \cle{Healthcare access} : accès aux soins de santé.
\item \cle{Literacy program} : programme d'alphabétisation.
\item \cle{Vocational training} : formation professionnelle. \eng{Vocational training helps young people find jobs.}
\item \cle{Sustainable development} : développement durable.
\end{itemize}
\end{definition}

\begin{attention}[Pièges de traduction]
\begin{itemize}
\item \eng{Sustainable} = \textbf{durable} (\eng{sustainable development} = développement durable). Ne dites pas « soutenable ».
\item \eng{Infrastructure} est \textbf{singulier} en anglais : \eng{The infrastructure \textbf{is} improving}, pas \eng{are}.
\item « Une ressource » = \eng{a resource}, mais on dit souvent \eng{resources} au pluriel : \eng{to mobilize resources}.
\end{itemize}
\end{attention}

\courssub{Carte mentale récapitulative}

\begin{popfigure}
\begin{tikzpicture}[mindmap, concept color=popblue!50, align=center, level 1/.style={sibling angle=120, level distance=3.2cm}, level 2/.style={sibling angle=50, level distance=2.6cm}]
\node[concept, align=center]{Fund-Raising\\Lexique}
child[concept color=popgreen!50] {node[concept, align=center]{Sources\\of funds}
child {node[concept, align=center]{Grants\\Subventions}}
child {node[concept, align=center]{Donations\\Dons}}
child {node[concept, align=center]{Sponsorships\\Parrainages}}
child {node[concept] {Crowdfunding}}}
child[concept color=poppink!60] {node[concept] {Actions}
child {node[concept, align=center]{Campaign\\Campagne}}
child {node[concept, align=center]{Event\\Événement}}
child {node[concept, align=center]{Outreach\\Sensibilisation}}}
child[concept color=poporange!60] {node[concept] {Gestion}
child {node[concept, align=center]{Accountability\\Reddition}}
child {node[concept] {Budget}}
child {node[concept] {Audit}}};
\end{tikzpicture}
\legende{Carte mentale du lexique de la collecte de fonds : trois branches à mémoriser.}
\end{popfigure}

\courssub{Les collocations essentielles}

En anglais, certains verbes « collent » à certains noms. On dit \eng{launch a campaign} mais jamais \eng{open a campaign}. Voici les collocations à connaître par cœur pour le Bac.

\begin{center}
\begin{tabularx}{\linewidth}{|l|X|X|}
\hline
\rowcolor{popblueL} Collocation anglaise & Traduction française & Exemple en contexte \\
\hline
to launch a campaign & lancer une campagne & \eng{They launched a campaign to build a well.} \\
\hline
to secure a grant / funding & obtenir une subvention / un financement & \eng{The NGO secured a grant from UNESCO.} \\
\hline
to raise awareness & sensibiliser (litt. « élever la conscience ») & \eng{Posters raise awareness about hygiene.} \\
\hline
to raise funds / money & collecter des fonds & \eng{We raised 5 million FCFA in one month.} \\
\hline
to mobilize resources & mobiliser les ressources & \eng{The chief mobilized resources for the school.} \\
\hline
to empower the community & autonomiser / renforcer la communauté & \eng{Education empowers the community.} \\
\hline
to yield results & produire des résultats, porter ses fruits & \eng{The campaign finally yielded results.} \\
\hline
\end{tabularx}
\end{center}

\begin{exemplebox}[Exemples traduits à retenir]
\begin{itemize}
\item \eng{To launch a campaign.} « Lancer une campagne. »
\item \eng{To secure funding.} « Obtenir un financement » (plus fort et plus formel que \eng{to get money}).
\item \eng{A raffle.} « Une tombola. »
\item \eng{An auction.} « Une vente aux enchères. »
\item \eng{Vocational training.} « La formation professionnelle. »
\item \eng{Sustainable.} « Durable » (et non « soutenable »).
\item \eng{The association organized a food drive for flood victims.} « L'association a organisé une collecte de denrées pour les victimes des inondations. »
\item \eng{Transparency and accountability attract more donors.} « La transparence et la reddition des comptes attirent davantage de donateurs. »
\end{itemize}
\end{exemplebox}

\begin{savaistu}[Les Development Associations du Cameroun anglophone]
Dans les régions du Northwest et du Southwest, les \textbf{Development Associations} (par exemple la \emph{Buea Development Association}) jouent le rôle principal de \eng{fund-raisers} pour l'eau potable, l'électricité, les routes et les salles communautaires. Leur originalité : elles organisent des \textbf{Annual Conventions} aux États-Unis ou en Europe, où les membres de la diaspora se réunissent, font des \eng{pledges} et envoient des \eng{contributions}. C'est un modèle remarquable de \eng{community-driven development} (développement impulsé par la communauté elle-même), souvent cité comme exemple d'\eng{empowerment} local.
\end{savaistu}

\courssub{Exercice résolu et entraînement}

\begin{exoresolu}[Complétez avec le mot du champ lexical qui convient]
Énoncé — Complétez chaque phrase avec : \eng{grant}, \eng{volunteers}, \eng{raffle}, \eng{outreach}, \eng{transparency}, \eng{vocational training}.
\begin{enumerate}
\item Fifty \ldots\ldots{} helped to clean the community hall without being paid.
\item The association won a \ldots\ldots{} of 10 million FCFA from an international foundation.
\item We sold tickets for a \ldots\ldots{} and the winner received a goat!
\item The \ldots\ldots{} team visited remote villages to explain the project.
\item Young people learn carpentry and sewing in our \ldots\ldots{} centre.
\item The treasurer publishes all accounts: this is what we call \ldots\ldots{}.
\end{enumerate}
\tcblower
Solution détaillée :
\begin{enumerate}
\item \textbf{volunteers} — ils travaillent sans être payés : ce sont des bénévoles.
\item \textbf{grant} — une somme accordée par une fondation, sans remboursement : une subvention.
\item \textbf{raffle} — on vend des tickets et on tire un gagnant au sort : une tombola.
\item \textbf{outreach} — l'équipe « va vers » les villages : c'est une action de terrain.
\item \textbf{vocational training} — la menuiserie et la couture relèvent de la formation professionnelle.
\item \textbf{transparency} — publier les comptes, c'est la transparence.
\end{enumerate}
\end{exoresolu}

\begin{exercice}[À vous de jouer]
\begin{enumerate}
\item Traduisez en anglais : « Les donateurs ont promis deux millions de FCFA, et la campagne a porté ses fruits. »
\item Trouvez le mot : un événement où l'on vend des objets au plus offrant (\eng{a\ldots}).
\item Vrai ou faux : \eng{The infrastructure of the village are very poor.} Corrigez si nécessaire.
\item Complétez avec la bonne collocation : \eng{The students \ldots\ldots{} awareness about malaria.} (\eng{raise} / \eng{rise} / \eng{make})
\end{enumerate}
\end{exercice}

\begin{corrige}[Exercice « À vous de jouer »]
\begin{enumerate}
\item \eng{Donors pledged two million FCFA, and the campaign yielded results.} (On utilise \eng{pledge} pour « promettre » un don et la collocation \eng{yield results}.)
\item \eng{An auction} (une vente aux enchères).
\item \textbf{Faux.} \eng{Infrastructure} est singulier : \eng{The infrastructure of the village \textbf{is} very poor.}
\item \eng{The students \textbf{raise} awareness about malaria.} Attention : \eng{raise} (transitif, « élever quelque chose ») et non \eng{rise} (intransitif, « s'élever »).
\end{enumerate}
\end{corrige}

\begin{aretenir}[L'essentiel du lexique en cinq points]
\begin{enumerate}
\item Les acteurs : \eng{organizers, donors, beneficiaries, volunteers, sponsors, stakeholders}.
\item Les actions : \eng{campaign, drive, auction, raffle, crowdfunding, outreach}.
\item L'argent : \eng{grant} (subvention), \eng{donation}, \eng{pledge} (promesse de don), \eng{budget}, \eng{audit}.
\item La bonne gestion : \eng{transparency} et \eng{accountability}.
\item Les collocations du Bac : \eng{launch a campaign, secure a grant, raise awareness, mobilize resources, empower the community, yield results}.
\end{enumerate}
\end{aretenir}

Ce lexique est maintenant votre boîte à outils : vous le réutiliserez dans la section 5 pour rédiger une lettre d'appel aux dons (\eng{fund-raising appeal letter}) et dans la section 6 pour la simulation de réunion communautaire.

\coursec{Grammaire en contexte : Discours rapporté \& Passif (Révision Bac)}

Dans la section précédente, nous avons enrichi notre vocabulaire sur la collecte de fonds et les ressources communautaires (\eng{fund-raising, donation, sponsor, grant, community hall}...). Reprenons maintenant le texte support et observons \cle{comment} ces faits sont exprimés grammaticalement. Deux phénomènes majeurs apparaissent dans tout compte-rendu d'activité communautaire — et ce sont deux classiques de l'épreuve du Baccalauréat : le \cle{passif} et le \cle{discours rapporté}.

\courssub{Observation : l'actif et le passif dans le texte}

Relisons attentivement ces deux phrases extraites de notre texte sur la campagne de la \eng{NDA} (\eng{Ndop Development Association}) :

\begin{exemplebox}[Deux phrases, deux points de vue]
\eng{The NDA \textbf{launched} a fund-raising campaign in March.}\\
« La NDA a lancé une campagne de collecte de fonds en mars. »

\eng{The community hall \textbf{was built} by volunteers in six months.}\\
« La salle communautaire a été construite par des bénévoles en six mois. »
\end{exemplebox}

Dans la première phrase, le sujet (\eng{The NDA}) \textbf{fait} l'action : c'est la \cle{voix active}. Dans la seconde, le sujet (\eng{The community hall}) \textbf{subit} l'action : c'est la \cle{voix passive}. Remarquez que dans la phrase passive, l'attention se porte sur le résultat (la salle), pas sur ceux qui ont agi.

\begin{propriete}[Le passif : forme et emploi]
\textbf{Forme : Sujet + BE (conjugué au temps voulu) + Participe passé (V-ed / 3\textsuperscript{e} colonne) + (by + agent)}

\begin{center}
\begin{tabular}{|l|l|l|}
\hline
\rowcolor{popblueL} Voix active & Voix passive & Traduction du passif \\
\hline
Volunteers \textbf{built} the hall. & The hall \textbf{was built} (by volunteers). & La salle a été construite. \\
\hline
The NDA \textbf{is raising} funds. & Funds \textbf{are being raised}. & Des fonds sont en train d'être collectés. \\
\hline
They \textbf{have raised} 2 million FCFA. & 2 million FCFA \textbf{has been raised}. & 2 millions de FCFA ont été collectés. \\
\hline
The council \textbf{will open} the centre. & The centre \textbf{will be opened}. & Le centre sera inauguré. \\
\hline
\end{tabular}
\end{center}

\textbf{Emplois du passif :}
\begin{itemize}
\item l'\textbf{agent est inconnu} : \eng{The money was stolen.} « L'argent a été volé. »
\item l'agent est \textbf{évident ou sans importance} : \eng{The hall was inaugurated last Saturday.} « La salle a été inaugurée samedi dernier. »
\item on veut un \textbf{style objectif et impersonnel} : rapports d'activités, comptes-rendus, articles de presse, procès-verbaux de réunion.
\end{itemize}

\textbf{Exception :} les verbes intransitifs (\eng{go, come, die, happen, arrive}) n'ont \textbf{jamais} de passif, car ils n'ont pas de complément d'objet direct. On ne peut donc pas dire \eng{The accident was happened} : la seule forme correcte est \eng{The accident happened.} « L'accident a eu lieu. »
\end{propriete}

\begin{attention}[Erreurs fréquentes des francophones avec le passif]
\begin{enumerate}
\item \textbf{Oublier BE} : \eng{The hall built last year.} (FAUX) $\rightarrow$ \eng{The hall \textbf{was} built last year.} En français, « la salle construite » peut se dire ; en anglais, \eng{BE} est \textbf{obligatoire}.
\item \textbf{Mauvais participe passé} : attention aux verbes irréguliers ! \eng{The money was raised} (régulier : \eng{raise-raised-raised}) mais \eng{The well was \textbf{dug} by the youth} (irrégulier : \eng{dig-dug-dug}). Révisez votre liste : \eng{build-built-built, make-made-made, give-gave-given, write-wrote-written, sell-sold-sold}.
\item \textbf{Traduire « on » par \eng{people} ou \eng{they}} : « On a construit une école » se rend mieux par le passif : \eng{A school was built.} Le passif anglais remplace très souvent le « on » français.
\end{enumerate}
\end{attention}

\courssub{Le discours rapporté (Reported Speech)}

Dans notre texte, Mama Agnes déclare : \eng{“This hall is our bank.”} Pour rédiger un résumé formel ou un compte-rendu de réunion, on ne cite pas mot à mot : on \textbf{rapporte} les paroles. \eng{She said (that) the hall was their bank.} « Elle a dit que la salle était leur banque. »

\begin{propriete}[La concordance des temps (backshift)]
Quand le verbe introducteur est au \textbf{passé} (\eng{said, told, announced, explained}), les temps de la citation reculent d'un cran vers le passé :

\begin{center}
\begin{tabular}{|l|l|l|}
\hline
\rowcolor{popgreenL} Discours direct & Discours rapporté & Exemple rapporté \\
\hline
Present simple & Past simple & \eng{“I \textbf{help} the group.”} $\rightarrow$ She said she \textbf{helped} the group. \\
\hline
Present continuous & Past continuous & \eng{“We \textbf{are collecting} money.”} $\rightarrow$ They said they \textbf{were collecting} money. \\
\hline
Past simple / Present perfect & Past perfect & \eng{“We \textbf{have raised} funds.”} $\rightarrow$ They said they \textbf{had raised} funds. \\
\hline
Will & Would & \eng{“We \textbf{will build} a well.”} $\rightarrow$ He said they \textbf{would build} a well. \\
\hline
Can & Could & \eng{“You \textbf{can} join us.”} $\rightarrow$ She said we \textbf{could} join them. \\
\hline
Must & Had to & \eng{“We \textbf{must} act now.”} $\rightarrow$ He said they \textbf{had to} act then. \\
\hline
\end{tabular}
\end{center}

\textbf{Changements de pronoms, possessifs et adverbes :}
\begin{itemize}
\item \eng{I} $\rightarrow$ \eng{he/she} ; \eng{we} $\rightarrow$ \eng{they} ; \eng{my} $\rightarrow$ \eng{his/her} ; \eng{our} $\rightarrow$ \eng{their}
\item \eng{this/these} $\rightarrow$ \eng{that/those} ; \eng{here} $\rightarrow$ \eng{there} ; \eng{now} $\rightarrow$ \eng{then} ; \eng{today} $\rightarrow$ \eng{that day} ; \eng{tomorrow} $\rightarrow$ \eng{the next day} ; \eng{yesterday} $\rightarrow$ \eng{the day before} ; \eng{ago} $\rightarrow$ \eng{before}
\end{itemize}

\textbf{Verbes introducteurs utiles au Bac :} \eng{say (that)}, \eng{tell + quelqu'un (that)}, \eng{announce (that)}, \eng{explain (that)}, \eng{promise (that)}, \eng{ask + objet + to + V} (demande), \eng{advise + objet + to + V} (conseil), \eng{order + objet + to + V} (ordre).
\end{propriete}

\begin{popfigure}
\begin{tikzpicture}
\node[draw=popblue, fill=popblueL, rounded corners, align=center, text width=5.2cm] (direct) at (0,0) {\textbf{Direct Speech}\\[2pt] \eng{“We \textbf{have raised} enough}\\ \eng{money,” \textbf{said} the treasurer.}};
\node[draw=popgreen, fill=popgreenL, rounded corners, align=center, text width=5.2cm, right=3.2cm of direct] (indirect) {\textbf{Reported Speech}\\[2pt] \eng{The treasurer \textbf{said} (that) they}\\ \eng{\textbf{had raised} enough money.}};
\draw[<->, very thick, dashed, poporange] (direct) -- node[above, align=center, font=\small] {Backshift :\\ Present Perfect $\rightarrow$ Past Perfect} (indirect);
\end{tikzpicture}
\legende{Transformation du discours direct en discours rapporté : recul des temps (backshift)}
\end{popfigure}

\begin{attention}[Quand NE PAS faire le backshift]
\begin{enumerate}
\item Le verbe introducteur est au \textbf{présent} : \eng{He says he \textbf{is} tired.} « Il dit qu'il est fatigué. » — pas de changement de temps.
\item La phrase exprime une \textbf{vérité générale ou scientifique} : \eng{The teacher said the earth \textbf{is} round.} « Le professeur a dit que la terre est ronde. »
\item Le fait est \textbf{encore vrai au moment où l'on parle} : \eng{She said the hall \textbf{is} still under construction.} « Elle a dit que la salle est encore en construction. »
\end{enumerate}
\end{attention}

\courssub{Passif et discours rapporté combinés : le style du compte-rendu}

Au Baccalauréat, on vous demande souvent de résumer un texte ou de rapporter des paroles dans un style formel. Les deux outils se combinent naturellement :

\begin{exemplebox}[Exemples traduits]
\eng{Active: “Local businesses \textbf{sponsored} the festival.”}\\
« Des entreprises locales ont parrainé le festival. »

\eng{Passive: “The festival \textbf{was sponsored} by local businesses.”}\\
« Le festival a été parrainé par des entreprises locales. »

\eng{Reported: “The president \textbf{announced} that the roof \textbf{had been completed}.”}\\
« Le président a annoncé que le toit avait été achevé. » (discours rapporté + passif au past perfect)
\end{exemplebox}

\begin{center}
\begin{tabularx}{\linewidth}{|X|X|X|}
\hline
\rowcolor{popblueL} Français & English & Remarque \\
\hline
On a collecté des fonds. & Funds were collected. & « On » français = passif anglais \\
\hline
La salle se construit vite. & The hall is being built quickly. & Passif au present continuous \\
\hline
Elle a dit que tout était prêt. & She said (that) everything was ready. & \eng{that} est facultatif \\
\hline
Il m'a demandé de l'aider. & He asked me to help him. & \eng{ask + objet + to + V} \\
\hline
\end{tabularx}
\end{center}

\courssub{Entraînement guidé}

\begin{exoresolu}[De l'actif au passif, du direct au rapporté]
\textbf{Énoncé.} (a) Mettez à la voix passive : \eng{The women of the group sold 500 bags of cassava.} (b) Rapportez : \eng{“We will finish the water project next month,” the engineer said.}
\tcblower
\textbf{Solution détaillée.}

(a) \textbf{Étape 1 :} identifier l'objet de l'action : \eng{500 bags of cassava}. \textbf{Étape 2 :} le placer en tête comme nouveau sujet. \textbf{Étape 3 :} conjuguer \eng{BE} au temps du verbe actif (\eng{sold} = past simple $\rightarrow$ \eng{were}, car \eng{bags} est pluriel). \textbf{Étape 4 :} ajouter le participe passé \eng{sold} (irrégulier : \eng{sell-sold-sold}).

Réponse : \eng{500 bags of cassava \textbf{were sold} by the women of the group.} « 500 sacs de manioc ont été vendus par les femmes du groupe. »

(b) \textbf{Étape 1 :} verbe introducteur \eng{said} au passé $\rightarrow$ backshift obligatoire. \textbf{Étape 2 :} \eng{will} $\rightarrow$ \eng{would} ; \eng{we} $\rightarrow$ \eng{they} ; \eng{next month} $\rightarrow$ \eng{the following month}.

Réponse : \eng{The engineer said (that) they \textbf{would finish} the water project \textbf{the following month}.} « L'ingénieur a dit qu'ils termineraient le projet d'adduction d'eau le mois suivant. »
\end{exoresolu}

\begin{exercice}[À vous de jouer]
\begin{enumerate}
\item Mettez à la voix passive : \eng{The community raised two million francs for the health centre.}
\item Rapportez : \eng{“This hall is our bank,” Mama Agnes said.}
\item Rapportez : \eng{“I can bring ten chairs tomorrow,” Peter promised.}
\item Corrigez la faute : \eng{The new well was build by the youth last year.}
\end{enumerate}
\end{exercice}

\begin{corrige}[Exercice : passif et discours rapporté]
\begin{enumerate}
\item \eng{Two million francs \textbf{were raised} by the community for the health centre.} (\eng{francs} pluriel $\rightarrow$ \eng{were} ; participe régulier \eng{raised}.)
\item \eng{Mama Agnes said (that) \textbf{that} hall \textbf{was their} bank.} (\eng{this} $\rightarrow$ \eng{that} ; \eng{is} $\rightarrow$ \eng{was} ; \eng{our} $\rightarrow$ \eng{their}.)
\item \eng{Peter promised (that) he \textbf{could bring} ten chairs \textbf{the next day}.} (\eng{can} $\rightarrow$ \eng{could} ; \eng{tomorrow} $\rightarrow$ \eng{the next day}.)
\item \eng{The new well was \textbf{built} by the youth last year.} (Participe passé irrégulier : \eng{build-built-built}, jamais \eng{build} après \eng{was}.)
\end{enumerate}
\end{corrige}

\begin{aretenir}[L'essentiel de la section]
\begin{itemize}
\item \textbf{Passif} = Sujet + \eng{BE} (conjugué) + participe passé (+ \eng{by} + agent). On l'utilise quand l'action ou le résultat compte plus que l'agent — style idéal pour rapports et comptes-rendus.
\item \textbf{Discours rapporté} : verbe introducteur au passé $\rightarrow$ backshift (present $\rightarrow$ past ; present perfect $\rightarrow$ past perfect ; \eng{will} $\rightarrow$ \eng{would} ; \eng{can} $\rightarrow$ \eng{could}) + changement des pronoms et des marqueurs de temps et de lieu.
\item \textbf{Pas de backshift} si le verbe introducteur est au présent ou si l'énoncé est une vérité générale.
\end{itemize}
\end{aretenir}

\coursec{Production Guidée : Écriture d'un Appel aux Dons (Fund-raising Appeal Letter)}

Dans la section précédente, nous avons révisé deux outils grammaticaux essentiels du programme de Terminale : le \cle{discours rapporté} (\eng{reported speech}) et le \cle{passif}. Ces deux structures ne sont pas de simples exercices de grammaire : ce sont les marqueurs du \textbf{ton formel et crédible}. Or, c'est précisément ce ton qu'exige le genre de texte que nous allons maintenant produire : la \textbf{lettre d'appel aux dons} (\eng{fund-raising appeal letter}), un texte persuasif par excellence, fréquemment demandé à l'épreuve d'expression écrite (\emph{composition}) du baccalauréat.

\courssub{Observer d'abord : un modèle authentique commenté}

Avant d'écrire, observons. Lis la lettre suivante deux fois : une première fois en \emph{skimming} (idée générale), une seconde fois en \emph{scanning} (repérage des cinq parties signalées en gras).

\begin{textebox}[Model Appeal Letter — Mbankomo Development Union (MDU)]
\textbf{Subject: Urgent Appeal: Clean Water for Mbankomo Village}

Dear Partner,

\textbf{Hook:} At present, 3,000 villagers in Mbankomo walk 5 km every day to fetch unsafe water from a polluted stream. Water-borne diseases are rampant, and many children miss school because of repeated stomach infections.

\textbf{Project:} The Mbankomo Development Union (MDU) has launched a project to drill a solar-powered borehole in the village square. The geological survey has already been completed, and the land has been secured by the traditional council.

\textbf{Impact:} This borehole will serve the primary school, the health centre and 400 households. According to the local nurse, diarrhoea cases could be reduced by half within one year. Girls, who spend hours fetching water, will be able to attend school regularly.

\textbf{Call to Action:} We need 8,500,000 FCFA to complete the project. A donation of 50,000 FCFA buys ten metres of piping; 10,000 FCFA pays for one bag of cement. You can donate via Mobile Money: 6XX XXX XXX (MDU Account) or at any branch of our partner bank.

\textbf{Trust:} A full financial report will be published on our website, and receipts will be issued for all donations. Together, we can give Mbankomo the clean water it deserves.

Yours sincerely,

The President, MDU
\end{textebox}

\textbf{Glossaire :}

\begin{center}
\begin{tabularx}{\linewidth}{|l|X|l|}\hline
\rowcolor{popblueL} \textbf{English} & \textbf{Français} & \textbf{Prononciation} \\ \hline
to fetch water & aller chercher de l'eau & « fètch » \\ \hline
unsafe & insalubre, non potable & « an-sèïfe » \\ \hline
water-borne diseases & maladies transmises par l'eau & « ouo-teu-born » \\ \hline
rampant & qui sévit, endémique & « ram-peunt » \\ \hline
borehole & forage, puits foré & « bo-rol » \\ \hline
survey & étude (ici : étude géologique) & « seu-vèï » \\ \hline
to secure (land) & obtenir, sécuriser (un terrain) & « si-kiou-eu » \\ \hline
household & ménage, foyer & « raousse-old » \\ \hline
piping & tuyauterie, canalisations & « paï-ping » \\ \hline
receipt & reçu & « ri-site » (p muet) \\ \hline
to issue & délivrer, émettre & « i-chou » \\ \hline
\end{tabularx}
\end{center}

\courssub{La structure type d'une lettre d'appel aux dons}

\begin{definition}[Fund-raising appeal letter]
Une \cle{lettre d'appel aux dons} est une lettre formelle et persuasive envoyée par une association, une ONG ou un comité de développement à des partenaires (entreprises, diaspora, particuliers) pour collecter des fonds (\eng{to raise funds}) destinés à l'acquisition d'une ressource ou d'un service communautaire : forage, bibliothèque, équipement informatique, bloc sanitaire, etc.
\end{definition}

L'observation du modèle montre que la lettre suit une architecture en \textbf{cinq mouvements}, précédée d'un en-tête et d'un objet. Cette progression va du problème vers la solution, puis vers le lecteur lui-même.

\begin{popfigure}
\begin{tikzpicture}[node distance=0.8cm]
\node[draw, rectangle, rounded corners, fill=poppinkL, align=center, text width=6.5cm] (hook) {\textbf{1. HOOK}\\ Problème urgent (accroche)};
\node[draw, rectangle, rounded corners, fill=poporangeL, align=center, text width=6.5cm, below=of hook] (proj) {\textbf{2. PROJECT}\\ Solution concrète};
\node[draw, rectangle, rounded corners, fill=popgoldL, align=center, text width=6.5cm, below=of proj] (impact) {\textbf{3. IMPACT}\\ Bénéficiaires, chiffres};
\node[draw, rectangle, rounded corners, fill=popgreenL, align=center, text width=6.5cm, below=of impact] (cta) {\textbf{4. CALL TO ACTION}\\ Montant, Mobile Money, banque};
\node[draw, rectangle, rounded corners, fill=popblueL, align=center, text width=6.5cm, below=of cta] (thx) {\textbf{5. THANKS \& TRUST}\\ Transparence, contact};
\draw[->, thick, popdark] (hook) -- (proj);
\draw[->, thick, popdark] (proj) -- (impact);
\draw[->, thick, popdark] (impact) -- (cta);
\draw[->, thick, popdark] (cta) -- (thx);
\end{tikzpicture}
\legende{La pyramide persuasive : du problème à la confiance}
\end{popfigure}

\begin{propriete}[Les cinq mouvements de la lettre d'appel]
\begin{enumerate}
\item \textbf{En-tête et objet} : nom de l'association (avec logo si possible), adresse, date, puis \eng{Subject:} suivi d'un titre court et percutant.
\item \textbf{Hook (accroche)} : le problème, présenté avec un chiffre ou une image forte. Il doit créer l'\emph{urgence} sans exagérer.
\item \textbf{Project (le projet)} : la solution concrète, présentée comme déjà en marche (\eng{has been launched}, \eng{the survey has been completed}) — le passif y est roi.
\item \textbf{Impact} : qui bénéficie du projet, combien de personnes, quels changements mesurables.
\item \textbf{Call to Action + Trust} : un appel à l'action précis (combien donner, comment donner) suivi de garanties de transparence et de remerciements.
\end{enumerate}
\end{propriete}

\begin{attention}[Le Call to Action : le cœur de la lettre]
L'erreur la plus grave est d'émouvoir le lecteur... puis de ne jamais lui dire \textbf{quoi faire}. Un appel à l'action efficace répond à trois questions : \emph{Combien ?} (\eng{We need 8,500,000 FCFA}), \emph{Que paie mon don ?} (\eng{50,000 FCFA buys ten metres of piping}), \emph{Comment donner ?} (\eng{Donate via Mobile Money: 6XX XXX XXX}). Sans ces trois éléments, la lettre échoue, même si elle est bien écrite.
\end{attention}

\courssub{La boîte à outils persuasive : connecteurs et formules}

Une lettre d'appel n'est pas une lettre amicale. Elle utilise un registre soutenu, des connecteurs logiques et des formules de persuasion. Voici la trousse de l'élève de Terminale :

\begin{center}
\begin{tabularx}{\linewidth}{|l|X|X|}\hline
\rowcolor{popblueL} \textbf{Fonction} & \textbf{English} & \textbf{Français} \\ \hline
Urgence & \eng{There is an urgent need for...} & Il y a un besoin urgent de... \\ \hline
Conséquence & \eng{Consequently, ... / Therefore, ...} & Par conséquent, ... / C'est pourquoi... \\ \hline
Solution & \eng{To address this issue, we have launched...} & Pour répondre à ce problème, nous avons lancé... \\ \hline
Promesse au donateur & \eng{Your support will enable us to...} & Votre soutien nous permettra de... \\ \hline
Solidarité & \eng{Together, we can...} & Ensemble, nous pouvons... \\ \hline
Transparence & \eng{Accountability is guaranteed. / Receipts will be issued.} & La transparence est garantie. / Des reçus seront délivrés. \\ \hline
Remerciement & \eng{We are grateful for your generosity.} & Nous sommes reconnaissants de votre générosité. \\ \hline
\end{tabularx}
\end{center}

\begin{exemplebox}[La grammaire de la section 4 au service de la persuasion]
\textbf{Passif} (le projet avance, les choses sont faites sérieusement) :

\eng{The geological survey has been completed.} « L'étude géologique a été menée à bien. »

\eng{Receipts will be issued for all donations.} « Des reçus seront délivrés pour tous les dons. »

\textbf{Discours rapporté} (des témoignages crédibilisent le projet) :

\eng{The nurse said that diarrhoea cases had dropped since the first well was built.} « L'infirmier a déclaré que les cas de diarrhée avaient diminué depuis la construction du premier puits. »

\eng{The headmaster explained that the school needed twenty computers.} « Le proviseur a expliqué que le lycée avait besoin de vingt ordinateurs. »
\end{exemplebox}

\begin{attention}[Erreurs fréquentes des francophones]
\begin{itemize}
\item \textbf{Mélange de registres} : ne passe pas du familier (\eng{Please help us, we are suffering!}) au formel dans la même lettre. Choisis le ton formel et respectueux du début à la fin.
\item \textbf{Promettre l'impossible} : évite \eng{We will solve all the problems of the village.} Préfère des engagements mesurables : \eng{The borehole will provide clean water to 400 households.}
\item \textbf{Accord du collectif} : \eng{community} est un nom collectif singulier. Écris \eng{The community \textbf{needs} water} et non \eng{The community need water}.
\item \textbf{Faux amis} : \eng{to assist} signifie « aider » (et non « assister à », qui se dit \eng{to attend}) ; \eng{a donation} est « un don » ; \eng{to achieve} est « réaliser, accomplir » (et non « achever » au sens de finir, qui se dit \eng{to complete}).
\item \textbf{Traduction mot à mot} : « faire un don » se dit \eng{to \textbf{make} a donation} (et non \eng{to do a donation}) ; « collecter des fonds » se dit \eng{to \textbf{raise} funds} (et non \eng{to rise funds} — \eng{to rise} est intransitif : \eng{The sun rises}).
\end{itemize}
\end{attention}

\courssub{Rédaction guidée : choisir un projet local}

\begin{methode}[Écrire sa lettre d'appel en cinq étapes]
\begin{enumerate}
\item \textbf{Choisir un projet local réaliste} : le forage d'eau à Bafia, une bibliothèque scolaire à Foumban, l'équipement informatique d'un lycée de Yaoundé, un bloc sanitaire au marché de Bafoussam.
\item \textbf{Remplir une fiche de préparation} : problème (1 chiffre fort), solution (2 actions déjà réalisées, au passif), impact (bénéficiaires + chiffres), budget total, prix symbolique d'un don, moyen de paiement.
\item \textbf{Rédiger paragraphe par paragraphe} en suivant la pyramide : Hook → Project → Impact → Call to Action → Trust. Impose-toi au moins deux phrases au passif et une phrase au discours rapporté.
\item \textbf{Relire en binôme} avec la grille d'évaluation ci-dessous (peer-review).
\item \textbf{Finaliser la mise en page} : en-tête avec nom et logo de l'association, références (adresse, téléphone), objet, formule d'appel (\eng{Dear Partner,}), formule de politesse (\eng{Yours sincerely,}), signature et fonction (\eng{The President, MDU}).
\end{enumerate}
\end{methode}

\begin{methode}[Grille d'auto-évaluation (Peer Review)]
Coche chaque critère après relecture de la lettre de ton binôme :
\begin{itemize}
\item[$\square$] Structure en cinq paragraphes respectée (Hook, Project, Impact, Call to Action, Trust) ?
\item[$\square$] Vocabulaire du thème utilisé (\eng{borehole, water-borne, rampant, secured, piping, receipt}) ?
\item[$\square$] Au moins deux phrases au passif ? (ex. \eng{The survey \textbf{was conducted}...})
\item[$\square$] Au moins une phrase au discours rapporté ? (ex. \eng{The nurse \textbf{said} that cases \textbf{had dropped}...})
\item[$\square$] Ton formel et persuasif, sans familiarité ?
\item[$\square$] Coordonnées et montants clairs dans le Call to Action ?
\end{itemize}
\end{methode}

\courssub{Entraînement : de la fiche au texte}

\begin{exoresolu}[Transformer une fiche de projet en paragraphes]
Énoncé — À partir de la fiche suivante, rédige les paragraphes « Hook » et « Call to Action » d'une lettre d'appel pour la bibliothèque de Foumban : 2,000 élèves ; aucune bibliothèque ; livres partagés à cinq ; budget 4,000,000 FCFA ; 5,000 FCFA = un manuel scolaire ; paiement Mobile Money 6XX XXX XXX (Foumban Students' Association).
\tcblower
Solution détaillée —

\textbf{Hook :} \eng{In Foumban, more than 2,000 students prepare for their examinations without a single school library. Textbooks are so scarce that one book is often shared by five learners. Consequently, many brilliant students fail simply because they have nowhere to study.}

« À Foumban, plus de 2 000 élèves préparent leurs examens sans la moindre bibliothèque scolaire. Les manuels sont si rares qu'un livre est souvent partagé par cinq apprenants. Par conséquent, beaucoup d'élèves brillants échouent simplement parce qu'ils n'ont nulle part où étudier. »

\textbf{Call to Action :} \eng{We need 4,000,000 FCFA to equip the reading room. A donation of only 5,000 FCFA buys one textbook for our students. Your support will enable hundreds of learners to prepare for their exams in good conditions. Please donate via Mobile Money: 6XX XXX XXX (Foumban Students' Association). Together, we can open the doors of knowledge to Foumban's youth.}

« Nous avons besoin de 4 000 000 FCFA pour équiper la salle de lecture. Un don de seulement 5 000 FCFA achète un manuel scolaire pour nos élèves. Votre soutien permettra à des centaines d'apprenants de préparer leurs examens dans de bonnes conditions. Merci de donner par Mobile Money : 6XX XXX XXX (Foumban Students' Association). Ensemble, nous pouvons ouvrir les portes du savoir à la jeunesse de Foumban. »

Points de méthode : le Hook commence par un chiffre (\eng{2,000 students}) et se termine par le connecteur \eng{Consequently} ; le Call to Action répond aux trois questions (combien, que paie le don, comment donner) et se clôt par la formule de solidarité \eng{Together, we can...}.
\end{exoresolu}

\begin{exercice}[À ton tour]
Choisis l'un des projets suivants : (a) le forage d'eau à Bafia ; (b) l'équipement informatique de ton lycée. Rédige une lettre d'appel complète de 150 à 180 mots en respectant la pyramide en cinq mouvements. Contraintes obligatoires : deux phrases au passif, une phrase au discours rapporté, trois connecteurs de la boîte à outils, un Call to Action chiffré. Échange ensuite ta lettre avec ton binôme et évalue-la avec la grille de peer-review.
\end{exercice}

\begin{corrige}[Exercice — éléments attendus]
La lettre doit comporter : un en-tête (nom de l'association, date), un objet (\eng{Subject: Urgent Appeal: ...}), la formule \eng{Dear Partner,}, cinq paragraphes identifiables, au moins deux passifs corrects (\eng{has been launched}, \eng{was donated}, \eng{will be published}), une phrase de discours rapporté avec recul des temps (\eng{The principal said that the school needed...}), un montant total, un prix symbolique par don, un moyen de paiement, et la clôture \eng{Yours sincerely,} + signature et fonction. Toute familiarité, promesse irréaliste ou Call to Action flou doit être sanctionnée dans la grille.
\end{corrige}

\begin{aretenir}[L'essentiel de la section]
Une lettre d'appel aux dons réussie repose sur trois piliers : une \textbf{structure} en cinq mouvements (Hook → Project → Impact → Call to Action → Trust), une \textbf{grammaire de la crédibilité} (passif pour montrer que le projet avance, discours rapporté pour apporter des témoignages), et un \textbf{appel à l'action précis} (combien, que paie le don, comment donner). Le ton reste formel et persuasif du début à la fin, et la transparence (\eng{receipts will be issued}) installe la confiance qui transforme le lecteur en donateur.
\end{aretenir}

\coursec{Production Orale : Simulation de Réunion Communautaire (Role-play)}

Dans la section précédente, vous avez appris à \emph{rédiger} une lettre d'appel aux dons (\eng{fund-raising appeal letter}) : un texte écrit, planifié, relu, corrigé. Mais une collecte de fonds ne se prépare pas seulement sur le papier. Avant d'envoyer la moindre lettre, il faut que la communauté se \emph{réunisse}, \emph{décide} et \emph{répartisse les tâches}. C'est exactement ce que nous allons pratiquer maintenant : prendre la parole en réunion, en anglais, avec aisance et politesse. L'écrit prépare la décision ; l'oral la fait naître.

\courssub{Mise en situation : la réunion de l'association de développement}

\begin{experience}[La grande réunion de fin d'année]
Vous êtes membres de la \eng{Buea Town Development Association}. Ordre du jour (\eng{agenda}) : préparer la grande collecte de fonds de fin d'année (\eng{end-of-year fund-raising drive}) destinée à financer l'achat de bancs et d'ordinateurs pour le lycée du quartier. La classe est divisée en groupes de sept élèves. Chaque élève tire un rôle au sort et dispose de cinq minutes pour préparer au moins trois interventions. La réunion dure dix minutes et doit aboutir à trois décisions concrètes : le type d'événement, la date, et la répartition des responsabilités.
\end{experience}

\textbf{Distribution des rôles (\eng{role cards}) :}

\begin{center}
\begin{tabularx}{\linewidth}{|X|X|X|}
\hline
\rowcolor{popblueL} Rôle & English & Mission dans la réunion \\
\hline
Président(e) & Chairperson & Ouvre et dirige la réunion, donne la parole, conclut. \\
\hline
Trésorier(ère) & Treasurer & Présente le budget, rappelle le coût des projets. \\
\hline
Secrétaire & Secretary & Prend les notes (\eng{minutes}) et résume les décisions. \\
\hline
Représentant(e) des femmes & Women's representative & Propose les aspects sociaux et logistiques. \\
\hline
Représentant(e) des jeunes & Youth representative & Propose l'événement (concert, match, \eng{cultural night}). \\
\hline
Chef traditionnel / Notable & Traditional ruler / Elder & Donne la bénédiction, tranche en cas de désaccord. \\
\hline
Donateur potentiel (invité) & Potential donor (guest) & Pose des questions, exige la transparence (\eng{accountability}). \\
\hline
\end{tabularx}
\end{center}

\begin{popfigure}
\begin{tikzpicture}
\node[draw, circle, fill=popblueL, align=center, minimum size=1.6cm] (chair) at (0,0) {Chair};
\node[draw, circle, fill=popgreenL, align=center, minimum size=1.6cm] (sec) at (-2.5,-2) {Secretary};
\node[draw, circle, fill=popgreenL, align=center, minimum size=1.6cm] (tres) at (2.5,-2) {Treasurer};
\node[draw, circle, fill=poporangeL, align=center, minimum size=1.6cm] (women) at (-3.5,-4.2) {Women};
\node[draw, circle, fill=poporangeL, align=center, minimum size=1.6cm] (youth) at (3.5,-4.2) {Youth};
\node[draw, circle, fill=poppinkL, align=center, minimum size=1.6cm] (donor) at (0,-6) {Donor};
\draw[thick, popline] (chair) -- (sec);
\draw[thick, popline] (chair) -- (tres);
\draw[thick, popline] (chair) -- (women);
\draw[thick, popline] (chair) -- (youth);
\draw[thick, dashed, poporange] (chair) -- (donor);
\end{tikzpicture}
\legende{Organigramme de la réunion : le président (\eng{chair}) est au centre ; le donateur, invité extérieur, est relié en pointillés.}
\end{popfigure}

\courssub{Observer d'abord : un extrait de réunion authentique}

Avant d'énoncer les règles, lisons comment une vraie réunion communautaire se déroule en anglais. Observez les formules en gras : elles reviennent dans presque toutes les réunions anglophones.

\begin{textebox}[Extract from the minutes of the Buea Town Development Association]
\textbf{Chairperson:} Good morning, everyone. \textbf{I declare this meeting open}. The \textbf{agenda} today has three items: the type of fund-raising event, the date, and the sharing of tasks. \textbf{Shall we start}?

\textbf{Youth representative:} Madam Chair, \textbf{I propose we organize} a cultural night with a football match in the afternoon. Young people will pay a small entrance fee.

\textbf{Treasurer:} \textbf{I see your point, but have you considered} the cost? We shall need a public address system and security. \textbf{I'm not sure that's feasible} with our current funds.

\textbf{Women's representative:} \textbf{Why don't we approach} local businesses for sponsorship? The brewery and the two microfinance banks supported the harvest festival last year.

\textbf{Chairperson:} \textbf{That's a brilliant idea}. John, \textbf{can you take charge of} logistics? Mary, \textbf{you will liaise with} the sponsors and prepare the appeal letters.

\textbf{Potential donor:} Excuse me, Madam Chair. Before I commit any money, I would like to know how the funds will be managed. \textbf{Could you clarify} the accountability procedures?

\textbf{Treasurer:} Certainly. Every franc will be recorded, and a financial report \textbf{will be read} at the next meeting.

\textbf{Chairperson:} \textbf{To sum up}, we have agreed on a cultural night on the second Saturday of December. John handles logistics, Mary contacts the sponsors, and the treasurer will present a budget. \textbf{Let's meet next week}, same time. \textbf{The meeting is adjourned}. Thank you all.
\end{textebox}

\textbf{Glossaire :}

\begin{center}
\begin{tabularx}{\linewidth}{|X|X|X|}
\hline
\rowcolor{popblueL} Mot anglais & Nature & Traduction \\
\hline
agenda & nom & ordre du jour \\
\hline
entrance fee & nom & droit d'entrée \\
\hline
sponsorship & nom & parrainage, mécénat \\
\hline
to liaise with & verbe & assurer la liaison avec \\
\hline
to commit (money) & verbe & engager (de l'argent) \\
\hline
accountability & nom & obligation de rendre compte \\
\hline
to adjourn & verbe & lever (la séance) \\
\hline
\end{tabularx}
\end{center}

\textbf{Questions de compréhension :}
\begin{enumerate}
\item What are the three items on the agenda?
\item Why does the treasurer doubt the youth representative's proposal at first?
\item What solution does the women's representative suggest?
\item What does the donor demand before giving money?
\item What two decisions are taken at the end of the meeting?
\end{enumerate}

\courssub{Les formules fonctionnelles (\eng{functional language})}

\begin{definition}[Functional language]
La \emph{langue fonctionnelle} est l'ensemble des expressions toutes faites qui permettent d'accomplir un acte de parole précis : ouvrir une réunion, proposer, accepter, refuser poliment, conclure. Ces formules sont \cle{fixes} : on ne les traduit pas mot à mot, on les mémorise comme des blocs.
\end{definition}

\begin{exemplebox}[Formules utiles par fonction]
\begin{itemize}
\item \textbf{Opening :} \eng{Good morning everyone. Let's get started. The agenda today is...} « Bonjour à tous. Commençons. L'ordre du jour est... »
\item \textbf{Chairing :} \eng{I declare the meeting open.} « Je déclare la séance ouverte. » \eng{John, you have the floor.} « John, vous avez la parole. »
\item \textbf{Proposing :} \eng{I suggest organizing a cultural night.} « Je suggère d'organiser une nuit culturelle. » \eng{Why don't we approach local businesses?} « Pourquoi ne pas solliciter les entreprises locales ? »
\item \textbf{Agreeing :} \eng{Absolutely. That's a great idea. I'm in favor of that.} « Absolument. C'est une excellente idée. J'y suis favorable. »
\item \textbf{Disagreeing politely :} \eng{I'm not sure that's feasible.} « Je ne suis pas sûr que ce soit réalisable. » \eng{I see your point, but have you considered the cost?} « Je comprends votre point de vue, mais avez-vous pensé au coût ? »
\item \textbf{Assigning tasks :} \eng{John, can you take charge of logistics?} « John, peux-tu prendre en charge la logistique ? » \eng{Mary will liaise with the sponsors.} « Mary assurera la liaison avec les sponsors. »
\item \textbf{Closing :} \eng{To sum up, we agreed on... Next meeting is on Friday.} « En résumé, nous avons décidé de... La prochaine réunion aura lieu vendredi. »
\end{itemize}
\end{exemplebox}

\begin{propriete}[La grammaire cachée derrière les formules de proposition]
Trois structures à maîtriser absolument :
\begin{itemize}
\item \eng{suggest} + \textbf{verbe en -ing} : \eng{I suggest organizing a concert.} « Je suggère d'organiser un concert. » (jamais \eng{I suggest to organize}).
\item \eng{Why don't we} + \textbf{base verbale} : \eng{Why don't we invite the mayor?} « Pourquoi n'inviterions-nous pas le maire ? »
\item \eng{I propose we} + \textbf{subjonctif} (base verbale sans -s à la 3\textsuperscript{e} personne) : \eng{I propose we organize a raffle.} / \eng{I suggest that the treasurer \textbf{present} a budget.} « Je suggère que le trésorier présente un budget. »
\end{itemize}
\end{propriete}

\begin{attention}[Erreurs fréquentes des francophones]
\begin{itemize}
\item Ne dites pas \emph{I am agree} : le verbe \eng{agree} n'utilise jamais \eng{be}. Dites \eng{I agree} « je suis d'accord ».
\item \emph{Assist} est un faux ami : \eng{assist} signifie « aider ». Pour « assister à une réunion », dites \eng{attend a meeting}.
\item Ne traduisez pas « prendre la parole » par \emph{take the word} : dites \eng{take the floor} ou \eng{May I say something?}
\item \eng{Actually} signifie « en fait », pas « actuellement » : pour « actuellement », dites \eng{currently}.
\end{itemize}
\end{attention}

\courssub{Focus prononciation : l'accent tonique et l'intonation}

En anglais, le sens est porté par les \cle{mots forts} (noms, verbes, adjectifs) ; les petits mots (articles, prépositions, auxiliaires) sont avalés. C'est le contraire du français, où toutes les syllabes ont à peu près le même poids.

\begin{aretenir}[Mots-clés de la leçon : accent tonique]
\begin{itemize}
\item \textbf{FUND-raising} : « FOND-reï-zing » — accent sur \emph{FOND}.
\item \textbf{comMUnity} : « keu-MIU-ni-ti » — accent sur \emph{MIU}.
\item \textbf{SPONsorship} : « SPON-seur-chip » — accent sur \emph{SPON}.
\item \textbf{acCOUNTaBILity} : « a-KAUN-ta-BI-li-ti » — accent principal sur \emph{BI}.
\item \textbf{deVElopment} : « di-VÉ-leup-ment » — accent sur \emph{VÉ}.
\end{itemize}
\end{aretenir}

\begin{propriete}[L'intonation : montante ou descendante ?]
\begin{itemize}
\item \textbf{Intonation montante} (la voix monte à la fin) : suggestions, questions ouvertes de politesse. \eng{Shall we start?} « On commence ? » — la voix monte sur \emph{start}.
\item \textbf{Intonation descendante} (la voix descend à la fin) : décisions, affirmations, ordres. \eng{We start at ten.} « Nous commençons à dix heures. » — la voix descend sur \emph{ten}.
\end{itemize}
Règle d'or en réunion : on \emph{suggère} avec une voix qui monte (souplesse), on \emph{décide} avec une voix qui descend (autorité).
\end{propriete}

\courssub{Entraînement guidé et auto-évaluation}

\begin{exoresolu}[Compléter une réunion]
Énoncé : complétez les interventions suivantes avec une formule fonctionnelle adaptée, puis jouez la scène à voix haute.
\begin{enumerate}
\item (ouverture) « Good morning everyone. \ldots\ The agenda today is the harvest festival. »
\item (proposition) « \ldots\ a raffle to raise money for the school library. »
\item (désaccord poli) « \ldots, but a raffle requires a licence from the council. »
\item (attribution de tâche) « Peter, \ldots\ the printing of the tickets? »
\item (clôture) « \ldots, we agreed on a raffle on December 15th. \ldots\ next Monday. »
\end{enumerate}
\tcblower
Solution détaillée :
\begin{enumerate}
\item \eng{I declare this meeting open.} (formule d'ouverture solennelle du président).
\item \eng{I suggest organizing} ou \eng{I propose we organize} — attention : \eng{suggest} + -ing, \eng{propose we} + base verbale.
\item \eng{I see your point} — formule de désaccord poli qui reconnaît d'abord l'argument de l'autre.
\item \eng{can you take charge of} — attribution directe mais polie d'une tâche.
\item \eng{To sum up} ; \eng{Let's meet} — résumé des décisions puis annonce de la prochaine réunion.
\end{enumerate}
À l'oral : montez la voix sur la suggestion (2), descendez-la sur la décision finale (5).
\end{exoresolu}

\begin{methode}[Réussir le jeu de rôle en cinq étapes]
\begin{enumerate}
\item \textbf{Préparer} : relire sa fiche de rôle et choisir trois formules fonctionnelles à placer.
\item \textbf{Écouter} : attendre que le président donne la parole (\eng{you have the floor}) avant d'intervenir.
\item \textbf{Interagir} : réagir à l'intervenant précédent (\eng{I see your point, but...}) au lieu de réciter sa phrase.
\item \textbf{Soigner la voix} : accentuer les mots porteurs de sens, monter l'intonation pour suggérer, la descendre pour décider.
\item \textbf{S'auto-évaluer} : si possible, enregistrer la réunion (téléphone) et réécouter en cochant la grille : fluidité, vocabulaire précis du thème (\eng{fund-raising, sponsorship, accountability}), interaction réelle.
\end{enumerate}
\end{methode}

\begin{exercice}[À vous de jouer]
En groupes de sept, tenez la réunion complète de la \eng{Buea Town Development Association}. Contraintes : chaque membre doit utiliser au moins deux formules fonctionnelles différentes ; le secrétaire rédige ensuite un compte rendu de cinq lignes commençant par \eng{It was decided that...} ; le donateur ne s'engage que si la transparence financière est garantie. Enregistrez la réunion et réécoutez-la avec la grille d'auto-correction.
\end{exercice}

\begin{savaistu}[Robert's Rules et l'arbre à palabres]
Dans la culture anglo-saxonne (Royaume-Uni, États-Unis, Nigeria, Ghana), les réunions formelles suivent une version simplifiée des \emph{Robert's Rules of Order} : une \eng{motion} (proposition) est faite, puis \eng{seconded} (appuyée par un second membre), puis discutée, puis mise au vote. Au Cameroun, on mélange souvent cette procédure formelle et la tradition orale de l'\emph{arbre à palabres}, où le consensus et la parole de l'aîné priment sur le vote. Savoir naviguer entre ces deux modèles — voter comme à Londres, consensualiser comme au village — est un véritable atout professionnel dans les ONG, les coopératives et les administrations camerounaises, où les réunions se tiennent de plus en plus en anglais.
\end{savaistu}

\begin{aretenir}[L'essentiel de la section]
Une réunion anglophone se conduit avec des \cle{formules fixes} : ouvrir (\eng{I declare the meeting open}), proposer (\eng{I suggest organizing}, \eng{Why don't we...}), nuancer (\eng{I see your point, but...}), déléguer (\eng{Can you take charge of...?}) et conclure (\eng{To sum up... Let's meet next week}). La prononciation repose sur l'\cle{accent tonique} des mots de sens et sur l'\cle{intonation} : montante pour suggérer, descendante pour décider. Ces compétences orales complètent la lettre d'appel aux dons rédigée dans la section précédente : l'écrit convainc à distance, l'oral décide en face à face.
\end{aretenir}

\newpage

\coursec{Activité d'intégration}

\courssub{Objectifs de l'activité}

À la fin de cette activité d'intégration, l'élève sera capable de :

\begin{itemize}
\item \textbf{Comprendre} un texte authentique sur une campagne de collecte de fonds communautaire (\eng{skimming} et \eng{scanning}) ;
\item \textbf{Mobiliser} le lexique spécialisé de la collecte de fonds et des ressources communautaires (\cle{fund-raising}, \cle{donor}, \cle{appeal}, \cle{target}, \cle{proceeds}, \cle{transparency}...) ;
\item \textbf{Réutiliser} la grammaire révisée (discours rapporté et voix passive) dans des productions personnelles ;
\item \textbf{Rédiger} une lettre formelle d'appel aux dons en cinq étapes ;
\item \textbf{Interagir oralement} dans une réunion communautaire simulée, avec des rôles distribués ;
\item \textbf{Synthétiser} un projet sur un support visuel (affiche) et le présenter devant la classe.
\end{itemize}

\courssub{Texte support : Modèle de campagne de collecte de fonds}

\begin{textebox}[Model Fund-raising Text — “A Library for Buea Community High School”]
“Education is the most powerful weapon which you can use to change the world,” Nelson Mandela once said. At Buea Community High School, this weapon is blunted by a lack of books. Our school library has only 50 outdated textbooks for 400 students. Many pupils have never held a storybook or a science encyclopedia.

The Parents-Teachers Association (PTA) has launched the “Books for Tomorrow” campaign. Our target is to raise 2,500,000 FCFA by the end of the academic year. The proceeds will buy 500 curriculum textbooks, 200 fiction titles, and fund the installation of solar-powered reading lamps. A local carpenter will build the shelves, and the books will be catalogued by volunteer students.

The PTA chairperson declared, “We cannot wait for government grants alone; the community must act now.” He added that the library would be named after the highest donor. Transparency is guaranteed: a financial report will be published monthly on the school notice board and the PTA WhatsApp group.

Every contribution counts. You can donate via Mobile Money (6XX XXX XXX), bank transfer to the PTA account, or in cash at the school bursar's office. Together, let us turn the page for our children's future.
\end{textebox}

\textbf{Traduction des passages clés :} « L'éducation est l'arme la plus puissante... » (N. Mandela). « Notre bibliothèque ne compte que 50 manuels obsolètes pour 400 élèves. » « L'APE a lancé la campagne “Livres pour Demain”. Objectif : 2 500 000 FCFA. » « Le président de l'APE a déclaré : “Nous ne pouvons pas attendre les subventions de l'État seul ; la communauté doit agir maintenant.” Il a ajouté que la bibliothèque porterait le nom du plus grand donateur. » « La transparence est garantie : un rapport financier sera publié chaque mois. »

\courssub{Situation d'intégration}

\begin{textebox}[Contexte de la mission]
Your school, Government High School Nkol-Eton (Centre Region, Cameroon), has launched its annual \emph{Community Development Week}. The Parents-Teachers Association (PTA) has announced a competition: \emph{“My Village / My Neighbourhood Has Talent”}. Each class group must design a complete fund-raising campaign for a real or realistic community need: repairing a classroom blackboard, buying a motorbike ambulance for the health centre, setting up a class library, or sinking a borehole for clean water.

The best project will receive a \emph{seed grant} of 100,000 FCFA from a local microfinance institution. Your group of four or five students has been selected to represent your class. You must submit a written file (project article + appeal letter + poster) and defend your campaign in a five-minute simulated launching meeting in front of a jury of teachers and classmates.
\end{textebox}

\textbf{Traduction du contexte :} Ton école lance sa Semaine du développement communautaire. L'association des parents d'élèves organise un concours : chaque groupe doit concevoir une campagne de collecte de fonds complète pour un besoin communautaire réel ou réaliste. Le meilleur projet recevra une aide de démarrage de 100 000 FCFA. Ton groupe doit remettre un dossier écrit et défendre la campagne lors d'une réunion de lancement simulée de cinq minutes.

\courssub{Consignes}

Travaillez en groupes de quatre ou cinq élèves. Répartissez les rôles dès le départ : \eng{project leader, secretary, treasurer, spokesperson, designer}.

\begin{enumerate}
\item \textbf{Choisir et comprendre le besoin.} Lisez rapidement (\eng{skimming}) le texte modèle ci-dessus, puis repérez par \eng{scanning} les informations clés : le besoin, le montant visé, les activités prévues, les bénéficiaires. Choisissez ensuite un besoin précis pour votre village ou quartier et complétez la fiche : \eng{What is the need? Who will benefit? How much is needed? By when?}

\item \textbf{Rédiger le texte de présentation du projet} (environ 150 mots, style article ou rapport). Structure attendue : titre accrocheur, présentation du problème, solution proposée, budget estimé, appel à l'action. Utilisez au moins \textbf{deux phrases à la voix passive} (ex. \eng{The borehole will be drilled by a local company.}) et \textbf{une phrase au discours rapporté} (ex. \eng{The quarter head said that clean water was a priority.}).

\item \textbf{Rédiger la lettre formelle d'appel aux dons} en suivant les cinq étapes vues en leçon : en-tête et adresses, salutation, présentation du besoin, demande précise (montant, moyens de paiement), formule de politesse et remerciement. Adressez-la à un donateur réaliste (commerçant, élite du village, ONG).

\item \textbf{Concevoir l'affiche (poster)} avec quatre rubriques obligatoires : \eng{Goal, Budget, Payment methods, Accountability}. Soyez clairs, visuels et honnêtes : la transparence convainc les donateurs.

\item \textbf{Simuler la réunion de lancement} (5 minutes). Chaque membre parle : le leader présente le projet, le trésorier détaille le budget, le porte-parole lance l'appel aux dons, un membre répond aux questions du public. Utilisez les formules fonctionnelles : \eng{“Let's get started.”}, \eng{“Could I have your attention, please?”}, \eng{“Every contribution counts.”}

\item \textbf{Présenter et évaluer.} Chaque groupe présente son projet devant la classe. Les autres groupes et le professeur notent avec la grille unique : Contenu, Langue, Communication, Créativité (sur 5 points chacun). Le groupe gagnant remporte le \eng{seed grant} symbolique.
\end{enumerate}

\courssub{Ressources à mobiliser : Lexique, Grammaire, Méthode}

\begin{aretenir}[Lexique utile : Fund-raising \& Community Resources]
\begin{itemize}
\item \cle{fund-raising} (n.) : collecte de fonds ; \cle{donor} (n.) : donateur ; \cle{appeal} (n.) : appel (aux dons) ; \cle{target / goal} (n.) : objectif (chiffré) ; \cle{proceeds} (n. plur.) : recettes, produit de la collecte ; \cle{contribution} (n.) : contribution, don ; \cle{accountability} (n.) : reddition de comptes, transparence ; \cle{borehole} (n.) : forage ; \cle{seed grant} (n.) : aide de démarrage, subvention d'amorçage ; \cle{pledge} (n./v.) : engagement de don / s'engager à donner ; \cle{outreach} (n.) : action de sensibilisation, démarche vers le public.
\item \textbf{Expressions clés :} \eng{to raise funds for}, \eng{to launch a campaign}, \eng{to donate to}, \eng{every contribution counts}, \eng{to keep records of donations}, \eng{to be accountable for}, \eng{to reach a target}.
\end{itemize}
\end{aretenir}

\begin{propriete}[Révision Bac : Voix Passive (Passive Voice)]
\textbf{Forme :} Sujet + \emph{be} (conjugué au temps requis) + \emph{participe passé} (+ \emph{by} + agent si nécessaire).
\begin{center}
\begin{tabular}{|l|l|l|}
\hline
\textbf{Temps} & \textbf{Actif} & \textbf{Passif} \\
\hline
Present Simple & They build schools. & Schools \textbf{are built}. \\
Past Simple & They built the well. & The well \textbf{was built}. \\
Present Perfect & They have raised 1M. & 1M \textbf{has been raised}. \\
Future (will) & We will drill a borehole. & A borehole \textbf{will be drilled}. \\
Modals (must/can) & We must finish it. & It \textbf{must be finished}. \\
\hline
\end{tabular}
\end{center}
\textbf{Emploi :} Mettre l'accent sur l'action/le résultat (le forage, l'argent) plutôt que sur l'acteur (l'entreprise, les donateurs). Très utilisé dans les rapports, articles, lettres formelles.
\textbf{Attention :} Accord du verbe \emph{be} avec le \textbf{sujet} (singulier/pluriel). \eng{The funds \textbf{were} collected} (pluriel).
\end{propriete}

\begin{propriete}[Révision Bac : Discours Rapporté (Reported Speech)]
\textbf{Principe :} Recul des temps (\emph{backshift}) si le verbe introducteur est au passé (\emph{said, told, declared, added}).
\begin{center}
\begin{tabular}{|l|l|l|}
\hline
\textbf{Direct} & \textbf{Rapporté} & \textbf{Exemple} \\
\hline
Present Simple & Past Simple & \eng{“We need money.”} $\rightarrow$ \eng{They said they \textbf{needed} money.} \\
Present Continuous & Past Continuous & \eng{“We are raising funds.”} $\rightarrow$ \eng{They said they \textbf{were raising} funds.} \\
Present Perfect & Past Perfect & \eng{“We have reached the target.”} $\rightarrow$ \eng{They said they \textbf{had reached} the target.} \\
\emph{will} & \emph{would} & \eng{“We will build it.”} $\rightarrow$ \eng{They said they \textbf{would} build it.} \\
\emph{can} & \emph{could} & \eng{“We can help.”} $\rightarrow$ \eng{They said they \textbf{could} help.} \\
\hline
\end{tabular}
\end{center}
\textbf{Autres changements :} \emph{now} $\rightarrow$ \emph{then}, \emph{today} $\rightarrow$ \emph{that day}, \emph{here} $\rightarrow$ \emph{there}, \emph{this/these} $\rightarrow$ \emph{that/those}. Le \emph{that} de liaison est optionnel.
\end{propriete}

\begin{methode}[La lettre formelle d'appel aux dons (Appeal Letter) en 5 étapes]
\begin{enumerate}
\item \textbf{En-tête (Heading) :} Adresse de l'expéditeur (en haut à droite), Date (sous l'adresse), Adresse du destinataire (à gauche, sous la date).
\item \textbf{Salutation :} \eng{Dear Sir/Madam,} (si nom inconnu) ou \eng{Dear Mr/Mrs [Nom],} (si connu). Pas de prénom seul.
\item \textbf{Corps de la lettre (Body) :}
      \begin{itemize}
      \item \emph{Paragraphe 1 :} Présenter l'organisme, le projet et le besoin urgent (\eng{We are writing to appeal for...}).
      \item \emph{Paragraphe 2 :} Donner des détails concrets (chiffres, bénéficiaires) et annoncer le montant visé (\eng{Our target is...}).
      \item \emph{Paragraphe 3 :} Formuler la demande précise et indiquer les modalités de paiement (\eng{Donations can be made via...}).
      \item \emph{Paragraphe 4 :} Rassurer sur la transparence (\eng{A financial report will be published...}).
      \end{itemize}
\item \textbf{Conclusion :} Remerciement anticipé (\eng{Thank you in advance for your generosity.} / \eng{We look forward to your support.}).
\item \textbf{Formule de politesse finale (Complimentary close) :} \eng{Yours faithfully,} (si \eng{Dear Sir/Madam}) ou \eng{Yours sincerely,} (si nom connu) + Signature + Nom tapé + Fonction.
\end{enumerate}
\end{methode}

\courssub{Production guidée : Modèle commenté (Corrigé type)}

\begin{exoresolu}[Dossier modèle : “Clean Water for Nkol-Eton”]
\textbf{Énoncé.} Produire le texte de présentation (article de projet, $\approx$ 150 mots) et la lettre d'appel aux dons pour le projet “Forage à Nkol-Eton”.

\tcblower

\textbf{1. Project Article (Modèle)}

\textit{\textbf{Clean Water for Nkol-Eton: Our Borehole Campaign}}

\textit{In Nkol-Eton, students walk two kilometres every morning to fetch water before classes. This problem affects over 600 pupils and reduces study time considerably. Our group has decided to launch a fund-raising campaign to sink a borehole near the school playground. The borehole \textbf{will be drilled} by a local company, and a hand pump \textbf{will be installed}. The total cost is estimated at 450,000 FCFA. The quarter head \textbf{said that clean water was} “the first priority of the village”. Donations can be made by mobile money or in cash to the school treasurer. A public record of all donations \textbf{will be displayed} on the notice board. Every contribution counts: together, we can give our school clean water before the end of the year.}

\textbf{Commentaire linguistique :}
\begin{itemize}
\item Structure respectée : Problème $\rightarrow$ Solution $\rightarrow$ Budget $\rightarrow$ Appel.
\item \textbf{3 voix passives} (au moins 2 requises) : \eng{will be drilled}, \eng{will be installed}, \eng{will be displayed}. Le sujet reçoit l'action.
\item \textbf{1 discours rapporté} correct : \eng{said that clean water was} (recul \emph{is} $\rightarrow$ \emph{was}). Les guillemets anglais “ ” sont conservés pour la citation imbriquée.
\item Lexique ciblé : \eng{fetch water, sink a borehole, hand pump, quarter head, notice board}.
\end{itemize}

\textbf{2. Appeal Letter (Modèle)}

\textit{Government High School Nkol-Eton\\
P.O. Box 12, Nkol-Eton\\
15th March 2025}

\textit{The Manager\\
Nkol-Eton Microfinance Institution\\
P.O. Box 45, Nkol-Eton}

\textit{Dear Sir/Madam,}

\textit{We are writing to appeal for your support for our borehole project, “Clean Water for Nkol-Eton”. Our 600 pupils currently have no access to clean water at school, which exposes them to water-borne diseases and wastes valuable study time.}

\textit{We need to raise 450,000 FCFA to drill and equip the borehole. So far, 180,000 FCFA \textbf{has already been collected} from the PTA and local businesses. Your donation, however small, would make a real difference. Contributions can be sent by Mobile Money to 6XX XX XX XX (account name: GHS Nkol-Eton Water Project) or handed to the school treasurer. A full financial report \textbf{will be published} at the end of the campaign.}

\textit{Thank you in advance for your generosity.}

\textit{Yours faithfully,}

\textit{(Signature)}

\textit{John Ngono\\
Project Leader, Clean Water Committee}

\textbf{Commentaire linguistique :}
\begin{itemize}
\item Format formel respecté (adresses, date, objet implicite, formules).
\item Ton professionnel : \eng{We are writing to appeal for...}, \eng{would make a real difference}.
\item \textbf{Voix passive} naturelle : \eng{has already been collected} (Present Perfect passif), \eng{will be published} (Future passif).
\item Précision : Montant total, montant acquis, numéro de téléphone, nom du compte.
\item Promesse de reddition de comptes (\eng{accountability}) explicite.
\end{itemize}
\end{exoresolu}

\courssub{Grille d'évaluation de l'activité (Sur 20)}

\begin{center}
\begin{tabularx}{\linewidth}{|X|c|c|c|c|}
\hline
\rowcolor{popblueL} \textbf{Critère} & \textbf{5 pts} & \textbf{4 pts} & \textbf{2-3 pts} & \textbf{0-1 pt} \\
\hline
\textbf{Contenu \& Cohérence} & Projet réaliste, complet, budget clair, besoin justifié. & Projet clair, 1-2 détails manquants. & Projet vague, budget flou. & Incomplet, hors sujet. \\
\hline
\textbf{Langue (Grammaire/Lexique)} & Passif \& Discours rapporté corrects ; vocabulaire précis, varié. & 1-2 erreurs mineures (temps, accord). & Erreurs fréquentes qui gênent la lecture. & Anglais très approximatif. \\
\hline
\textbf{Communication (Écrit/Oral)} & Lettre formelle parfaite ; oral fluide, rôles tenus, formules utilisées. & Bonnes productions, légère hésitation orale. & Difficulté à s'exprimer, lecture du texte. & Incompréhensible ou mutisme. \\
\hline
\textbf{Créativité \& Présentation} & Affiche visuelle, originale, claire ; présentation dynamique. & Affiche correcte, présentation standard. & Affiche minimaliste, lecture monotone. & Pas d'affiche, pas d'effort. \\
\hline
\end{tabularx}
\end{center}

\courssub{Bilan de la leçon}

\begin{aretenir}[Ce que j'ai appris / Dois savoir faire]
\begin{itemize}
\item Je sais \textbf{lire efficacement} un texte de collecte de fonds : \eng{skimming} pour l'idée générale, \eng{scanning} pour les chiffres, dates, noms.
\item Je maîtrise le \textbf{lexique} de la solidarité communautaire (\cle{fund-raising, donor, proceeds, accountability, seed grant}...) et je l'emploie en contexte.
\item Je réinvestis la \textbf{voix passive} (forme \emph{be} + V3, accord sujet) et le \textbf{discours rapporté} (recul des temps, pronoms, indicateurs spatio-temporels) dans mes écrits.
\item Je sais \textbf{rédiger une lettre formelle} d'appel aux dons (5 étapes : en-tête, salutation, corps structuré, conclusion, formule finale).
\item Je sais \textbf{présenter un projet} oralement en équipe, avec des rôles définis et des formules de communication (\eng{“Let's get started”, “Every contribution counts”}).
\item J'ai compris que la \textbf{transparence} (\eng{accountability}) est la clé de la confiance des donateurs et de la réussite d'une campagne.
\end{itemize}
\end{aretenir}

\begin{attention}[Erreurs fréquentes à éviter (Pièges Bac)]
\begin{itemize}
\item \textbf{Infinitif de but :} Ne dites pas \eng{*to raise money \textbf{for to} buy...} ; dites \eng{to raise money \textbf{to} buy...} ou \eng{to raise funds \textbf{for the purchase of}...}.
\item \textbf{Discours rapporté :} N'oubliez jamais le \textbf{recul du temps} : \eng{He said he \textbf{would} help} (et non \emph{will}) ; \eng{They said they \textbf{had} finished} (et non \emph{have}).
\item \textbf{Voix passive :} Le verbe \emph{be} \textbf{doit s'accorder} avec le sujet : \eng{The funds \textbf{were} collected} (pluriel) ; \eng{The money \textbf{was} collected} (singulier/indénombrable).
\item \textbf{Faux ami :} \eng{to assist} = aider ; pour “assister à une réunion”, dites \eng{to \textbf{attend} a meeting}.
\item \textbf{Prépositions :} \eng{to donate \textbf{to} a cause} (pas *donate for) ; \eng{to appeal \textbf{for} support} (pas *appeal support) ; \eng{to contribute \textbf{to} a fund}.
\item \textbf{Orthographe :} \eng{fund-raising} (trait d'union souvent utilisé), \eng{cheque} (UK) / \eng{check} (US), \eng{programme} (UK) / \eng{program} (US). Au Cameroun, on suit la norme britannique.
\end{itemize}
\end{attention}

\newpage

\coursec{Méthodes et exercices résolus}

\coursec{Lesson 19 — Reading: Texts on helping with fund-raising activities/support groups in the acquisition of community resources/services}

\courssub{Méthode 1 — Lire et comprendre un texte sur les levées de fonds (skimming et scanning)}

\begin{methode}[Comprendre un texte sur les fund-raising activities]
\begin{enumerate}
\item \cle{Lire le titre et la source} : ils annoncent le thème (une collecte de fonds, un groupe de soutien, un projet communautaire : école, forage, centre de santé). Pose-toi la question : \emph{who is raising funds, for what purpose, and how?}
\item \cle{Skimming} (lecture rapide, 1 minute) : lis le premier et le dernier paragraphe, puis la première phrase de chaque paragraphe. Objectif : dégager l'\emph{idée générale} (main idea), pas les détails.
\item \cle{Repérer les mots du champ lexical} : \eng{to raise funds, donation, donor, charity, committee, borehole, support group, contribution, target, to launch an appeal}. Ils confirment le thème.
\item \cle{Scanning} (lecture ciblée) : pour répondre à une question précise (chiffre, date, nom, montant), repère le \emph{mot-clé} de la question, localise-le dans le texte, puis lis la phrase entière autour.
\item \cle{Déduire le sens des mots inconnus} : utilise le contexte (la phrase avant et après), la famille du mot (\eng{donation} $\to$ \eng{to donate}), les préfixes et suffixes (\eng{re-}, \eng{-ful}, \eng{-less}).
\item \cle{Reformuler} : au Bac, une bonne réponse de compréhension est rédigée avec \emph{tes propres mots}, jamais recopiée mot à mot du texte.
\end{enumerate}
\end{methode}

\begin{exoresolu}[Reading comprehension — type Bac]
\textbf{Texte} : \eng{“Last Saturday, the women of Nkongsamba launched a fund-raising campaign to equip the local health centre. Through door-to-door collections, a charity football match and generous donations from the diaspora, they raised two million CFA francs in a single day — half of their target. The support group announced that the remaining funds would be collected before the end of the year.”}

\textbf{Questions} : (a) What is the text about? (b) How much money was collected, and by what means? (c) Find in the text a word meaning “people living abroad”. (d) What does “target” mean here?
\tcblower
\textbf{(a)} Skimming : le texte parle d'une \emph{collecte de fonds organisée par des femmes pour équiper un centre de santé}. Réponse rédigée : \eng{The text is about a fund-raising campaign organised by the women of Nkongsamba to equip their local health centre.}

\textbf{(b)} Scanning : je repère le mot-clé \eng{two million CFA francs} et les moyens listés. Réponse : \eng{They raised two million CFA francs in one day, through door-to-door collections, a charity football match and donations from the diaspora.}

\textbf{(c)} Le mot cherché est \eng{the diaspora} (« la diaspora »). Indice de contexte : \eng{donations from} + pays étrangers.

\textbf{(d)} Déduction par le contexte : \eng{half of their target} et \eng{the remaining funds} montrent que \eng{target} signifie \emph{l'objectif, la somme visée} (faux ami : ce n'est pas une « cible » à atteindre physiquement). Réponse : \eng{“Target” means the total amount of money they want to collect (their goal).}
\end{exoresolu}

\courssub{Méthode 2 — Répondre aux questions de compréhension et d'interprétation}

\begin{methode}[Rédiger des réponses de compréhension au Bac]
\begin{enumerate}
\item \cle{Identifier le type de question} :
\begin{itemize}
\item \emph{Wh-question} (\eng{Who? What? Why? How?}) $\to$ réponse précise, en phrase complète.
\item \emph{Yes/No question} (\eng{Did they...?}) $\to$ \eng{Yes, they did. / No, they didn't.} + justification tirée du texte.
\item \emph{True/False + justify} $\to$ cite ou paraphrase la ligne qui justifie.
\item \emph{Question d'interprétation} (\eng{In your opinion...? Do you think...?}) $\to$ donne ton avis avec \eng{I think / In my opinion / I believe that...} + un argument.
\end{itemize}
\item \cle{Reprendre le temps verbal de la question} : question au prétérit $\to$ réponse au prétérit ; question au présent $\to$ réponse au présent.
\item \cle{Accorder sujet et verbe} et vérifier la place des mots : sujet + verbe + complément.
\item \cle{Éviter la copie brute} : reformule (\eng{collected} $\to$ \eng{raised} ; \eng{gave money} $\to$ \eng{made donations}).
\item \cle{Relire} chaque réponse : orthographe, majuscule aux noms propres, point final.
\end{enumerate}
\end{methode}

\begin{exoresolu}[Questions variées sur le texte de Nkongsamba]
\textbf{Questions} : (a) Did the women reach their target on Saturday? Justify. (b) Why did the support group organise a football match? (c) True or False: “The campaign is finished.” Justify. (d) In your opinion, is the diaspora important for community projects? (2 sentences)
\tcblower
\textbf{(a)} Question Yes/No au prétérit $\to$ \eng{No, they didn't. They raised only half of their target, so they still need more funds.} (Justification paraphrasée : \eng{half of their target}.)

\textbf{(b)} Question \eng{Why} $\to$ réponse avec \eng{because} ou \eng{to} + base verbale : \eng{They organised a football match to collect money for the health centre.} (But exprimé par \eng{to} + infinitif.)

\textbf{(c)} \eng{False. The text says the remaining funds “would be collected before the end of the year”, so the campaign is still going on.} Attention au temps : le texte annonce une action future $\to$ la campagne n'est pas terminée.

\textbf{(d)} Question d'opinion $\to$ \eng{In my opinion, the diaspora is very important for community projects because its members can send money that the village does not have. I believe their generous donations help schools and health centres to develop faster.} (Deux phrases, avis + argument, comme demandé.)
\end{exoresolu}

\courssub{Méthode 3 — Résumer un texte et donner son opinion}

\begin{methode}[Résumer et réagir en quelques phrases]
\begin{enumerate}
\item \cle{Relever les 3 ou 4 informations essentielles} : qui ? quoi ? comment ? résultat ?
\item \cle{Rédiger un résumé de 2 à 3 phrases} avec des connecteurs : \eng{First... Then... Finally...} ou \eng{The text deals with... / It explains that... / It concludes that...}
\item \cle{Ne pas recopier} de phrases entières : utilise des synonymes et des structures personnelles.
\item \cle{Donner son opinion} ensuite, clairement séparée du résumé : \eng{In my opinion... / I find this initiative admirable because... / I think every community should...}
\item \cle{Conclure} par une phrase d'ouverture ou un souhait : \eng{I hope their project will succeed.}
\end{enumerate}
Structure type à retenir : \eng{The text deals with + noun / The article is about + -ing... According to the writer,... In my opinion,...}
\end{methode}

\begin{exoresolu}[Summary and personal opinion — type Bac]
\textbf{Énoncé} : Summarise the Nkongsamba text in two or three sentences, then give your opinion on fund-raising activities in your community. (about 50 words)
\tcblower
\textbf{Étape 1} : informations essentielles = femmes de Nkongsamba / collecte pour le centre de santé / plusieurs moyens (porte-à-porte, match, diaspora) / 2 millions FCFA, moitié de l'objectif.

\textbf{Étape 2} : résumé rédigé : \eng{The text deals with a fund-raising campaign launched by the women of Nkongsamba to equip their health centre. Thanks to door-to-door collections, a charity football match and donations from the diaspora, they raised two million CFA francs, which is half of their target.}

\textbf{Étape 3} : opinion : \eng{In my opinion, such activities are very useful because they help communities to get resources that the State cannot always provide. I think every village should have a support group like this one.}

\textbf{Vérification} : environ 60 mots, résumé au présent de vérité générale + prétérit pour les faits, opinion introduite par \eng{In my opinion}. Aucune phrase copiée du texte.
\end{exoresolu}

\courssub{Méthode 4 — Maîtriser le lexique des levées de fonds et des ressources communautaires}

\begin{methode}[Utiliser le vocabulaire du fund-raising sans faute]
\begin{enumerate}
\item \cle{Apprendre les mots par familles} :
\begin{itemize}
\item Collecter : \eng{to raise funds, to collect money, a fund-raising campaign, a charity event, to launch an appeal}
\item Donner : \eng{to donate, to make a donation, a donor, a contribution, to contribute to}
\item Soutenir : \eng{a support group, to support, volunteers, to volunteer, solidarity}
\item Ressources : \eng{community resources, a borehole (forage), a health centre, equipment, a classroom block, a water supply}
\end{itemize}
\item \cle{Mémoriser les collocations} : on dit \eng{to raise funds} (jamais \eng{to rise funds}), \eng{to reach a target}, \eng{to equip a health centre}, \eng{to launch a campaign}.
\item \cle{Attention aux prépositions} : \eng{to contribute \textbf{to} a project}, \eng{to rely \textbf{on} donors}, \eng{to be responsible \textbf{for} the collection}, \eng{to participate \textbf{in} an event}.
\item \cle{Distinguer les faux amis} : \eng{to achieve} = accomplir (pas « achever » au sens de tuer) ; \eng{a donation} = un don ; \eng{actually} = en fait (pas « actuellement » $\to$ \eng{currently}).
\item \cle{Réutiliser le lexique} dans tes propres phrases pour l'ancrer : rédige 3 phrases sur un projet de ton village.
\end{enumerate}
\end{methode}

\begin{exemplebox}[Tableau récapitulatif : Vocabulaire clé du thème]
\begin{center}
\begin{tabularx}{\linewidth}{|X|X|X|}
\hline
\rowcolor{popblueL} \textbf{English} & \textbf{Français} & \textbf{Exemple / Collocation} \\
\hline
\eng{to raise funds} & lever des fonds & \eng{raise funds for a borehole} \\
\eng{a fund-raising campaign} & une campagne de collecte & \eng{launch a fund-raising campaign} \\
\eng{a donation / to donate} & un don / faire un don & \eng{make a generous donation} \\
\eng{a donor} & un donateur & \eng{local and foreign donors} \\
\eng{a support group} & un groupe de soutien & \eng{the village support group} \\
\eng{a borehole} & un forage (eau) & \eng{drill a borehole} \\
\eng{a target} & un objectif (chiffre) & \eng{reach the target} \\
\eng{the diaspora} & la diaspora & \eng{donations from the diaspora} \\
\eng{to contribute to} & contribuer à & \eng{contribute 5,000 francs to the project} \\
\eng{volunteers} & des bénévoles & \eng{volunteers worked for free} \\
\hline
\end{tabularx}
\end{center}
\end{exemplebox}

\begin{exoresolu}[Vocabulary in context]
\textbf{Énoncé} : Complete with words from the lesson: \emph{raise, donation, volunteers, target, borehole, contribute}. (a) The village needs a new \_\_\_ to provide clean water. (b) Many \_\_\_ helped to build the classroom block. (c) Every family was asked to \_\_\_ 5,000 francs to the project. (d) Their \_\_\_ is ten million francs. (e) A generous businessman made a large \_\_\_. (f) They organised a concert to \_\_\_ funds.
\tcblower
\textbf{(a)} \eng{borehole} — un forage fournit de l'eau potable (\eng{clean water}).
\textbf{(b)} \eng{volunteers} — ce sont des bénévoles qui aident gratuitement ; pluriel car \eng{many}.
\textbf{(c)} \eng{contribute} — structure : \eng{to contribute + somme + to a project} ; après \eng{was asked to}, on met la base verbale.
\textbf{(d)} \eng{target} — l'objectif financier visé ; singulier avec \eng{their} car un seul objectif commun.
\textbf{(e)} \eng{donation} — après l'adjectif \eng{large}, il faut un nom ; \eng{to make a donation} = faire un don.
\textbf{(f)} \eng{raise} — collocation obligatoire : \eng{to raise funds} ; après \eng{to} (but), base verbale. Piège : \eng{to rise} (s'élever) est intransitif et ne prend jamais de complément.
\end{exoresolu}

\courssub{Méthode 5 — Réviser le discours rapporté et le passif en contexte (Révision Bac)}

\begin{propriete}[Règles de transformation : Discours rapporté et Voix passive]
\textbf{1. Discours rapporté (Reported Speech)} avec un verbe introducteur au prétérit (\eng{said, announced, explained, added}) :
\begin{itemize}
\item \cle{Recule les temps} (backshifting) :
\begin{center}
\begin{tabular}{|l|l|}
\hline
\textbf{Direct} & \textbf{Rapporté} \\
\hline
Present Simple (\eng{do/does}) & Past Simple (\eng{did}) \\
Present Continuous (\eng{am/is/are doing}) & Past Continuous (\eng{was/were doing}) \\
Present Perfect (\eng{have/has done}) & Past Perfect (\eng{had done}) \\
Past Simple (\eng{did}) & Past Perfect (\eng{had done}) \\
\eng{will} & \eng{would} \\
\eng{can} & \eng{could} \\
\eng{must} & \eng{had to} \\
\hline
\end{tabular}
\end{center}
\item \cle{Change les repères de temps et de lieu} : \eng{today $\to$ that day ; now $\to$ then ; here $\to$ there ; this $\to$ that ; tomorrow $\to$ the next day ; yesterday $\to$ the day before ; last week $\to$ the week before}.
\item \cle{Adapte les pronoms et possessifs} : \eng{I $\to$ he/she ; we $\to$ they ; my $\to$ his/her ; our $\to$ their}.
\end{itemize}

\textbf{2. Voix passive (Passive Voice)} :
\begin{itemize}
\item Objet du verbe actif $\to$ Sujet du verbe passif.
\item Verbe = \emph{be} (conjugué au temps de la phrase active) + \emph{participe passé} du verbe.
\item Agent (sujet actif) introduit par \eng{by} (souvent omis s'il est évident ou inconnu).
\item \cle{Accord} : \emph{be} s'accorde avec le nouveau sujet (singulier/pluriel).
\item \cle{Verbes intransitifs} (\eng{happen, arrive, die, sleep, go}) : \textbf{jamais} de passif (pas d'objet).
\end{itemize}
\end{propriete}

\begin{exoresolu}[Grammar transformations — type Bac]
\textbf{Énoncé} : (a) Put into reported speech: The chairwoman said: “We will collect the remaining funds before the end of the year.” (b) Put into the passive: “The women of Nkongsamba equipped the health centre.” (c) Put into the passive: “The diaspora has donated two million CFA francs.”
\tcblower
\textbf{(a)} Verbe introducteur au prétérit (\eng{said}) $\to$ recul des temps et changement des pronoms/repères : \eng{The chairwoman said (that) they would collect the remaining funds before the end of the year.} Justifications : \eng{we $\to$ they} ; \eng{will $\to$ would} ; \eng{the end of the year} inchangé (repère absolu, discours rapporté la même année).

\textbf{(b)} Objet \eng{the health centre} devient sujet ; verbe actif au prétérit $\to$ \emph{be} au prétérit (\eng{was}) + participe passé \eng{equipped} : \eng{The health centre was equipped by the women of Nkongsamba.} Accord singulier : \eng{was} (pas \eng{were}).

\textbf{(c)} Temps actif = present perfect (\eng{has donated}) $\to$ passif = \eng{has been} + participe passé : \eng{Two million CFA francs have been donated by the diaspora.} Sujet \eng{two million CFA francs} pluriel $\to$ \eng{have been}. Participe passé de \eng{donate} : \eng{donated} (verbe régulier).
\end{exoresolu}

\courssub{Méthode 6 — Rédiger un appel aux dons (fund-raising appeal letter)}

\begin{methode}[Écrire une lettre d'appel aux dons (formal letter)]
\begin{enumerate}
\item \cle{Respecter la mise en page formelle} :
\begin{itemize}
\item Adresse de l'expéditeur (haut gauche) + Date (haut droite).
\item Adresse du destinataire (gauche, sous la date).
\item Objet (\eng{Re: / Subject:}) clair : \eng{Fund-raising Appeal for the Construction of a Borehole}.
\item Formule d'appel : \eng{Dear Sir or Madam,} (si nom inconnu) ou \eng{Dear Mr/Mrs [Nom],}.
\item Paragraphes courts, séparés par une ligne blanche.
\item Formule de politesse finale : \eng{Yours faithfully,} (si \eng{Dear Sir or Madam}) ou \eng{Yours sincerely,} (si nom connu).
\item Signature manuscrite + Nom imprimé + Fonction.
\end{itemize}
\item \cle{Paragraphe 1 (Introduction)} : se présenter et annoncer le but. \eng{I am writing on behalf of the support group of [Village] to launch an appeal for funds.}
\item \cle{Paragraphe 2 (Le projet et le besoin)} : expliquer le problème et la solution. \eng{Our community lacks clean water. Women walk 5 km daily. We plan to construct a borehole near the school. The estimated cost is [Montant] CFA francs.}
\item \cle{Paragraphe 3 (État d'avancement)} : ce qui est fait + ce qui manque. \eng{So far, we have raised [Montant] through [activités]. We still need [Montant] to reach our target.}
\item \cle{Paragraphe 4 (Appel et politesse)} : conditionnel de politesse + modalités de don. \eng{We would be grateful for any contribution, however small. Donations can be sent to [adresse/compte]. Thank you for your generosity.}
\item \cle{Temps et modes clés} : Present Perfect (\eng{we have raised}) pour le passé récent/acquis ; Conditionnel (\eng{would be, could help}) pour la politesse ; Futur (\eng{will benefit}) pour l'impact futur.
\end{enumerate}
\end{methode}

\begin{exoresolu}[Guided writing — appeal letter (about 120 words)]
\textbf{Énoncé} : You are the secretary of the support group of your village. Write a letter to a local businessman asking him to contribute to the construction of a borehole.
\tcblower
\textbf{Modèle rédigé} :

\eng{Mbalmayo, 12th March 2025}

\eng{The Manager}

\eng{Etoundi Enterprises}

\eng{P.O. Box 145, Mbalmayo}

\eng{\textbf{Re: Fund-raising Appeal for the Construction of a Borehole}}

\eng{Dear Sir or Madam,}

\eng{I am writing on behalf of the support group of our village to appeal for your generous contribution. Our community of about three thousand people has no clean water supply, and the women walk five kilometres every day to fetch water.}

\eng{We have decided to construct a borehole near the primary school. The project will cost four million CFA francs. So far, we have raised two and a half million through a charity concert and door-to-door collections.}

\eng{We would be very grateful if you could help us to reach our target. Any donation, however small, will be warmly welcomed.}

\eng{Thank you for your kindness and generosity.}

\eng{Yours faithfully,}

\eng{(Signature)}

\eng{Ngono Marie, Secretary}

\textbf{Justifications} : Mise en page formelle (adresses, objet, \eng{Dear Sir or Madam} $\to$ \eng{Yours faithfully}) ; Paragraphe 1 : présentation + but ; Paragraphe 2 : besoin + chiffres + \eng{will cost} (futur) ; Paragraphe 3 : \eng{have raised} (present perfect = acquis) ; Paragraphe 4 : conditionnel \eng{would be grateful if you could} (politesse) + \eng{however small} (concession).
\end{exoresolu}

\courssub{Méthode 7 — Expression orale : Simuler une réunion communautaire (Role-play)}

\begin{experience}[Simulation de réunion communautaire — Community Meeting Role-play]
\textbf{Contexte} : Le groupe de soutien du village organise une assemblée générale pour présenter l'avancement du projet de forage et motiver de nouvelles contributions.
\textbf{Rôles (groupes de 4-5 élèves)} :
\begin{enumerate}
\item \textbf{The Chairperson} : ouvre la séance, présente le rapport moral/financier (\eng{We have raised... We need...}).
\item \textbf{The Treasurer} : donne les chiffres exacts (\eng{Total target: 10 million. Collected: 4 million. Balance: 6 million}).
\item \textbf{A Village Elder} : insiste sur l'urgence (\eng{Our women suffer. The dry season is coming. We must act now.}).
\item \textbf{A Youth Representative} : propose des idées nouvelles (\eng{Let's organise a cultural night / a sponsored walk. We can use social media to reach the diaspora.}).
\item \textbf{A Sceptical Villager} : pose des questions critiques (\eng{How will the money be managed? What if the pump breaks down? Is there transparency?}).
\end{enumerate}
\textbf{Consignes linguistiques} :
\begin{itemize}
\item Utiliser le vocabulaire de la leçon : \eng{target, borehole, diaspora, contribution, transparency, accountability, to raise funds, to reach a goal}.
\item Structures utiles : \eng{I propose that we... / I suggest + V-ing / Why don't we...? / In my view... / I am concerned about... / We must ensure that...}
\item La présidence doit conclure : \eng{Let's vote on this proposal. / The meeting is adjourned.}
\end{itemize}
\textbf{Déroulement} : 5 min préparation (notes en anglais) $\to$ 10 min jeu de rôle $\to$ 5 min feedback (professeur + pairs) sur la prononciation, le vocabulaire, la grammaire (conditionnels, modaux, passif).
\end{experience}

\begin{savaistu}[Le saviez-vous ? — Le \eng{Njangi} / \eng{Tontine} au Cameroun]
Au Cameroun, le système traditionnel d'entraide financière s'appelle le \eng{Njangi} (ou \eng{Tontine}, \eng{Meeting}, \eng{Esusu} au Nigeria). C'est une forme ancestrale de \eng{community fund-raising} : les membres cotisent chaque semaine/mois, et la cagnotte est donnée à tour de rôle à l'un d'eux pour un gros projet (construction, commerce, frais scolaires, santé). Ce système repose sur la \eng{trust} (confiance) et la \eng{solidarity} (solidarité), sans contrat écrit ni banque. Aujourd'hui, beaucoup de \eng{support groups} modernes s'inspirent du \eng{Njangi} pour lancer des \eng{appeals} plus larges (forages, salles de classe, centres de santé) en y ajoutant la transparence (comptes rendus, trésoriers élus) et l'ouverture à la \eng{diaspora} via \eng{Mobile Money} (Orange Money, MTN MoMo). C'est un bel exemple d'\eng{endogenous development} (développement endogène) !
\end{savaistu}

\begin{attention}[Les 5 erreurs qui coûtent des points au Bac]
\begin{enumerate}
\item \cle{Confondre \eng{to raise} et \eng{to rise}} : on \emph{lève des fonds} avec \eng{to raise funds} (verbe transitif : \eng{raise, raised, raised}) ; \eng{to rise} (s'élever, augmenter — \eng{rise, rose, risen}) est intransitif et ne prend \textbf{jamais} de complément. Jamais \eng{*rise money} ni \eng{*risen funds}.
\item \cle{Faux ami \eng{actually}} : il signifie « en fait, en réalité » (\eng{Actually, I don't know him.}). Pour « actuellement », dis \eng{currently} ou \eng{at present}. De même, \eng{a target} = un objectif chiffré (pas une « cible » de tir dans ce contexte) ; \eng{to achieve a target} = atteindre un objectif.
\item \cle{Oublier le recul des temps au discours rapporté} : après \eng{said/announced} au prétérit, \eng{will} devient \eng{would}, \eng{can} devient \eng{could}, le présent devient prétérit, le prétérit/present perfect deviennent past perfect. Change aussi les pronoms (\eng{we $\to$ they}) et les repères (\eng{tomorrow $\to$ the next day}).
\item \cle{Mal former le passif} : il faut \emph{be} (au bon temps, accordé au sujet) + \textbf{participe passé}. \eng{The health centre \textbf{was equipped}} (singulier, prétérit) ; \eng{The funds \textbf{have been collected}} (pluriel, present perfect). Attention aux participes irréguliers : \eng{given, built, sent, made, spent, paid}.
\item \cle{Prépositions et ordre des mots calqués du français} : \eng{contribute \textbf{to}}, \eng{participate \textbf{in}}, \eng{rely \textbf{on}}, \eng{listen \textbf{to}}, \eng{wait \textbf{for}}. Adverbe de manière : \eng{They \textbf{generously donated} money} (avant le verbe) ou \eng{They donated money \textbf{generously}} (à la fin), mais \textbf{jamais} \eng{*donated generously money}. Majuscule obligatoire à \eng{I}, aux noms propres (villages, jours, mois, nationalités).
\end{enumerate}
\end{attention}

\begin{exercice}[Entraînement — Compréhension, Grammaire, Vocabulaire, Expression]
\textbf{A. Reading Comprehension} \\
Read the following text and answer the questions.

\eng{“The Youth Association of Buea Town has launched a massive fund-raising drive to renovate the community hall, which was damaged by a storm last year. The President, Mr. Ekema, explained that the hall serves as a meeting point, a vaccination centre and a venue for cultural events. ‘We cannot wait for the government to do everything,’ he said. ‘We have decided to help ourselves.’ So far, they have organised a sponsored walk and a cultural night, raising 1.5 million CFA francs. The target is 5 million. Mr. Ekema added that the diaspora branch of the association had promised to match every franc raised locally before December.”}

\textbf{Questions} :
\begin{enumerate}
\item What is the purpose of the fund-raising drive? (1 sentence)
\item How much money has been raised so far? What is the target?
\item Find in the text a synonym for “renovate” (paragraph 1) and a word meaning “people from the same origin living abroad” (last paragraph).
\item True or False? Justify by quoting the text.
\begin{enumerate}
\item The community hall is only used for meetings.
\item The government has already started the renovation.
\item The diaspora will double the local contributions.
\end{enumerate}
\item In your opinion, why is self-help important for a community? (2 sentences)
\end{enumerate}

\textbf{B. Grammar} \\
Rewrite the following sentences as indicated.
\begin{enumerate}
\item The President said: “We will finish the work before Christmas.” (Put into reported speech)
\item The youth organised a sponsored walk. (Put into the passive voice)
\item The diaspora has promised a large donation. (Put into the passive voice)
\item “We are collecting donations today,” the treasurer announced. (Put into reported speech)
\item The storm destroyed the roof last year. (Put into the passive voice)
\end{enumerate}

\textbf{C. Vocabulary} \\
Fill in the blanks with the correct word from the box. Two words are extra.
\eng{\{ borehole, target, donation, volunteers, contribute, launch, rise, equipment, diaspora, charity \}}
\begin{enumerate}
\item The support group decided to \_\_\_\_\_\_ a new appeal for funds next Monday.
\item Every member of the community was asked to \_\_\_\_\_\_ at least 2,000 francs.
\item The \_\_\_\_\_\_ for the new classroom block is ten million CFA francs.
\item Many young \_\_\_\_\_\_ helped to carry bricks and sand to the site.
\item A generous \_\_\_\_\_\_ from a local businessman allowed them to buy cement.
\item The \_\_\_\_\_\_ in Douala and Yaoundé sent money through Mobile Money.
\end{enumerate}

\textbf{D. Writing} \\
You are the Treasurer of your school's Environmental Club. Write a formal letter (120–150 words) to the Principal of a partner school abroad, asking for a donation to buy waste bins and create a school garden. Use the layout and expressions studied in Méthode 6.
\end{exercice}

\begin{corrige}[Exercice Entraînement]
\textbf{A. Reading Comprehension}
\begin{enumerate}
\item The purpose is to renovate the community hall damaged by a storm.
\item They have raised 1.5 million CFA francs so far. The target is 5 million CFA francs.
\item Synonym for “renovate”: \eng{renovate} (texte) $\to$ \eng{repair} / \eng{restore} (non dans le texte, mais sens). Dans le texte : \eng{renovate} est le mot utilisé. Pour “people from same origin living abroad” : \eng{the diaspora} (ou \eng{diaspora branch}).
\item (a) False. \eng{The hall serves as a meeting point, a vaccination centre and a venue for cultural events.}
\item (b) False. \eng{‘We cannot wait for the government to do everything,’ he said.}
\item (c) True. \eng{the diaspora branch... had promised to match every franc raised locally.}
\item (Sample) In my opinion, self-help is important because it empowers the community to solve its own problems without depending entirely on the government. It also strengthens unity and solidarity among the villagers.
\end{enumerate}

\textbf{B. Grammar}
\begin{enumerate}
\item The President said (that) they would finish the work before Christmas.
\item A sponsored walk was organised by the youth.
\item A large donation has been promised by the diaspora.
\item The treasurer announced that they were collecting donations that day.
\item The roof was destroyed by the storm last year.
\end{enumerate}

\textbf{C. Vocabulary}
\begin{enumerate}
\item \eng{launch}
\item \eng{contribute}
\item \eng{target}
\item \eng{volunteers}
\item \eng{donation}
\item \eng{diaspora}
\end{enumerate}
(Mots non utilisés : \eng{borehole, rise, equipment, charity})

\textbf{D. Writing} (Modèle indicatif)
\eng{Yaoundé, 20th October 2025}

\eng{The Principal}

\eng{Greenfield High School}

\eng{London, UK}

\eng{\textbf{Re: Appeal for Donation — School Garden and Waste Management Project}}

\eng{Dear Sir or Madam,}

\eng{I am writing on behalf of the Environmental Club of Government Bilingual High School, Yaoundé. Our club aims to promote a clean and green school environment. We have launched a project to purchase fifty waste bins and create a school garden to teach students sustainable agriculture.}

\eng{The total cost of the project is estimated at three million CFA francs. So far, we have raised one million through school-based activities like a “Green Day” and sales of recycled crafts. However, we still have a shortfall of two million francs.}

\eng{We would be extremely grateful if your school could support our initiative with a financial contribution. Any amount would make a significant difference. We will send you a detailed report and photos once the project is completed.}

\eng{Thank you very much for considering our request.}

\eng{Yours faithfully,}

\eng{(Signature)}

\eng{Ayuk John, Treasurer, Environmental Club}
\end{corrige}

\newpage

\coursec{Exercices}

\courssub{Je vérifie mes connaissances}

\begin{exercice}[QCM — Vocabulaire du fund-raising]
Choisis la bonne réponse pour chaque phrase.
\begin{enumerate}
\item Money collected for a good cause is called:
\begin{enumerate}
\item a refund \quad b) a donation \quad c) a salary \quad d) a debt
\end{enumerate}
\item A person who gives money regularly to a charity is:
\begin{enumerate}
\item a sponsor \quad b) a customer \quad c) a debtor \quad d) a burglar
\end{enumerate}
\item To \eng{raise funds} means:
\begin{enumerate}
\item to spend money \quad b) to collect money \quad c) to lose money \quad d) to hide money
\end{enumerate}
\item A \eng{fund-raising event} at school could be:
\begin{enumerate}
\item a strike \quad b) a charity concert \quad c) an exam \quad d) a detention
\end{enumerate}
\end{enumerate}
\end{exercice}

\begin{exercice}[Vrai ou Faux — Compréhension de notions]
Lis les affirmations suivantes et dis si elles sont \eng{True} ou \eng{False}, puis justifie ta réponse en une phrase en anglais.
\begin{enumerate}
\item \eng{A support group only helps people financially.}
\item \eng{Transparency is important when a community collects money.}
\item \eng{Crowdfunding means borrowing money from a bank.}
\item \eng{A borehole provides a village with clean water.}
\item \eng{Fund-raising activities always require the help of the government.}
\end{enumerate}
\end{exercice}

\begin{exercice}[Vocabulaire — Trouve le mot]
Complète chaque définition par le mot anglais qui convient, choisi dans la liste : \eng{pledge — borehole — charity — volunteer — committee — grant}.
\begin{enumerate}
\item A person who works without being paid: \dotfill
\item A deep hole drilled to get water: \dotfill
\item A formal promise to give money: \dotfill
\item An organisation that helps people in need: \dotfill
\item A group of people elected to manage a project: \dotfill
\item Money given by an institution for a specific project, which does not have to be paid back: \dotfill
\end{enumerate}
\end{exercice}

\begin{exercice}[Grammaire — Conjugaison à compléter]
Complète les phrases avec le verbe entre parenthèses à la forme correcte (passif ou discours rapporté).
\begin{enumerate}
\item Last year, a new classroom block \dotfill (build) by the parents' association.
\item The chairman said that the money \dotfill (keep) in a bank account.
\item The funds \dotfill (raise) during the charity walk next Saturday.
\item The treasurer announced that they \dotfill (publish) the accounts every month.
\item So far, two boreholes \dotfill (drill) in the village.
\end{enumerate}
\end{exercice}

\courssub{Je m'entraîne}

\begin{exercice}[Traduction — Français vers anglais]
Traduis les phrases suivantes en anglais correct.
\begin{enumerate}
\item La communauté a collecté de l'argent pour construire un forage.
\item Les habitants ont dit qu'ils organiseraient une vente de charité samedi.
\item Une somme importante a été promise par une entreprise de la région.
\item Les comptes de l'association doivent être publiés chaque année.
\item Nous avons besoin de bénévoles pour distribuer les repas.
\end{enumerate}
\end{exercice}

\begin{exercice}[Transformation de phrases]
Réécris chaque phrase selon la consigne entre parenthèses.
\begin{enumerate}
\item \eng{The villagers raised two million francs for the health centre.} (mets au passif)
\item \eng{The manager said: “We will sponsor your project.”} (mets au discours rapporté)
\item \eng{“The committee has opened a bank account,” the treasurer explained.} (mets au discours rapporté)
\item \eng{A local company donated fifty bags of cement.} (mets au passif)
\item \eng{The secretary asked: “When will the event take place?”} (mets au discours rapporté)
\end{enumerate}
\end{exercice}

\begin{exercice}[Texte à trous — Running an association]
Complète le texte suivant avec les mots de la liste : \eng{accountability — grant — pledge — outreach — stakeholders — sustainable}.

\begin{textebox}[Our association, one year later]
When our association received a (1) \dotfill from the city council, all the (2) \dotfill — members, donors and local authorities — were invited to a meeting. The president insisted on (3) \dotfill : every franc spent would be recorded and published. Several members made a (4) \dotfill to contribute monthly. Thanks to our (5) \dotfill programme, we visited remote neighbourhoods to inform families about our services. Our goal is to make the project (6) \dotfill , so that it can continue even without external funding.
\end{textebox}
\end{exercice}

\begin{exercice}[Appariement — Mots et définitions]
Associe chaque mot anglais (1--6) à sa définition (a--f).
\begin{enumerate}
\item crowdfunding
\item borehole
\item sponsor
\item handover ceremony
\item raffle
\item budget
\end{enumerate}
\begin{enumerate}[label=\alph*)]
\item an event where a finished project is officially given to the community
\item a detailed plan of how money will be spent
\item collecting small amounts of money from many people, usually online
\item a person or company that pays for an event in exchange for publicity
\item a deep narrow hole made in the ground to find water
\item a competition in which people buy tickets to win prizes, often for charity
\end{enumerate}
\end{exercice}

\courssub{Je me prépare au Bac}

\begin{exercice}[Type Bac — Compréhension de l'écrit]
Lis le texte suivant, puis réponds aux questions en anglais, avec tes propres mots autant que possible.

\begin{textebox}[Crowdfunding for a tech hub in Lagos]
When twenty-six-year-old Adaeze Obi returned to Lagos after her studies in Accra, she was shocked to discover that young graduates in her neighbourhood had nowhere to learn digital skills. Internet cafés were expensive, and the few computer training centres charged fees that most families could not afford. Adaeze decided to act. With three friends, she launched an online crowdfunding campaign called “Code Lagos Youth”, aiming to raise the equivalent of eight thousand euros to open a small tech hub in a rented shop.

The campaign page told the story of Chinedu, a nineteen-year-old who had taught himself basic programming on a borrowed smartphone but had nowhere to practise. Within two weeks, more than four hundred donors from Nigeria and abroad had contributed. Some gave the price of a cup of coffee; a Nigerian company based in London pledged a larger amount and promised to double every donation received during the final week.

Not everything went smoothly. Some residents doubted that the money would be used properly. To answer these concerns, Adaeze published a detailed budget online and promised monthly financial reports. “Trust is our most valuable resource,” she told a local radio station. “If people can see where every naira goes, they will give again.”

Six months later, the tech hub opened its doors with twelve computers, a fast internet connection and two volunteer trainers. Forty young people enrolled in the first free course. The local traditional ruler, who had initially been sceptical, attended the opening ceremony and offered to pay the rent for the second year. For Adaeze, the real success is not the money raised but the community that has formed around the project: “The hub belongs to everyone here. That is why it will survive.”
\end{textebox}

\textbf{Questions}
\begin{enumerate}
\item Say whether the following statements are \eng{True} or \eng{False}. Justify each answer with a quotation from the text.
\begin{enumerate}
\item Adaeze started her project because computer training was too expensive for local families.
\item All the donors gave large amounts of money.
\item The traditional ruler supported the project from the very beginning.
\end{enumerate}
\item Answer these questions in your own words.
\begin{enumerate}
\item What problem did young graduates in Adaeze's neighbourhood face?
\item How did Adaeze react to the residents' doubts about the use of the money?
\item According to Adaeze, what is the real success of the project?
\end{enumerate}
\item Find in the text words or expressions equivalent to the following:
\begin{enumerate}
\item promised to give (paragraph 2)
\item money collected (paragraph 4)
\item registered, signed up (paragraph 5)
\end{enumerate}
\item In about 50 words, summarise how Adaeze's project succeeded.
\end{enumerate}
\end{exercice}

\begin{exercice}[Type Bac — Langue : grammaire et traduction]
\textbf{A. Grammar.} Transform the following sentences as indicated.
\begin{enumerate}
\item \eng{The youth club organised a fund-raising gala last month.} (Put into the passive voice.)
\item \eng{The director said: “We shall build a new library next year.”} (Put into reported speech.)
\item \eng{“Have you received the grant?” the journalist asked the president.} (Put into reported speech.)
\item \eng{The community has already drilled two boreholes.} (Put into the passive voice.)
\item \eng{If the company sponsors us, we will buy new computers.} (Rewrite beginning with \eng{Unless}.)
\end{enumerate}

\textbf{B. Translation.} Traduis en anglais.
\begin{enumerate}
\item Les villageois ont annoncé qu'ils construiraient un centre de santé avec l'argent collecté.
\item Chaque don, même petit, est important pour notre projet.
\item Les comptes de l'association ont été publiés pour garantir la transparence.
\end{enumerate}
\end{exercice}

\begin{exercice}[Type Bac — Expression écrite]
\textbf{Topic.} Your village development association wants to build a community library. As the secretary of the association, write a formal appeal letter to the manager of a mining company operating nearby, requesting sponsorship. In your letter, you should:
\begin{itemize}
\item present your association and the project;
\item explain why the library is useful for the community;
\item specify what kind of support you are requesting (money, books, materials);
\item reassure the company about transparency and mention possible publicity for the sponsor.
\end{itemize}
Write between \textbf{200 and 250 words}. Use the appropriate layout of a formal letter (addresses, date, salutation, subscription).
\end{exercice}

\newpage

\coursec{Corrigés des exercices}

\begin{corrige}[Exercice 1 — QCM : Vocabulaire du fund-raising]
\begin{enumerate}
\item \textbf{b) a donation} — \eng{Money collected for a good cause is called a donation.} (Un \eng{refund} est un remboursement, un \eng{salary} un salaire, une \eng{debt} une dette.)
\item \textbf{a) a sponsor} — \eng{A person who gives money regularly to support a charity is a sponsor.} (On peut aussi dire \eng{a donor} ou \eng{a regular contributor}, mais parmi les choix proposés, seul \eng{sponsor} convient. Remarque : dans un autre contexte, un \eng{sponsor} est aussi une entreprise qui finance un événement en échange de publicité.)
\item \textbf{b) to collect money} — \eng{To raise funds means to collect money for a cause.} Attention : \eng{raise} est ici un faux ami de « élever » ; dans ce contexte, il signifie « collecter, réunir ». Verbe irrégulier ? Non : \emph{raise -- raised -- raised} (régulier), à ne pas confondre avec \emph{rise -- rose -- risen} (« s'élever »).
\item \textbf{b) a charity concert} — \eng{A fund-raising event at school could be a charity concert.} (Une \eng{strike} est une grève, une \eng{detention} une retenue : rien à voir avec une collecte de fonds.)
\end{enumerate}
\end{corrige}

\begin{corrige}[Exercice 2 — Vrai ou Faux]
\begin{enumerate}
\item \textbf{False.} \eng{A support group can also offer emotional support, advice, training or practical help, not only money.}
\item \textbf{True.} \eng{Transparency is important because donors and members must know how the money collected is used.}
\item \textbf{False.} \eng{Crowdfunding means collecting small amounts of money from many people, usually online; it is not a bank loan.}
\item \textbf{True.} \eng{A borehole is a deep hole drilled into the ground to provide a village with clean water.}
\item \textbf{False.} \eng{Fund-raising activities can be organised by communities, associations or individuals without any government help.}
\end{enumerate}
\end{corrige}

\begin{corrige}[Exercice 3 — Vocabulaire : Trouve le mot]
\begin{enumerate}
\item \textbf{volunteer} — une personne qui travaille sans être payée.
\item \textbf{borehole} — un trou profond foré pour obtenir de l'eau (prononciation : « bo-hôle »).
\item \textbf{pledge} — une promesse formelle de donner de l'argent.
\item \textbf{charity} — une organisation qui aide les personnes dans le besoin.
\item \textbf{committee} — un groupe de personnes élues pour gérer un projet.
\item \textbf{grant} — une subvention, c'est-à-dire une somme donnée pour un projet précis et qu'il n'est pas nécessaire de rembourser.
\end{enumerate}
\end{corrige}

\begin{corrige}[Exercice 4 — Grammaire : passif et discours rapporté]
\begin{enumerate}
\item \eng{Last year, a new classroom block \textbf{was built} by the parents' association.} — Passif au prétérit simple : \emph{was/were + participe passé}. \eng{Build} est irrégulier : \emph{build -- built -- built}.
\item \eng{The chairman said that the money \textbf{would be kept} in a bank account.} — Discours rapporté : recul du futur (\emph{will} $\rightarrow$ \emph{would}) + passif. \eng{Keep} est irrégulier : \emph{keep -- kept -- kept}.
\item \eng{The funds \textbf{will be raised} during the charity walk next Saturday.} — Futur passif : \emph{will be + participe passé} (l'indice \eng{next Saturday} impose le futur).
\item \eng{The treasurer announced that they \textbf{would publish} the accounts every month.} — Discours rapporté : \emph{will} $\rightarrow$ \emph{would} après un verbe introducteur au passé (\eng{announced}).
\item \eng{So far, two boreholes \textbf{have been drilled} in the village.} — Present perfect passif : \emph{have/has been + participe passé} ; l'indice \eng{so far} (« jusqu'à présent ») impose le present perfect.
\end{enumerate}
\end{corrige}

\begin{corrige}[Exercice 5 — Traduction : français vers anglais]
\begin{enumerate}
\item \eng{The community raised money to build a borehole.} — « Collecter de l'argent » pour une cause se dit \eng{raise money} ; on peut aussi dire \eng{collected money for a borehole}, mais \eng{raised money to build} est la formulation la plus naturelle.
\item \eng{The inhabitants said (that) they would organise a charity sale on Saturday.} — Discours rapporté : futur dans le passé \eng{would organise} ; « vente de charité » = \eng{charity sale}.
\item \eng{A large sum of money was promised by a local company.} — Passif au prétérit ; « une entreprise de la région » = \eng{a local company}.
\item \eng{The association's accounts must be published every year.} — Passif avec modal : \emph{must be + participe passé}.
\item \eng{We need volunteers to distribute the meals.} — « Bénévoles » = \eng{volunteers} ; attention à ne pas écrire \eng{voluntaries}.
\end{enumerate}
\end{corrige}

\begin{corrige}[Exercice 6 — Transformation de phrases]
\begin{enumerate}
\item \eng{Two million francs \textbf{were raised} by the villagers for the health centre.} — Passif au prétérit ; le sujet devient \eng{two million francs} (pluriel $\rightarrow$ \emph{were}).
\item \eng{The manager said (that) they would sponsor our project.} — Recul des temps : \emph{will} $\rightarrow$ \emph{would} ; \emph{we} $\rightarrow$ \emph{they} ; \emph{your} $\rightarrow$ \emph{our}.
\item \eng{The treasurer explained (that) the committee had opened a bank account.} — Present perfect $\rightarrow$ past perfect (\emph{has opened} $\rightarrow$ \emph{had opened}).
\item \eng{Fifty bags of cement \textbf{were donated} by a local company.} — Passif au prétérit, pluriel (\emph{bags} $\rightarrow$ \emph{were}).
\item \eng{The secretary asked when the event would take place.} — Question rapportée : pas d'inversion, pas de \eng{if} car c'est une question en \emph{wh-} ; \emph{will} $\rightarrow$ \emph{would}.
\end{enumerate}
\end{corrige}

\begin{corrige}[Exercice 7 — Texte à trous : Running an association]
\begin{enumerate}
\item \textbf{grant} — \eng{received a grant from the city council} (une subvention du conseil municipal).
\item \textbf{stakeholders} — les parties prenantes : membres, donateurs et autorités locales.
\item \textbf{accountability} — \eng{insisted on accountability} : l'obligation de rendre compte ; l'indice est la phrase suivante (chaque franc dépensé serait enregistré et publié).
\item \textbf{pledge} — \eng{made a pledge to contribute monthly} : une promesse de contribution mensuelle.
\item \textbf{outreach} — \eng{outreach programme} : programme de proximité / de sensibilisation sur le terrain (visiter les quartiers éloignés).
\item \textbf{sustainable} — \eng{make the project sustainable} : rendre le projet durable, viable sans financement extérieur.
\end{enumerate}
\end{corrige}

\begin{corrige}[Exercice 8 — Appariement : mots et définitions]
\begin{center}
\begin{tabularx}{\linewidth}{|l|X|}\hline
\rowcolor{popblueL} Mot & Définition \\ \hline
1. crowdfunding & c) collecting small amounts of money from many people, usually online \\ \hline
2. borehole & e) a deep narrow hole made in the ground to find water \\ \hline
3. sponsor & d) a person or company that pays for an event in exchange for publicity \\ \hline
4. handover ceremony & a) an event where a finished project is officially given to the community \\ \hline
5. raffle & f) a competition in which people buy tickets to win prizes, often for charity \\ \hline
6. budget & b) a detailed plan of how money will be spent \\ \hline
\end{tabularx}
\end{center}
Réponses : \textbf{1-c, 2-e, 3-d, 4-a, 5-f, 6-b.}
\end{corrige}

\begin{corrige}[Exercice 9 — Type Bac : Compréhension de l'écrit]
\textbf{Barème indicatif (sur 20) :} question 1 : 6 pts (1 pt par réponse + 1 pt par citation) ; question 2 : 6 pts (2 pts par réponse) ; question 3 : 3 pts (1 pt par mot) ; question 4 : 5 pts (contenu 3 pts, langue 2 pts).

\textbf{1. True or False with quotations}
\begin{enumerate}
\item \textbf{True.} \eng{“The few computer training centres charged fees that most families could not afford.”}
\item \textbf{False.} \eng{“Some gave the price of a cup of coffee”} — la plupart des dons étaient modestes ; seule une entreprise a promis une somme plus importante.
\item \textbf{False.} \eng{“The local traditional ruler, who had initially been sceptical...”} — il n'a soutenu le projet qu'à la cérémonie d'ouverture.
\end{enumerate}

\textbf{2. Answers in your own words (réponses modèles)}
\begin{enumerate}
\item \eng{They had no place where they could learn digital skills: internet cafés were costly and computer training centres charged fees that their families could not pay.}
\item \eng{She reacted by being transparent: she published a detailed budget online and promised to give financial reports every month so that people could see how the money was spent.}
\item \eng{For her, the real success is not the money but the community that has formed around the project, because the hub belongs to everyone and will therefore survive.}
\end{enumerate}

\textbf{3. Equivalents}
\begin{enumerate}
\item \eng{pledged} (« a Nigerian company based in London pledged a larger amount »)
\item \eng{the money raised} (« the real success is not the money raised »)
\item \eng{enrolled} (« Forty young people enrolled in the first free course »)
\end{enumerate}

\textbf{4. Summary (environ 50 mots) — proposition :}

\eng{Adaeze launched an online crowdfunding campaign to open a tech hub for young graduates in Lagos. More than four hundred donors contributed, and a company doubled the donations. To gain people's trust, she published her budget and monthly reports. Six months later, the hub opened with twelve computers and forty students enrolled in a free course.} (50 mots)

\textbf{Points de vigilance :} citer le texte entre guillemets pour les justifications ; ne pas recopier de longues phrases pour les questions 2 ; respecter la limite de mots au résumé ; accorder les temps (prétérit pour le récit).
\end{corrige}

\begin{corrige}[Exercice 10 — Type Bac : Langue (grammaire et traduction)]
\textbf{Barème indicatif :} partie A : 5 pts (1 pt par transformation) ; partie B : 6 pts (2 pts par phrase).

\textbf{A. Grammar}
\begin{enumerate}
\item \eng{A fund-raising gala was organised by the youth club last month.} — Passif au prétérit : \emph{was + organised}.
\item \eng{The director said (that) they would build a new library the following year.} — \emph{shall/will} $\rightarrow$ \emph{would} ; \emph{next year} $\rightarrow$ \emph{the following year}.
\item \eng{The journalist asked the president if/whether he (or she) had received the grant.} — Question rapportée en yes/no : \eng{if/whether} + ordre sujet-verbe ; present perfect $\rightarrow$ past perfect.
\item \eng{Two boreholes have already been drilled by the community.} — Present perfect passif : \emph{have been + drilled}.
\item \eng{Unless the company sponsors us, we will not buy new computers.} — \eng{Unless} = \emph{if ... not} : « À moins que l'entreprise ne nous parraine, nous n'achèterons pas de nouveaux ordinateurs. » Après \eng{unless}, le verbe est à la forme affirmative (le sens négatif est contenu dans \eng{unless}) ; la négation porte sur la proposition principale (\emph{will not buy}).
\end{enumerate}

\textbf{B. Translation}
\begin{enumerate}
\item \eng{The villagers announced (that) they would build a health centre with the money (that had been) raised.} — Discours rapporté : \emph{would build} ; « l'argent collecté » = \eng{the money raised/collected}.
\item \eng{Every donation, even a small one, is important for our project.} — Ou : \eng{Every gift, however small, matters to our project.}
\item \eng{The association's accounts were published to guarantee transparency.} — Passif au prétérit ; « garantir la transparence » = \eng{to guarantee/ensure transparency}.
\end{enumerate}

\textbf{Points de vigilance :} au discours rapporté, toujours faire reculer les temps et adapter les pronoms et les expressions de temps ; au passif, vérifier l'accord de l'auxiliaire \emph{be} avec le nouveau sujet ; ne jamais traduire « collecter de l'argent » par \eng{win money} (faux ami avec « gagner »).
\end{corrige}

\begin{corrige}[Exercice 11 — Type Bac : Expression écrite (appeal letter)]
\textbf{Barème indicatif (sur 20) :} mise en page de la lettre formelle (adresses, date, salutation, formule de politesse) : 4 pts ; respect des quatre contenus exigés : 8 pts (2 pts chacun) ; correction grammaticale et vocabulaire : 5 pts ; cohérence et respect du nombre de mots (200--250) : 3 pts.

\textbf{Proposition de rédaction modèle :}

\begin{textebox}[Model appeal letter]
Mambanda Village Development Association\\
P.O. Box 125, Kumba\\
South West Region\\[0.3em]
15th March 2025\\[0.3em]
The General Manager\\
Riverside Mining Company\\
Kumba\\[0.3em]
Dear Sir or Madam,\\[0.3em]
\textbf{Request for sponsorship: community library project}\\[0.3em]
I am writing on behalf of the Mambanda Village Development Association, of which I am the secretary. Our association has decided to build a community library in the village, and we are looking for sponsors to help us complete this project.

At present, our young people have no quiet place to study and no access to books other than their school textbooks. A library would allow pupils to prepare for their examinations, help adults improve their literacy and give the whole community access to information and digital resources.

We would be grateful if your company could support us with a financial contribution of any amount, or with donations of books, shelves, tables and chairs.

We wish to assure you that all funds received will be managed with complete transparency: our accounts are audited and published every year. In return for your generosity, we would be pleased to display your company's name on the building and to mention your support at the opening ceremony and in the local press.

We hope you will consider our request favourably and remain at your disposal for any further information.\\[0.3em]
Yours faithfully,\\[0.3em]
Ngong Clarisse Ebai\\
Secretary, Mambanda Village Development Association
\end{textebox}

(Environ 220 mots pour le corps de la lettre.)

\textbf{Points de vigilance :}
\begin{itemize}
\item Mise en page : adresse de l'expéditeur en haut à gauche, date, adresse du destinataire, \eng{Dear Sir or Madam,} si le nom est inconnu, et \eng{Yours faithfully,} en conclusion (jamais \eng{Yours sincerely} avec \eng{Dear Sir}).
\item Ne pas traduire « demander » par \eng{demand} (trop agressif) : utiliser \eng{request}, \eng{ask for}, \eng{we would be grateful if...}.
\item « Parrainage » = \eng{sponsorship} ; « subvention » = \eng{grant} ; « rendre compte » = \eng{accountability}.
\item Éviter les phrases trop longues ; un paragraphe par idée (projet, utilité, demande, transparence).
\item Vérifier la longueur : compter les mots du corps de la lettre uniquement.
\end{itemize}
\end{corrige}

\newpage

\coursec{Fiche bilan}

\begin{popfigure}
\begin{tikzpicture}[mindmap, grow cyclic, every node/.style={concept, font=\small},
root concept/.append style={concept color=popblue, font=\small\bfseries},
level 1/.append style={level distance=3.4cm, sibling angle=51.4},
level 2/.append style={level distance=2.1cm, sibling angle=40, font=\footnotesize}]
\node[root concept]{Fund-raising \& Community Resources}
  child[concept color=poporange]{ node {Reading strategies}
    child { node {skimming} }
    child { node {scanning} }
    child { node {guessing words} }
  }
  child[concept color=popgreen]{ node {Key vocabulary}
    child { node {to raise funds} }
    child { node {donor, charity} }
    child { node {borehole, hall} }
  }
  child[concept color=poppurple]{ node {Reported speech}
    child { node {tense backshift} }
    child { node {say / tell / ask} }
  }
  child[concept color=poppink]{ node {Passive voice}
    child { node {be + past participle} }
    child { node {by + agent} }
  }
  child[concept color=popgold]{ node {Appeal letter}
    child { node {layout + tone} }
    child { node {persuasion} }
  }
  child[concept color=popblue!60]{ node {Speaking}
    child { node {meeting role-play} }
    child { node {giving opinions} }
  }
  child[concept color=popgreen!60]{ node {Comprehension}
    child { node {answer + justify} }
    child { node {summarise} }
  };
\end{tikzpicture}
\legende{Carte mentale de la leçon : lire, comprendre et réutiliser un texte sur la collecte de fonds communautaire.}
\end{popfigure}

\begin{bilan}
\textbf{1. Lexique clé du thème (à connaître par cœur)}

\begin{center}
\begin{tabularx}{\linewidth}{|X|X|}\hline
\rowcolor{popblueL} \textbf{English} & \textbf{Français} \\ \hline
to raise funds / a fund-raising event & collecter des fonds / une collecte de fonds \\ \hline
a donor / to donate / a donation & un donateur / faire un don / un don \\ \hline
a charity / a support group & une œuvre caritative / un groupe de soutien \\ \hline
a borehole / a community hall & un forage / une salle communautaire \\ \hline
to launch an appeal & lancer un appel (aux dons) \\ \hline
a contribution / to contribute & une contribution / contribuer \\ \hline
volunteers / to volunteer & des bénévoles / se porter volontaire \\ \hline
a target / to reach a target & un objectif / atteindre un objectif \\ \hline
\end{tabularx}
\end{center}

\textbf{2. Stratégies de lecture}
\begin{itemize}
\item \cle{Skimming} : lecture rapide pour l'idée générale (titre, premier et dernier paragraphes, premières phrases).
\item \cle{Scanning} : repérage d'une information précise (chiffre, nom, date, lieu).
\item \cle{Guessing} : déduire le sens d'un mot inconnu grâce au contexte, à la famille de mots et aux faux amis écartés.
\end{itemize}

\textbf{3. Grammaire à connaître par cœur}

\begin{center}
\begin{tabularx}{\linewidth}{|X|X|}\hline
\rowcolor{popgreenL} \textbf{Règle} & \textbf{Exemple} \\ \hline
Discours rapporté : recul des temps (present $\rightarrow$ past, will $\rightarrow$ would, past $\rightarrow$ past perfect) & “We \emph{will build} a hall,” he said. $\rightarrow$ He said they \emph{would build} a hall. \\ \hline
Verbes introducteurs : say (to sb), tell sb, ask sb, advise sb to + V & The chairman \emph{told} the villagers \emph{to contribute}. \\ \hline
Passif : be (au temps voulu) + participe passé (+ by + agent) & A borehole \emph{was drilled} by the support group. \\ \hline
Passif sans agent quand l'agent est évident ou inconnu & The money \emph{has been collected}. \\ \hline
\end{tabularx}
\end{center}

\textbf{4. Méthodes express}
\begin{itemize}
\item \cle{Répondre à une question} : réponse complète en anglais + citation courte du texte comme justification.
\item \cle{Résumer} : 3 à 5 phrases, ses propres mots, aucune phrase recopiée.
\item \cle{Lettre d'appel aux dons} : adresse et date, salutation, introduction (qui ? pourquoi ?), corps (projet, montant visé, utilisation des fonds), appel clair, formule de politesse, signature.
\item \cle{Réunion orale} : saluer, donner son avis (\eng{In my opinion...}, \eng{I suggest that...}), réagir (\eng{I agree / I disagree because...}), conclure.
\end{itemize}
\end{bilan}

\begin{aretenir}[Checklist avant le Bac]
Coche chaque point que tu maîtrises réellement :
\begin{itemize}
\item[$\square$] Je sais faire un skimming en moins de deux minutes et dégager l'idée générale d'un texte.
\item[$\square$] Je sais repérer rapidement un chiffre, un nom ou une date (scanning).
\item[$\square$] Je sais déduire le sens d'un mot inconnu à partir du contexte.
\item[$\square$] Je connais au moins quinze mots du lexique de la collecte de fonds et des ressources communautaires.
\item[$\square$] Je réponds aux questions par une phrase complète et je justifie par le texte.
\item[$\square$] Je sais résumer un texte en trois à cinq phrases sans recopier.
\item[$\square$] Je maîtrise le recul des temps au discours rapporté (will $\rightarrow$ would, etc.).
\item[$\square$] Je distingue \eng{say}, \eng{tell} et \eng{ask} sans erreur de construction.
\item[$\square$] Je transforme une phrase active en passive à tous les temps courants.
\item[$\square$] Je sais rédiger une lettre d'appel aux dons avec la mise en page correcte.
\item[$\square$] Je sais donner mon opinion et réagir poliment lors d'une réunion en anglais.
\item[$\square$] Je relis ma production pour vérifier l'accord sujet-verbe et les temps.
\end{itemize}
\end{aretenir}

\findecours

\end{document}
